\documentclass[11pt]{amsart}
\usepackage[margin=1.15in]{geometry}
\usepackage{amsmath,amssymb,amsthm}
\usepackage{booktabs}
\usepackage[hidelinks]{hyperref}

% ---------------------------------------------------------------------
% Theorem environments and macros: same conventions as the companion
% manuscript monotonicity/monotone-coupling.tex (numbering by section).
% ---------------------------------------------------------------------
\newtheorem{theorem}{Theorem}[section]
\newtheorem{proposition}[theorem]{Proposition}
\newtheorem{lemma}[theorem]{Lemma}
\newtheorem{corollary}[theorem]{Corollary}
\theoremstyle{definition}
\newtheorem{remark}[theorem]{Remark}
\newtheorem{definition}[theorem]{Definition}
\numberwithin{equation}{section}
\raggedbottom
\emergencystretch=3em

\newcommand{\eps}{\varepsilon}
\newcommand{\PP}{\mathbb P}
\newcommand{\EE}{\mathbb E}
\newcommand{\cE}{\mathcal E}
\newcommand{\re}{\operatorname{Re}}
\newcommand{\im}{\operatorname{Im}}
% visible marker for bibliographic details that still have to be checked
\newcommand{\todocite}[1]{\textsf{[TODO-cite: #1]}}
% notation of this article
\newcommand{\cA}{\mathcal A}          % compact group generated by the ordinates
\newcommand{\cZ}{\mathcal Z}          % set of distinct positive ordinates
\newcommand{\cK}{\mathcal K}          % certified ordinates outside the head
\newcommand{\cM}{\mathcal M}          % Cesaro mean of the race
\newcommand{\cN}{\mathcal N}          % normalised Cesaro mean
\newcommand{\cL}{\mathcal L}          % explicit lower-bound function
\newcommand{\hA}{\lambda_{\cA}}       % Haar measure on the group
\newcommand{\vth}{\vartheta_{o}}      % theta function over primes not dividing q
\newcommand{\dlow}{\underline{d}}
\newcommand{\dup}{\overline{d}}
\newcommand{\Ein}{E^{\theta}}
% status tags for numerical constants (Section 2.5)
\newcommand{\certA}{\textsf{[cert-A]}}
\newcommand{\certB}{\textsf{[cert-B]}}
\newcommand{\numtag}{\textsf{[num]}}

\begin{document}

\title[Weighted prime races under GRH alone]{Weighted prime races
under GRH alone: lead inheritance, convex order and explicit bias
bounds}

\author{Raghav Sharma}
\email{raghavsharma1729@gmail.com}
\thanks{Draft of November 2026, not yet submitted. This work was
carried out with substantial AI assistance (Anthropic's Claude): the
derivations, the proofs, the computations and the literature review
were AI-produced; what has and what has not been independently
verified is recorded in Sections~\ref{sec:comp} and~\ref{sec:gaps}
and in Appendix~\ref{sec:review}; the author takes responsibility for
the content. Comments are very welcome. Contact:
raghavsharma1729@gmail.com. Code and data:
\url{https://github.com/raghavs1729/weighted-prime-races}.}
\date{November 2026}
\subjclass[2020]{Primary 11N13; Secondary 11M26, 11Y35, 60E15}
\keywords{Chebyshev's bias, prime number races, weighted prime sums,
logarithmic density, convex order, explicit formula, Ces\`aro mean}

\begin{abstract}
For a real primitive character $\chi$ and a real exponent $m$ let
$D_m(x)=-\sum_{p\le x}\chi(p)p^m$ (for $\chi=\chi_4$: primes
$\equiv3$ against primes $\equiv1$ mod $4$, each weighted by $p^m$).
We study these races assuming GRH$(\chi)$ and, where stated,
$L(\tfrac12,\chi)\ne0$ (unconditional for $q=3,4$), but \emph{no}
linear independence hypothesis on the zeros.
(1) For $q=4$, under GRH$(\chi_4)$, the Ces\`aro mean
$x^{-1}\int_0^x(\pi(t;4,3)-\pi(t;4,1))\,dt$ is positive for
\emph{every} $x>3$, with an explicit lower bound, and for all
$m\le0$, $m<m'\le1$ and every $x>3$, $D_{m'}(x)\ge0$ implies
$D_m(x)>0$; the same holds mod~$3$ for every $x>2$. Eventual
nonnegativity of the Ces\`aro mean is equivalent to GRH$(\chi_q)$.
(2) The normalised race at a lighter weight is exactly the average of
the race at a heavier weight against an explicit probability kernel
supported on the past. Consequently the limiting distributions
increase strictly in the convex order as the weight increases.
(3) For every $m>-\tfrac12$ both lead sets have positive lower
logarithmic density; the variance of the limiting distribution
involves the squared multiplicities of the zeros, and we bound the
excess over the closed form explicitly. A one-sided Chebyshev
(Cantelli) inequality and a new bound based on the \emph{finite}
absence of low-order integer relations among certified zeros give:
under GRH$(\chi_4)$ alone, $\{x:\pi(x;4,3)>\pi(x;4,1)\}$ has lower
logarithmic density $>0.8842$, and under GRH$(\chi_3)$ alone the
corresponding set mod~$3$ has lower logarithmic density $>0.9112$
(the values under linear independence are $0.9959$ and $0.9991$).
(4) On the range where leads are inherited, upper and lower densities
are monotone in the weight and ties matter for at most countably many
weights. A final section states the analogue for
$\sum_{p\le x}p^m$ against $\mathrm{li}(x^{m+1})$. This work was
carried out with substantial AI assistance (see the first footnote
and Section~\ref{sec:gaps}).
\end{abstract}

\maketitle

\section{Introduction}\label{sec:intro}

\noindent\textbf{Disclosure.} This article was produced with
substantial AI assistance (Anthropic's Claude): the derivations,
proofs, computations and literature review are AI-produced. The two
working notes behind Sections~\ref{sec:resolvent}--\ref{sec:ties}
were each read by an adversarial AI referee and repaired;
Appendix~\ref{sec:review} summarises the referees' findings and their
resolution. The note behind Section~\ref{sec:zeta} was not refereed.
Section~\ref{sec:comp} lists every computation together with its
rigor class, Section~\ref{sec:formal} says what has and what has not
been formalized, and Section~\ref{sec:gaps} lists every place where a
statement is weaker than one might hope, depends on a computation, or
rests on a quoted result that was not re-derived here.

\subsection{Setting}\label{subsec:setting}

Throughout, $\chi$ is a real primitive Dirichlet character modulo
$q\ge3$ and $\chi^*:=-\chi$, so that $\chi^*(p)=+1$ on quadratic
nonresidues and $-1$ on residues. For real $m$ put
\[
D_m(x):=\sum_{p\le x}\chi^*(p)\,p^m,\qquad
R_m(x):=\sum_{p\le x}\chi^*(p)\,p^m\log p .
\]
For $\chi=\chi_4$, $D_0(x)=\pi(x;4,3)-\pi(x;4,1)$ is Chebyshev's race;
for $\chi=\chi_3$, $D_0(x)=\pi(x;3,2)-\pi(x;3,1)$. We write
\[
M=m+\tfrac12,\qquad M'=m'+\tfrac12,\qquad\kappa=m'-m=M'-M,
\]
and, except in Section~\ref{sec:explicit}, restrict to $m>-\tfrac12$,
i.e.\ $M>0$. The normalised races are
\begin{equation}\label{eq:normalised}
E_M(x):=2M\,x^{-M}(\log x)\,D_m(x),\qquad
\Ein_M(x):=2M\,x^{-M}\,R_m(x)\qquad(x>1);
\end{equation}
these are the functions $\cE_{m,0}$, $\cE_{m,1}$ of the companion
manuscript~\cite{MC} and $E_m'$, $E_m$ of~\cite{WR}.

The nontrivial zeros of $L(s,\chi)$ are written $\rho=\beta+i\gamma$
and are listed with multiplicity. \emph{GRH$(\chi)$} is the statement
that $\beta=\tfrac12$ for all of them; no hypothesis on $\zeta$ is
used anywhere in this article except in
Remark~\ref{rem:oneclass} and in Section~\ref{sec:zeta}, where it is
stated. We put
$n_0:=\operatorname{ord}_{s=1/2}L(s,\chi)\ge0$, and
\emph{NCZ$(\chi)$} is the statement $n_0=0$, i.e.\
$L(\tfrac12,\chi)\ne0$. NCZ is unconditional for $q\in\{3,4\}$ (the
series \eqref{eq:Lpositive} below are positive for every $\sigma>0$),
and for every real primitive $\chi$ of conductor $\le10^9$, resp.\
$\le10^{10}$ if $\chi$ is odd, by Platt--Trudgian
\cite[Theorem~1]{PT26}. \emph{LI$(\chi)$}, the linear independence
over $\mathbb Q$ of the positive ordinates, is \textbf{not} assumed
anywhere in this article; the few places where a statement under LI
is quoted for comparison say so explicitly.

Under GRH$(\chi)$, $\cZ$ denotes the set of \emph{distinct} positive
ordinates and $m_\gamma\ge1$ the multiplicity of $\gamma\in\cZ$. Sums
written $\sum_{\gamma>0}$ run over the positive ordinates \emph{with
multiplicity}; sums written $\sum_{\gamma\in\cZ}$ run over distinct
ordinates. For $M,M'>0$ put
\begin{gather}
r_\gamma(M):=\frac{4M}{M+i\gamma},\qquad
c_\gamma(M):=|r_\gamma(M)|=\frac{4M}{\sqrt{M^2+\gamma^2}},\qquad
w_M(\gamma):=\frac{8M^2}{M^2+\gamma^2}=\tfrac12c_\gamma(M)^2,
\notag\\
\sigma_M^2:=\sum_{\gamma>0}w_M(\gamma)=\sum_{\gamma\in\cZ}m_\gamma\,
w_M(\gamma),\qquad
V_M:=\sum_{\gamma\in\cZ}m_\gamma^2\,w_M(\gamma),
\label{eq:defV}\\
\lambda_\chi(M,M'):=\frac{1+2n_0}{MM'}
-\sum_{\gamma>0}\frac{4}{|M+i\gamma|\,|M'+i\gamma|},
\label{eq:deflambda}\\
C(\chi):=\sum_{\gamma>0}\frac{2}{\gamma^2},\qquad
C_2(\chi):=\sum_{\gamma\in\cZ}\frac{2m_\gamma^2}{\gamma^2}.\notag
\end{gather}
Thus $V_M\ge\sigma_M^2$, with equality if and only if all zeros are
simple. For a measurable $S\subset[1,\infty)$, $\dlow(S)$ and
$\dup(S)$ denote the lower and upper logarithmic densities, the
$\liminf$ and $\limsup$ as $X\to\infty$ of
$(\log X)^{-1}\int_{S\cap[1,X]}dt/t$.

\subsection{Results}\label{subsec:results}

The asymptotic backbone is Theorem~A of the companion
manuscript~\cite{MC}, restated as Theorem~\ref{thm:MCA} below: under
GRH$(\chi)$ alone, if $\lambda_\chi(M,M')>0$ then
$D_{m'}(x)\ge0\Rightarrow D_m(x)>0$ for all $x\ge x_0(\chi,m,m')$,
with $x_0$ not explicit. We do not re-prove it. The results of the
present article are of three kinds: statements for \emph{every} $x$
for the two smallest moduli, structural statements about the family
of limiting distributions, and explicit lower bounds for the bias.
The headline statements are collected in the five theorems below;
each is a combination of results stated and proved in later sections,
to which the proofs refer.

For $q\in\{4,3\}$ let $\chi_q$ be the nonprincipal character mod $q$,
$p_0$ the least prime not dividing $q$ ($p_0=3$ for $q=4$, $p_0=2$
for $q=3$), and
\begin{equation}\label{eq:defTheta1}
\begin{gathered}
\Theta_1(t):=\sum_{p\le t}\chi_q^*(p)\,(\log p)\,(t-p)=\int_0^tR_0(u)\,du,\\
\cM(x):=\sum_{p\le x}\chi_q^*(p)\Bigl(1-\frac px\Bigr)
=D_0(x)-\frac{D_1(x)}{x}=\frac1x\int_0^xD_0(t)\,dt .
\end{gathered}
\end{equation}
So $\cM$ is the Ces\`aro mean of the race; for $q=4$,
$\cM(x)=x^{-1}\int_0^x(\pi(t;4,3)-\pi(t;4,1))\,dt$.

\begin{theorem}[every $x$; $q=4$ and $q=3$]\label{thm:allx}
Let $q\in\{4,3\}$ and assume GRH$(\chi_q)$. Let $\cL_{\chi_q}$ be
the explicit function \eqref{eq:defL}.
\begin{enumerate}
\item[(i)] $\Theta_1(t)\ge\cL_{\chi_q}(t)\,t^{3/2}$ for every
$t\ge p_0$, and $\Theta_1(t)>0$ for every $t>p_0$. Here
$\cL_{\chi_q}(t)>0$ for $t\ge10^4$ and
$\cL_{\chi_q}(t)\to0.51104$ $(q=4)$, $0.55338$ $(q=3)$ as
$t\to\infty$.
\item[(ii)] For every $x>p_0$,
\[
\cM(x)\;\ge\;\frac{\Theta_1(x)}{x\log x}\;>\;0 .
\]
That is, $\int_0^x(\pi(t;4,3)-\pi(t;4,1))\,dt>0$ for every $x>3$
under GRH$(\chi_4)$, and
$\int_0^x(\pi(t;3,2)-\pi(t;3,1))\,dt>0$ for every $x>2$ under
GRH$(\chi_3)$.
\item[(iii)] $\cM(x)>\sqrt x/(2\log x)$ for every $x\ge23$ if $q=4$,
and for every $x\ge6$ if $q=3$.
\item[(iv)] For all real $m\le0$ and $m<m'\le1$, with $\kappa=m'-m$,
and every $x>p_0$,
\[
D_m(x)-x^{-\kappa}D_{m'}(x)\;\ge\;
\frac{\kappa\,x^{m-1}\,\Theta_1(x)}{\log x}\;>\;0 .
\]
In particular, at every $x>p_0$: $D_{m'}(x)\ge0$ implies $D_m(x)>0$,
and $D_m(x)\le0$ implies $D_{m'}(x)<0$.
\item[(v)] Conversely, and with no hypothesis: if
$\int_0^xD_0(t)\,dt\ge0$ for all sufficiently large $x$, then
GRH$(\chi_q)$ holds. Hence GRH$(\chi_q)$, the positivity of
$\int_0^xD_0(t)\,dt$ for every $x>p_0$, and its nonnegativity for
all sufficiently large $x$ are equivalent.
\end{enumerate}
\end{theorem}

Parts (i)--(iv) are Theorems~\ref{thm:Theta1}, \ref{thm:cesaro},
\ref{thm:cone} and Corollaries~\ref{cor:half}, \ref{cor:inherit};
part (v) is Theorem~\ref{thm:equiv}. The proofs combine an analytic
argument for $t\ge10^4$ with a finite verification for $t\le10^4$
done in exact integer arithmetic; (iii) uses, in addition, a
ball-arithmetic computation up to $1.6\cdot10^8$ $(q=4)$, resp.\
$2.6\cdot10^7$ $(q=3)$, a double-precision evaluation of an explicit
function on a grid, and a cited explicit bound for $\theta(x)$
(Sections~\ref{sec:explicit} and~\ref{sec:comp} say exactly what is
certified and how). In words, (iv) with $(m,m')=(0,1)$ says: under
GRH$(\chi_4)$, at \emph{every} $x>3$, a lead or tie of the class
$3$ mod $4$ in the race weighted by $p$ is a strict lead in
Chebyshev's race, and a lead or tie of the class $1$ mod $4$ in
Chebyshev's race is a strict lead of that class in the race weighted
by $p$.

\begin{theorem}[averaging identity and convex order]\label{thm:order}
Let $\chi$ be a real primitive character and $0<M<M'$. Let
$\mu_{M,M'}$ be the probability measure
$\frac{M}{M'}\delta_0+(1-\frac{M}{M'})\,Me^{-Mu}\,du$ on $[0,\infty)$.
\begin{enumerate}
\item[(a)] Unconditionally, for every real $y$,
$\Ein_M(e^y)=\int_{[0,\infty)}\Ein_{M'}(e^{y-u})\,\mu_{M,M'}(du)$.
\item[(b)] Assume GRH$(\chi)$ and NCZ$(\chi)$. Then $E_M(e^y)$ and
$\Ein_M(e^y)$ have a common limiting distribution $\nu_M$, of mean
$1$, and for every convex $\varphi:\mathbb R\to\mathbb R$
\[
\int\varphi\,d\nu_M\;\le\;\int\varphi\,d\nu_{M'},
\]
with strict inequality if $\varphi$ is not affine and the right-hand
side is finite. In particular the variance $V_M$, every absolute
central moment, and the mean deficit $\int x^-\,d\nu_M$ are strictly
increasing functions of $M$.
\end{enumerate}
\end{theorem}

This is Theorem~\ref{thm:resolvent}(a), Theorem~\ref{thm:limit},
Theorems~\ref{thm:cx}, \ref{thm:strict} and Corollary~\ref{cor:mono}.

\begin{theorem}[both lead sets; variance; one-sided bias bound]
\label{thm:bias}
Let $\chi$ be a real primitive character, assume GRH$(\chi)$ and
NCZ$(\chi)$, and let $m>-\tfrac12$. Then
\[
\dlow\{x:D_m(x)<0\}>0,\qquad
\dlow\{x:D_m(x)>0\}\;\ge\;\frac1{1+V_M}>0,
\]
where $V_M=\sum_{\gamma\in\cZ}m_\gamma^2\,8M^2/(M^2+\gamma^2)$ is the
variance of $\nu_M$; moreover
$V_M\ge\sigma_M^2=4M\frac{\Lambda'}{\Lambda}(M+\tfrac12,\chi)$ with
equality if and only if all zeros of $L(s,\chi)$ are simple, and
$\dlow\{D_m>0\}\to1$ as $m\to-\tfrac12^+$.
\end{theorem}

This is Theorems~\ref{thm:leadsets}, \ref{thm:var} and
\ref{thm:cantelli}.

\begin{theorem}[Chebyshev's race under GRH alone]\label{thm:cheb}
\begin{enumerate}
\item[(a)] Assume GRH$(\chi_4)$. Then
$\dlow\{x:\pi(x;4,3)>\pi(x;4,1)\}>0.884276434$.
\item[(b)] Assume GRH$(\chi_3)$. Then
$\dlow\{x:\pi(x;3,2)>\pi(x;3,1)\}>0.911240824$.
\end{enumerate}
Both constants are certified in ball arithmetic from the zero lists
of the repository (status \certA{} of Section~\ref{subsec:tags}). If
in addition all zeros of $L(s,\chi_q)$ up to height $10^8/q$ are
simple, the bounds improve to $0.922718033$ and $0.944555790$
(status \certB). The one-sided Chebyshev inequality of
Theorem~\ref{thm:bias} alone gives $0.8549$ and $0.8885$ \certA.
\end{theorem}

This is Theorems~\ref{thm:cantelli-q} and \ref{thm:NRq}; the bound
behind (a), (b) is Theorem~\ref{thm:NR}, which converts the
\emph{finite} statement ``the $N$ lowest ordinates satisfy no integer
relation of order $\le4$, and resonate with no other certified
ordinate at the orders that occur'' into a density bound. For
comparison, under GRH and LI the densities exist and equal
$0.9959\ldots$ and $0.9991\ldots$ \cite{RS94}.

\begin{theorem}[monotonicity and ties]\label{thm:tiesintro}
Assume GRH$(\chi)$ and let $I\subset(-\tfrac12,\infty)$ be an
interval of weights such that $\lambda_\chi(M,M')>0$ for all $m<m'$
in $I$ (for instance $I=(-\tfrac12,M_\sigma(\chi)-\tfrac12]$, where
$M_\sigma(\chi)$ is the root of $\sigma_M^2=2$). Then the four
functions $m\mapsto\dlow\{D_m>0\}$, $\dup\{D_m>0\}$,
$\dlow\{D_m\ge0\}$, $\dup\{D_m\ge0\}$ are nonincreasing on $I$, and
for all but countably many $m\in I$ counting ties for or against the
leader changes neither the upper nor the lower logarithmic density.
\end{theorem}

This is Theorem~\ref{thm:ties}; the monotonicity is
\cite[the corollary on densities without LI]{MC}, the statement about ties is new here.

\subsection{What is new and what is not}\label{subsec:new}

\emph{Theorem~\ref{thm:allx}.} Nothing in the method is new. The
mechanism --- a positive bias from prime squares beating an
absolutely convergent sum over zeros after one integration, and
Landau's theorem for the converse --- is that of Hardy--Littlewood
and Landau \cite{HL18,Lan18} for the Abel means
$\sum_p\chi_4(p)e^{-p/x}$, of Pintz \cite{Pin91} and Johnston
\cite{Joh23} for $\pi-\mathrm{li}$, and of Zhao \cite{Zhao25} for the
error terms in Mertens' theorems (including quadratic characters, for
all sufficiently large~$X$); see also Fujii \cite{Fuj88} and
Platt--Trudgian \cite{PT19}. Johnston's theorem \cite[Theorem~1.1]{Joh23}
states that RH is equivalent to $\int_2^x(\pi(t)-\mathrm{li}(t))\,dt<0$
for all $x>2$. \emph{His proof is the template for
Theorem~\ref{thm:allx}}: an explicit formula for $\psi_1$, the bias
from $\psi-\theta$, a second-order partial summation (his Lemma~2.7,
after Ingham \cite[pp.~103--104]{Ing32}) to pass from $\psi$ to $\pi$,
a finite verification, and an $\Omega$-result for the converse; and
\emph{his Section~6 lists the Chebyshev bias for primes in arithmetic
progressions as further work}. Theorem~\ref{thm:allx} is the direct
character analogue of his theorem and should be read as such. The
summation identity (Lemma~\ref{lem:ingham}) is Ingham's; the Landau
converse (v) is the Landau mechanism for one more kernel and is the
case $(m,m')=(0,1)$ of \cite[the corollary ``equivalence with GRH for $q\in\{3,4\}$'', part~(a)]{MC}. What we did not
find in what we read --- and our reading is limited: the abstract
only of Liu--Wang--Shi \cite{LWS26}, who announce a classification of
the residue-class pairs for which an integrated prime counting error
is eventually negative; not Suzuki \cite{Suz25}, Fujii, or the
originals of Hardy--Littlewood, Landau, Pintz and Ingham --- is:
\begin{enumerate}
\item[(a)] the statement itself for the Ces\`aro mean of Chebyshev's
race, for \emph{every} $x>3$ (resp.\ $x>2$), with explicit constants;
\item[(b)] its reading as inheritance of leads across weights at
every $x$ (Theorem~\ref{thm:allx}(iv));
\item[(c)] the observation that the integrated $\theta$-race
$\Theta_1$ is one-signed, so that the positive kernel of Ingham's
partial summation can be dropped with no loss --- in particular
without the loss of a factor $1/\log x$ --- and its two uses, the
cone $m\le0$, $m<m'\le1$ and the reduction for an arbitrary pair of
weights (Lemma~\ref{lem:pairs}).
\end{enumerate}
In Johnston's setting the integrated error is not one-signed and the
kernel term has to be estimated; this is why the analytic part of
his argument starts at $10^{19}$ and ours at $10^4$.

\emph{Theorem~\ref{thm:order}.} The identity (a) is one line of
partial summation and underlies \cite[Theorem~A]{MC}. Its reading as
an average against a probability kernel, the semigroup law
(Theorem~\ref{thm:resolvent}(d)) and the convex order with its
strictness were not found in the literature searched (searches for
``convex order'', ``stochastic order'', ``majorization'' together
with prime races or Rubinstein--Sarnak). The representation of the
limiting distribution as the push-forward of Haar measure on the
closure of a one-parameter subgroup, without LI, is known
(Rubinstein--Sarnak \cite[\S2]{RS94}, Devin \cite[\S4]{Dev20},
Bailleul \cite{Bai21}); we claim no novelty for
Theorem~\ref{thm:limit}.

\emph{Theorems~\ref{thm:bias} and \ref{thm:cheb}.} \emph{The general
Chebyshev-type inequality is in Devin and Fiorilli.} Fiorilli
\cite[Lemma~2.7]{Fio14} proves, under GRH and without LI,
$\dlow\ge1-2\operatorname{Var}/\mathrm{mean}^2$ when the mean is large
compared with the standard deviation, and Devin
\cite[Corollary~5.8]{Dev20} proves $\dlow\ge1-\operatorname{Var}/
\mathrm{mean}^2$ whenever the mean is positive; the variance with
\emph{squared} multiplicities is \cite[Lemma~2.5]{Fio14} and
\cite[Theorem~2.1]{Dev20}. For $\chi_4$ Devin's corollary reads
$\dlow\ge1-V_{1/2}$, which would be $0.844$ if the zeros are simple.
We found no paper that evaluates such a bound for $q=4$ or $q=3$, and
none that controls the multiplicities so as to make the number
unconditional under GRH. What is added here is: the one-sided
(Cantelli) form, which improves $1-V$ to $1/(1+V)$; the explicit
bound for the multiplicity excess $V_M-\sigma_M^2$
(Lemmas~\ref{lem:mult}, \ref{lem:JA},
Propositions~\ref{prop:Vplus}, \ref{prop:tierB}); the certified
evaluation; the family of weights; and the two improvements over
Cantelli (Theorems~\ref{thm:peel} and~\ref{thm:NR}). Positivity of
both densities under GRH alone is known at $m=0$
(Rubinstein--Sarnak \cite[Theorem~1.2]{RS94}, Devin
\cite[Theorem~2.5(1)]{Dev20}, and Kaczorowski \cite{Kac93} for
natural lower densities); for $m\ne0$ we did not find it stated.
Numerical verifications of the absence of low-order relations among
zeros of $\zeta$ are due to Best--Trudgian \cite{BT15} (not read by
us); we did not find such checks used to \emph{prove} a density
bound.

\emph{The endpoint $m=-\tfrac12$.} Hayani \cite{Hay25} and,
independently, Sheth \cite{Sheth24} prove density-one statements at
the single critical weight $p^{-1/2}$ under GRH alone.
Proposition~\ref{prop:endpoint} is the quantitative approach to that
endpoint from inside the range $m>-\tfrac12$.

\emph{Sums of powers of primes.} Sums $\sum_{p\le x}p^k$ have been
studied in comparison with $\pi(x^{k+1})$ by Gerard--Washington
\cite{GW18}, Washington \cite{Was24} and, in arithmetic progressions,
by Boran, Byun, Li, Miller and Reyes \cite{BBLMR24}. Those works
compare one weighted sum with an unweighted count and do not race two
residue classes with a common weight. Section~\ref{sec:zeta} states
the analogue of the present results for $\sum_{p\le x}p^m$ against
$\mathrm{li}(x^{m+1})$.

\begin{remark}[on the status of GRH]\label{rem:GRHstatus}
A proof of GRH for Dirichlet $L$-functions was announced on
20~November 2026 \cite{Alp26}, accompanied by a Lean formalization
that \cite{Lean26} reports to have checked; no human-readable proof
has been refereed. We have not examined either, and do not use them.
\emph{Every result of this article remains stated conditionally on
GRH$(\chi)$.} If that proof is confirmed,
Theorems~\ref{thm:allx} and \ref{thm:tiesintro} and the \certA{}
statements of Theorem~\ref{thm:cheb} become unconditional, and
Theorems~\ref{thm:order}(b) and~\ref{thm:bias} become unconditional
for every real primitive $\chi$ for which NCZ$(\chi)$ is known.
Neither NCZ nor LI is addressed by that work, and the Lean
development described in Section~\ref{sec:formal} does not depend on
it.
\end{remark}

\subsection{Organisation}
Section~\ref{sec:prelim} collects notation, the quoted explicit
formulas and the status tags. Section~\ref{sec:resolvent} proves the
averaging identity. Section~\ref{sec:limit} identifies the limiting
distribution and proves the convex order;
Section~\ref{sec:leadsets} treats the two lead sets.
Section~\ref{sec:variance} computes the variance and bounds the
multiplicity excess; Section~\ref{sec:biasbounds} contains the
Cantelli and non-resonance bounds. Section~\ref{sec:explicit} proves
the statements for every $x$ and Section~\ref{sec:landau} the Landau
converse. Section~\ref{sec:ties} treats monotonicity and ties.
Section~\ref{sec:zeta} states the zeta analogue.
Section~\ref{sec:comp} lists the computations and their rigor class,
Section~\ref{sec:formal} says exactly what has and what has not been
formalized in Lean, Section~\ref{sec:gaps} lists the gaps and
caveats, and
Appendix~\ref{sec:review} the review history.

\section{Notation, quoted explicit formulas and conventions}
\label{sec:prelim}

Implied constants may depend on $\chi$ and on the fixed exponents
$m,m'$, never on $x$, $t$, $u$, $y$, $T$ or on the zero under
consideration. We write $y=\log x$ throughout,
\[
f_M(y):=E_M(e^y),\qquad g_M(y):=\Ein_M(e^y),
\]
and note that both vanish for $y<\log2$ (there are no primes below
$2$). Further,
\begin{gather*}
\psi(x,\chi):=\sum_{n\le x}\Lambda(n)\chi(n),\qquad
\psi_m(x):=\sum_{n\le x}\Lambda(n)\chi(n)n^m,\\
\psi_1(x,\chi):=\int_0^x\psi(t,\chi)\,dt
=\sum_{n\le x}\Lambda(n)\chi(n)(x-n).
\end{gather*}
The digamma function is written $\Gamma'/\Gamma$.

\subsection{Zeros}\label{subsec:zeros}
We use the following standard facts \cite[\S\S12, 16]{Dav00}. The
nontrivial zeros satisfy $0<\beta<1$, they are symmetric under
$\rho\mapsto\bar\rho$ because $\chi$ is real, $\sum_\rho|\rho|^{-2}
<\infty$, and, with $N(T)$ the number of zeros with $0<\gamma\le T$
(counted with multiplicity),
\begin{equation}\label{eq:zerocount}
N(t+1)-N(t)\le A_0\log(t+3)\quad(t\ge0),\qquad
N(T)\sim\frac{T}{2\pi}\log T ,
\end{equation}
for a constant $A_0=A_0(\chi)$. Consequently, under GRH$(\chi)$,
$\sum_{\gamma>0}\gamma^{-2}<\infty$, all the series in
\eqref{eq:defV}--\eqref{eq:deflambda} that are sums over zeros with
multiplicity converge, and $\cZ$ is infinite. (Finiteness of $V_M$
and $C_2(\chi)$, which involve $m_\gamma^2$, follows from
\eqref{eq:zerocount}, since $m_\gamma\le A_0\log(\gamma+3)$; an
explicit bound is Lemma~\ref{lem:mult}.)

For $q\in\{4,3\}$ the Dirichlet series converges for real $\sigma>0$
and, grouping consecutive terms,
\begin{equation}\label{eq:Lpositive}
\begin{aligned}
L(\sigma,\chi_4)&=\sum_{n\ge0}\bigl((4n+1)^{-\sigma}-(4n+3)^{-\sigma}\bigr)>0,\\
L(\sigma,\chi_3)&=\sum_{n\ge0}\bigl((3n+1)^{-\sigma}-(3n+2)^{-\sigma}\bigr)>0
\end{aligned}
\qquad(\sigma>0).
\end{equation}
In particular NCZ$(\chi_4)$ and NCZ$(\chi_3)$ hold unconditionally,
and neither $L$-function has a real zero in $(0,1)$.

\subsection{Quoted inputs}\label{subsec:quoted}
We quote four results of the first paper \cite{WR} of this series, in
the form in which they are stated or used in the companion
manuscript \cite{MC}. Theorem numbers of \cite{WR} refer to its
August 2026 version. The manuscript \cite{MC} was being revised while
this article was written, so its results other than Theorem~A are
cited by name rather than by number.

\begin{lemma}[{\cite[Lemma~1]{WR}}; truncated explicit formula, all
truncations]\label{lem:WRclassical}
Unconditionally, for all $x\ge2$ and $T\ge2$,
\[
\psi(x,\chi)=-\sum_{|\gamma|\le T}\frac{x^\rho}{\rho}
+O\Bigl(\frac{x\log^2(xT)}{T}+\log x\Bigr),
\]
with no inequality assumed between $T$ and $x$.
\end{lemma}

\begin{lemma}[{\cite[Lemmas~2 and~3]{WR}}; shifted formula and
prime-power separation]\label{lem:WRshift}
Assume GRH$(\chi)$ and NCZ$(\chi)$ and fix $m>-\tfrac12$. Uniformly in
$t\ge2$ and $T\ge2$,
\[
\psi_m(t)=-\sum_{|\gamma|\le T}\frac{t^{\rho+m}}{\rho+m}
+O\Bigl(\frac{t^{m+1}\log^2(tT)}{T}+(1+t^m)\log t\Bigr),
\]
and, as $t\to\infty$,
\[
R_m(t)=-\psi_m(t)+\frac{t^{M}}{2M}+o(t^{M}).
\]
\end{lemma}

\begin{lemma}[{\cite[proofs of Theorem~1(i) and of
Proposition~1]{WR}}; {\cite[Section ``Inputs quoted from the companion paper'']{MC}}]\label{lem:WRexp}
Assume GRH$(\chi)$ and NCZ$(\chi)$ and fix $M>0$. For
$h\in\{f_M,g_M\}$, all $u\ge2$ and all $T\ge2$,
\begin{equation}\label{eq:WRexp}
h(u)\;=\;1+\re\sum_{0<\gamma\le T}r_\gamma(M)\,e^{i\gamma u}
+\cE(u,T),\qquad
|\cE(u,T)|\ll\frac{u\,e^{u/2}\log^2(e^uT)}{T}+d(u),
\end{equation}
where $d$ is a fixed bounded function with $d(u)\ll1/u$.
\end{lemma}

(In \cite{MC}, where this lemma is quoted in the same form, the function $d$ is written out:
$d(u)=(1-b)/u+u^2e^{-\min(M,1/2)u}+ue^{-u/6}+ue^{-c\sqrt u}$ with
$b=1$ for $g_M$ and $b=0$ for $f_M$.) Subtracting the two instances
of \eqref{eq:WRexp} with the same $T$ and letting $T\to\infty$ gives,
under GRH$(\chi)$ and NCZ$(\chi)$,
\begin{equation}\label{eq:fminusg}
f_M(y)-g_M(y)=O(1/y)\qquad(y\ge2),
\end{equation}
pointwise; this is \cite[the theorem on logarithmic weights, part~(i)]{MC} for the two exponents of
the logarithm that occur here.

\begin{proposition}[{\cite[Proposition~2 and eq.~(8)]{WR}};
{\cite[the proposition ``variance identity'']{MC}}; Hadamard product and variance
identity]\label{prop:WRvar}
Assume GRH$(\chi)$ and NCZ$(\chi)$, let $\chi(-1)=(-1)^a$,
$\Lambda(s,\chi)=(q/\pi)^{(s+a)/2}\Gamma\bigl(\tfrac{s+a}2\bigr)
L(s,\chi)$ and $\Xi(z):=\Lambda(\tfrac12+z,\chi)$. Then
\[
\frac{\Xi(z)}{\Xi(0)}=\prod_{\gamma>0}\Bigl(1+\frac{z^2}{\gamma^2}\Bigr),
\]
the product converging absolutely and locally uniformly on
$\mathbb C$. Hence, for every complex $z$ with $\re z>0$,
\begin{equation}\label{eq:XiLogDer}
\frac{\Xi'}{\Xi}(z)=\sum_{\gamma>0}\frac{2z}{z^2+\gamma^2}
=\sum_{\gamma>0}\Bigl(\frac1{z-i\gamma}+\frac1{z+i\gamma}\Bigr),
\end{equation}
and in particular
\[
\sigma_M^2=4M\,\frac{\Lambda'}{\Lambda}\bigl(M+\tfrac12,\chi\bigr)
\quad(M>0),
\qquad
C(\chi)=(\log\Lambda)''\bigl(\tfrac12,\chi\bigr).
\]
\end{proposition}

(In \cite{WR} the product and the first identity are proved for
$\chi_4$ and it is observed, at eq.~(8) there, that the proof applies
verbatim to a general real primitive $\chi$. The passage from the
product to \eqref{eq:XiLogDer} at complex $z$ is logarithmic
differentiation of a locally uniformly convergent product of entire
functions, valid wherever $\Xi\ne0$, in particular for $\re z>0$.)

\subsection{The asymptotic backbone}\label{subsec:backbone}
The following theorem of the companion manuscript is used, not
re-proved.

\begin{theorem}[{\cite[Theorem~A]{MC}}; inheritance of leads, GRH
only]\label{thm:MCA}
Let $\chi$ be a real primitive character and assume GRH$(\chi)$
(neither LI$(\chi)$ nor NCZ$(\chi)$ is assumed). Let
$-\tfrac12<m<m'$.
\begin{enumerate}
\item[(i)] For all $x\ge2$,
\begin{align*}
D_m(x)-x^{-\kappa}D_{m'}(x)
&=\kappa\int_2^xt^{-\kappa-1}D_{m'}(t)\,dt\\
&=\frac{\kappa\,x^{M}}{\log x}\Bigl[\frac{1+2n_0}{2MM'}
+2\,\re\sum_{\gamma>0}\frac{x^{i\gamma}}{(M+i\gamma)(M'+i\gamma)}\Bigr]
+O\Bigl(\frac{x^{M}}{\log^2x}\Bigr),
\end{align*}
where the first equality is an exact identity and the series
converges absolutely.
\item[(ii)] If $\lambda_\chi(M,M')>0$, there is
$x_0=x_0(\chi,m,m')$ such that for all $x\ge x_0$
\[
D_m(x)-x^{-\kappa}D_{m'}(x)\;\ge\;
\frac{\kappa\,\lambda_\chi(M,M')}{4}\cdot\frac{x^{M}}{\log x}\;>\;0 .
\]
In particular, for every $x\ge x_0$: $D_{m'}(x)\ge0$ implies
$D_m(x)>0$, and $D_m(x)\le0$ implies $D_{m'}(x)<0$.
\item[(iii)] $\lambda_\chi(M,M')\ge
\bigl(2+4n_0-\sigma_M\sigma_{M'}\bigr)/(2MM')$. In particular
$\lambda_\chi(M,M')>0$ whenever $\sigma_M\sigma_{M'}<2+4n_0$, hence
whenever $\sigma_{M'}^2\le2$.
\end{enumerate}
\end{theorem}

Let $M_\sigma(\chi)$ be the root of $\sigma_M^2=2$. By (iii),
$\lambda_\chi(M,M')>0$ for all $0<M<M'\le M_\sigma(\chi)$
\cite[the remark on the range of validity of Theorem~A]{MC}. Numerically
$M_\sigma(\chi_4)=1.8239\ldots$ and $M_\sigma(\chi_3)=2.1277\ldots$
\numtag{} (closed form of Proposition~\ref{prop:WRvar}; the
certificate files of \cite{MC} enclose them in
$[1.82391477,1.82391478)$ and $[2.12775095,2.12775096)$). The
threshold $x_0$ of (ii) is not explicit; Section~\ref{sec:explicit}
shows that for $q\in\{4,3\}$ and the pairs with $m\le0$, $m<m'\le1$
the two implications of (ii) hold for every $x>p_0$.

\subsection{Densities and limiting distributions}
\label{subsec:densities}
A locally bounded measurable $\phi:[0,\infty)\to\mathbb R$ \emph{has
the limiting distribution} $\nu$ (a Borel probability measure on
$\mathbb R$) if the laws $\nu_Y$ of $\phi(y)$, $y$ uniform on $[0,Y]$,
converge weakly to $\nu$ as $Y\to\infty$. By the portmanteau theorem
\cite[Theorem~2.1]{Bil99}, then
$\limsup_Y\nu_Y(B)\le\nu(B)$ for closed $B$ and
$\liminf_Y\nu_Y(B)\ge\nu(B)$ for open $B$. For $S\subset[1,\infty)$
and $X=e^Y$, $(\log X)^{-1}\int_{S\cap[1,X]}dt/t
=Y^{-1}\,\mathrm{meas}\{y\in[0,Y]:e^y\in S\}$, so logarithmic
densities of sets $\{x:E_M(x)\in B\}$ are limits of $\nu_Y(B)$. For
$x>1$ the signs of $D_m(x)$ and $E_M(x)$ agree, and those of
$R_m(x)$ and $\Ein_M(x)$ agree.

\subsection{Status tags for numerical constants}\label{subsec:tags}
Numerical constants in Sections~\ref{sec:variance}
and~\ref{sec:biasbounds} carry one of three tags.
\begin{itemize}
\item \certA: enclosed in ball arithmetic (Arb) from the certified
zero lists of the repository,
\texttt{prototype/explore/chi4\_zeros.txt} ($511$ zeros, complete and
simple up to height $T_4=640.165707$) and
\texttt{chi3\_zeros.txt} ($537$ zeros, up to $T_3=700.024747$); each
ordinate is enclosed with radius $10^{-20}$, and completeness and
simplicity are certified by sign changes together with a winding
count (\texttt{certify\_arb.py --K 0}).
\item \certB: as \certA, but assuming in addition that all zeros of
$L(s,\chi_q)$ up to height $10^8/q$ are simple
(Remark~\ref{rem:tierB}).
\item \numtag: numerical estimate only.
\end{itemize}
The rigor classes of the computations of Section~\ref{sec:explicit}
(exact integers, ball arithmetic, \texttt{long double} with stated
margins, double precision) are stated where they are used and
collected in Section~\ref{sec:comp}.

\section{The averaging (resolvent) identity}\label{sec:resolvent}

For $0<M<M'$ let $\kappa=M'-M$ and
\begin{equation}\label{eq:defmu}
\mu=\mu_{M,M'}:=\frac{M}{M'}\,\delta_0
+\Bigl(1-\frac{M}{M'}\Bigr)\,Me^{-Mu}\,du\qquad\text{on }[0,\infty),
\end{equation}
a \emph{probability} measure, with Fourier transform
$\hat\mu(\xi):=\int e^{-i\xi u}\,\mu(du)$. Write
$K_Mh(y):=\int_0^\infty Me^{-Mu}h(y-u)\,du$ (the resolvent
$M(M+d/dy)^{-1}$).

\begin{theorem}\label{thm:resolvent}
Let $\chi$ be a real primitive character and $0<M<M'$. Parts (a) and
(d) are unconditional; parts (b) and (c) assume GRH$(\chi)$ and
NCZ$(\chi)$.
\begin{enumerate}
\item[(a)] \emph{Exact identity for the $\theta$-form.} For every
real $y$,
\begin{align*}
\Ein_M(e^y)&=\int_{[0,\infty)}\Ein_{M'}(e^{y-u})\,\mu_{M,M'}(du)\\
&=\frac{M}{M'}\,\Ein_{M'}(e^y)+\Bigl(1-\frac{M}{M'}\Bigr)
\int_0^\infty Me^{-Mu}\,\Ein_{M'}(e^{y-u})\,du .
\end{align*}
\item[(b)] \emph{Unlogged race, pointwise.} For $y\ge2$,
\[
E_M(e^y)=\int_{[0,\infty)}E_{M'}(e^{y-u})\,\mu_{M,M'}(du)
+O_{\chi,M,M'}(1/y).
\]
\item[(c)] \emph{The smoothed piece is an absolutely convergent zero
sum.} As $y\to\infty$,
\begin{align*}
K_Mg_{M'}(y)&=2MM'e^{-My}\int_2^{e^y}t^{-\kappa-1}R_{m'}(t)\,dt\\
&=1+\re\sum_{\gamma>0}\frac{4MM'\,e^{i\gamma y}}{(M+i\gamma)(M'+i\gamma)}
+o(1).
\end{align*}
\item[(d)] \emph{Multiplier and semigroup.}
$\hat\mu_{M,M'}(\xi)=\dfrac{M}{M'}\cdot\dfrac{M'+i\xi}{M+i\xi}$, so
that $|\hat\mu_{M,M'}|\le1$ and
$r_\gamma(M)=\hat\mu_{M,M'}(\gamma)\,r_\gamma(M')$ for every real
$\gamma$; and $\mu_{M,M'}*\mu_{M',M''}=\mu_{M,M''}$ for $M<M'<M''$.
\end{enumerate}
\end{theorem}

\begin{proof}
(a) Riemann--Stieltjes partial summation ($R_{m'}(t)=0$ for $t<2$):
\[
R_m(x)=\int_{2^-}^xt^{-\kappa}\,dR_{m'}(t)
=x^{-\kappa}R_{m'}(x)+\kappa\int_2^xt^{-\kappa-1}R_{m'}(t)\,dt .
\]
Multiply by $2Mx^{-M}$. The first term becomes
$(M/M')\,\Ein_{M'}(x)$. In the second substitute $t=e^{y-u}$ and
$R_{m'}(t)=t^{M'}g_{M'}(\log t)/(2M')$:
\[
2M\kappa\,x^{-M}\int_2^xt^{M-1}\cdot t^{-M'}R_{m'}(t)\,dt
=\frac{\kappa}{M'}\int_0^{y-\log2}Me^{-Mu}\,g_{M'}(y-u)\,du ,
\]
and $\kappa/M'=1-M/M'$; the integral may be extended to
$u\in[0,\infty)$ because $g_{M'}=0$ below $\log2$. For $y<\log2$
both sides vanish.

(b) By \eqref{eq:fminusg}, $f_{M'}(v)-g_{M'}(v)=O(1/v)$ for $v\ge2$;
both functions are bounded on $[\log2,2]$ and vanish below $\log2$.
Hence, for $y\ge4$ (so that $y-u\ge2$ on $u\le y/2$),
\[
\Bigl|\int(f_{M'}-g_{M'})(y-u)\,\mu(du)\Bigr|
\ll\int_0^{y/2}Me^{-Mu}\,\frac2y\,du+e^{-My/2}+\frac1y\ll\frac1y ;
\]
for $2\le y<4$ the bound is trivial. Combine with (a) and with
\eqref{eq:fminusg} at weight $M$.

(c) The first equality is the substitution of (a). For the second,
multiply the formula for $\psi_{m'}(t)$ of Lemma~\ref{lem:WRshift} by
$t^{-\kappa-1}$, integrate over $[2,x]$ and let $T\to\infty$ with $x$
fixed: the truncation error tends to $0$, and since
$\rho+m'-\kappa=\rho+m$,
\[
\int_2^xt^{-\kappa-1}\psi_{m'}(t)\,dt
=-\sum_\rho\frac{x^{\rho+m}-2^{\rho+m}}{(\rho+m)(\rho+m')}
+O\bigl((1+x^m)\log^2x\bigr),
\]
the sum over all zeros converging absolutely
($\sum|\gamma|^{-2}<\infty$). After multiplication by $x^{-M}$ the
terms with $2^{\rho+m}$ and the error are $o(1)$. By
Lemma~\ref{lem:WRshift} again,
$R_{m'}(t)=-\psi_{m'}(t)+t^{M'}/(2M')+o(t^{M'})$, and
\[
x^{-M}\int_2^xt^{-\kappa-1}\Bigl(\frac{t^{M'}}{2M'}+o(t^{M'})\Bigr)dt
=\frac1{2MM'}+o(1)
\]
because $M>0$. Under GRH, $\rho+m=M+i\gamma$; pairing $\rho$ with
$\bar\rho$ gives the claim.

(d) $\int_0^\infty Me^{-Mu}e^{-i\xi u}\,du=M/(M+i\xi)$, so
$\hat\mu(\xi)=\frac{M}{M'}+(1-\frac{M}{M'})\frac{M}{M+i\xi}
=\frac{M}{M'}\cdot\frac{M'+i\xi}{M+i\xi}$, of modulus $\le1$ because
$M\le M'$, and
$\hat\mu(\gamma)\cdot4M'/(M'+i\gamma)=4M/(M+i\gamma)$. The multipliers
multiply, and a finite measure on $[0,\infty)$ is determined by its
Fourier transform.
\end{proof}

\begin{remark}[reading]\label{rem:reading}
$E_M$ is the output of the one-pole low-pass filter $K_M$ applied to
a weight-independent driving signal; passing from $M'$ to $M<M'$ is
averaging against a probability kernel supported on the past. All
the LI-free monotonicity statements of this article are consequences
of this positivity. Theorem~\ref{thm:resolvent}(d) is the precise
sense in which the family $(\nu_M)$ of limiting distributions is a
semigroup of contractions.
\end{remark}

\begin{remark}[exact unlogged form]\label{rem:exactunlogged}
For the unlogged race the same partial summation gives, exactly, for
$y>\log2$,
\[
E_M(e^y)=\frac{M}{M'}E_{M'}(e^y)+\Bigl(1-\frac{M}{M'}\Bigr)
\int_0^{y-\log2}Me^{-Mu}\cdot\frac{y}{y-u}\cdot E_{M'}(e^{y-u})\,du ;
\]
the factor $y/(y-u)=\log x/\log t$ is the source of the $O(1/y)$ in
(b). (Checked numerically, relative residual $\le1.3\cdot10^{-13}$,
by \texttt{grhonly/sieve\_checks.py}.)
\end{remark}

\begin{corollary}[$\theta$-form of {\cite[Theorem~A]{MC}}]
\label{cor:thetaA}
Assume GRH$(\chi)$ and NCZ$(\chi)$ and let $0<M<M'$. As
$x\to\infty$,
\[
\frac{\Ein_M(x)}{M}-\frac{\Ein_{M'}(x)}{M'}
=\frac{\kappa}{MM'}\,K_Mg_{M'}(\log x)
\;\ge\;\kappa\bigl(\lambda_\chi(M,M')-o(1)\bigr),
\]
and the same lower bound holds for $E_M/M-E_{M'}/M'$. So if
$\lambda_\chi(M,M')>0$ there is $x_0$ such that for all $x\ge x_0$:
$D_{m'}(x)\ge0\Rightarrow D_m(x)>0$ and
$R_{m'}(x)\ge0\Rightarrow R_m(x)>0$.
\end{corollary}

\begin{proof}
Theorem~\ref{thm:resolvent}(a), (c) and
$\bigl|\re\sum_{\gamma>0}4MM'e^{i\gamma y}/((M+i\gamma)(M'+i\gamma))\bigr|
\le1-MM'\lambda_\chi(M,M')$ (with $n_0=0$); for the unlogged race
use \eqref{eq:fminusg}.
\end{proof}

The statement for $D_m$ is Theorem~\ref{thm:MCA}(ii), proved in
\cite{MC} directly, with an error $O(x^M/\log^2x)$ and without
NCZ; nothing is added to it here except the short $\theta$-form
proof and the conclusion for $R_m$.

\section{The limiting distribution and the convex order}
\label{sec:limit}

\subsection{A mean-value lemma for all even moments}

\begin{lemma}\label{lem:moments}
Let $0<\lambda_1\le\lambda_2\le\cdots$ with
$\#\{j:t\le\lambda_j<t+1\}\le A_0\log(t+3)$ for all $t\ge0$, and let
$|c_j|\le A_1/\lambda_j$. For $3\le T<X$ put
$S(y)=\sum_{T<\lambda_j\le X}c_je^{i\lambda_jy}$. Then for every
integer $k\ge1$ and every $Y\ge1$,
\[
\frac1Y\int_0^Y|S(y)|^{2k}\,dy\;\le\;
C_k\,(A_0A_1)^{2k}\,\frac{(\log T)^{4k-1}}{T},
\]
with $C_k$ depending on $k$ only (uniformly in $X$ and $Y$).
\end{lemma}

\begin{proof}
$S^k=\sum c_{j_1}\cdots c_{j_k}e^{i\Omega y}$ with
$\Omega=\lambda_{j_1}+\dots+\lambda_{j_k}$. For integers $n\ge0$ let
$W_k(n)$ be the sum of $|c_{j_1}\cdots c_{j_k}|$ over the $k$-tuples
with $n\le\Omega<n+1$. Since
$|\frac1Y\int_0^Ye^{i(\Omega-\Omega')y}\,dy|\le
\min(1,2/(Y|\Omega-\Omega'|))\le\min(1,2/|\Omega-\Omega'|)$ for
$Y\ge1$,
\[
\frac1Y\int_0^Y|S^k|^2\,dy\le\sum_{n,n'}W_k(n)W_k(n')\,w(n-n'),
\qquad
w(d)=\begin{cases}1&(|d|\le3),\\ 2/(|d|-1)&(|d|\ge4).\end{cases}
\]
No orthogonality is used: exact coincidences $\Omega=\Omega'$,
whatever their origin, are counted with weight $1$.

\emph{Claim:} $W_k(n)\le B_k\log^{2k-1}(n+3)/(n+1)$ with
$B_k=C_k'(A_0A_1)^k$, and $W_k(n)=0$ for $n<kT-1$. For $k=1$:
$W_1(n)\le A_0\log(n+3)\cdot A_1/\max(n,T)\le2A_0A_1\log(n+3)/(n+1)$.
Inductively, a $k$-tuple with $\Omega\in[n,n+1)$ has
$\lambda_{j_1}\in[n_1,n_1+1)$ and the remaining sum in
$[n_2,n_2+1)$ with $n_1+n_2\in\{n-1,n\}$, so
\begin{align*}
W_k(n)&\le\sum_{n_1}W_1(n_1)\bigl[W_{k-1}(n-n_1)+W_{k-1}(n-1-n_1)\bigr]\\
&\le4B_1B_{k-1}\log^{2k-2}(n+3)\sum_{n_1=0}^{n}
\frac1{(n_1+1)(n-n_1+1)},
\end{align*}
and the last sum is $2H_{n+1}/(n+2)\le2(1+\log(n+1))/(n+2)$, where
$H_n$ is the $n$-th harmonic number. The vanishing for $n+1\le kT$ is
clear since every $\lambda_j>T$.

Now fix $n$; splitting $n'$ into $|n'-n|\le n/2$, $n'>3n/2$ and
$n'<n/2$ and using the claim,
$\sum_{n'}W_k(n')w(n-n')\ll B_k\log^{2k}(n+3)/(n+1)$. Hence the
double sum is
\[
\ll B_k^2\sum_{n\ge kT-1}\frac{\log^{4k-1}(n+3)}{(n+1)^2}
\ll B_k^2\,\frac{\log^{4k-1}T}{T}.\qedhere
\]
\end{proof}

\subsection{The compact group and the identification of
\texorpdfstring{$\nu_M$}{nu\_M}}

Assume GRH$(\chi)$. Let $\mathbb T$ be the unit circle,
$\iota:\mathbb R\to\mathbb T^{\cZ}$,
$\iota(y)=(e^{i\gamma y})_{\gamma\in\cZ}$, and
$\cA:=\overline{\iota(\mathbb R)}$, a compact connected metrizable
abelian group with normalised Haar measure $\hA$. Let
$e_\gamma:\cA\to\mathbb T$, $e_\gamma(\omega)=\omega_\gamma$, be the
coordinate characters and $\tau_s(\omega)=\omega\cdot\iota(s)$ the
translation flow. The characters of $\cA$ are the restrictions of
$e_{\mathbf n}(\omega)=\prod_\gamma\omega_\gamma^{n_\gamma}$,
$\mathbf n=(n_\gamma)$ a finitely supported integer vector, and
$e_{\mathbf n}(\iota(y))=e^{i\langle\mathbf n,\gamma\rangle y}$ with
$\langle\mathbf n,\gamma\rangle:=\sum_\gamma n_\gamma\gamma$, the
\emph{frequency} of $e_{\mathbf n}$.

\begin{lemma}[Kronecker--Weyl, no LI]\label{lem:KW}
\begin{enumerate}
\item[(a)] For continuous $h:\cA\to\mathbb C$,
$\frac1Y\int_0^Yh(\iota(y))\,dy\to\int_{\cA}h\,d\hA$.
\item[(b)] Two characters of $\cA$ with the same frequency coincide;
the characters $1$, $e_\gamma$, $\bar e_\gamma$ $(\gamma\in\cZ)$ are
pairwise distinct on $\cA$, hence orthonormal in $L^2(\cA)$; and the
integral over $\cA$ of a character is $1$ if its frequency is $0$
and $0$ otherwise.
\item[(c)] The flow $(\tau_s)$ is ergodic on $(\cA,\hA)$.
\item[(d)] Every nonempty open subset of $\cA$ has positive Haar
measure, and the push-forward of $\hA$ under a nontrivial character
is Haar measure on $\mathbb T$.
\end{enumerate}
\end{lemma}

\begin{proof}
If $\langle\mathbf n,\gamma\rangle=0$ then $e_{\mathbf n}\equiv1$ on
$\iota(\mathbb R)$, hence on $\cA$; if
$\langle\mathbf n,\gamma\rangle\ne0$, $e_{\mathbf n}$ is a nontrivial
character, $\int_{\cA}e_{\mathbf n}\,d\hA=0$, and the time average of
$e^{i\langle\mathbf n,\gamma\rangle y}$ tends to $0$. So (a) holds for
characters, hence for their span, which is uniformly dense in
$C(\cA)$ (Stone--Weierstrass). (b): applying the first sentence to a
quotient of two characters gives the first and third assertions;
$e_\gamma$, $e_{\gamma'}^{\pm1}$, $1$ have frequencies $\gamma$,
$\pm\gamma'$, $0$, pairwise different. (c): if $h\in L^2(\cA)$ is
flow-invariant then
$\hat h(e_{\mathbf n})(1-e^{i\langle\mathbf n,\gamma\rangle s})=0$
for all $s$, so $\hat h(e_{\mathbf n})=0$ unless
$\langle\mathbf n,\gamma\rangle=0$, i.e.\ unless $e_{\mathbf n}$ is
trivial; so $h$ is constant. (d): translation invariance and
compactness; the image of a connected group under a nontrivial
character is $\mathbb T$, and the image measure is invariant.
\end{proof}

For $T\ge2$ let
\begin{gather*}
F_{M,T}(\omega):=1+\re\sum_{\gamma\in\cZ,\ \gamma\le T}
m_\gamma\,r_\gamma(M)\,\omega_\gamma,\\
P_{M,T}(y):=F_{M,T}(\iota(y))
=1+\re\sum_{0<\gamma\le T}r_\gamma(M)\,e^{i\gamma y}.
\end{gather*}
Since $\sum_{\gamma\in\cZ}m_\gamma^2|r_\gamma(M)|^2<\infty$,
Lemma~\ref{lem:KW}(b) gives $F_{M,T}\to F_M$ in $L^2(\cA)$, where
\begin{equation}\label{eq:FM}
\begin{gathered}
F_M:=1+\re\sum_{\gamma\in\cZ}m_\gamma\,r_\gamma(M)\,e_\gamma\in L^2(\cA),
\qquad\int F_M\,d\hA=1,\\
\|F_M-1\|_2^2=\tfrac12\sum_{\gamma\in\cZ}m_\gamma^2|r_\gamma(M)|^2=V_M .
\end{gathered}
\end{equation}

\begin{theorem}\label{thm:limit}
Assume GRH$(\chi)$ and NCZ$(\chi)$, let $M>0$, $h\in\{f_M,g_M\}$, and
let $\nu_{M,Y}$ be the law of $h(y)$ for $y$ uniform on $[0,Y]$.
\begin{enumerate}
\item[(a)] For every $p<\infty$,
$\lim_{T\to\infty}\limsup_{Y\to\infty}\frac1Y\int_0^Y
|h(y)-P_{M,T}(y)|^p\,dy=0$; that is, $h$ is Besicovitch
$B^p$-almost periodic for every $p$.
\item[(b)] As $Y\to\infty$, $\nu_{M,Y}$ converges in every
Wasserstein distance $W_p$ $(p<\infty)$ to
\[
\nu_M:=\text{the law of }F_M\text{ under }\hA .
\]
Equivalently, for every continuous $\varphi$ with
$|\varphi(x)|\le C(1+|x|)^p$ for some $p$,
\[
\lim_{Y\to\infty}\frac1Y\int_0^Y\varphi\bigl(E_M(e^y)\bigr)\,dy
=\int\varphi\,d\nu_M=\int_{\cA}\varphi(F_M)\,d\hA\qquad(\text{finite}),
\]
and likewise for $\Ein_M$. In particular $F_M\in L^p(\cA)$ for all
$p<\infty$, and $\nu_M$ has mean $1$ and variance $V_M$.
\end{enumerate}
\end{theorem}

\begin{proof}
(a) It suffices to take $p=2k$. On $2\le y\le Y$ use
\eqref{eq:WRexp} with truncation $e^Y$:
$h-P_{M,T}=\re S_{T,e^Y}(y)+\cE(y,e^Y)$, where
$S_{T,X}(y)=\sum_{T<\gamma\le X}r_\gamma(M)e^{i\gamma y}$. By
Lemma~\ref{lem:moments} (with the positive ordinates listed with
multiplicity, $c_j=r_\gamma(M)$, $A_1=4M$, and
\eqref{eq:zerocount}) the first term has $2k$-th Ces\`aro mean
$\ll\log^{4k-1}T/T$ uniformly in $Y$. For the second,
$ue^{u/2}\log^2(e^ue^Y)e^{-Y}\le4Y^3e^{-Y/2}\to0$ uniformly in
$u\le Y$, and $d$ is bounded with $d(u)\to0$, so its $2k$-th Ces\`aro
mean tends to $0$. On $[0,2]$ both functions are bounded,
contributing $O(1/Y)$.

(b) Let $\nu^T_Y$ and $\nu^T$ be the laws of $P_{M,T}(y)$
($y$ uniform on $[0,Y]$) and of $F_{M,T}$ under $\hA$. By
Lemma~\ref{lem:KW}(a), $\nu^T_Y\to\nu^T$ weakly, and all these
measures live on one compact interval, so
$W_p(\nu^T_Y,\nu^T)\to0$. Coupling through the same $y$,
$W_p(\nu_{M,Y},\nu^T_Y)^p\le\frac1Y\int_0^Y|h-P_{M,T}|^p$. Hence
$\limsup_YW_p(\nu_{M,Y},\nu^T)\le\eta_T$ with $\eta_T\to0$ by (a); so
$(\nu^T)_T$ is $W_p$-Cauchy, with a limit $\nu^\infty$ (completeness
of the Wasserstein space), and $\nu_{M,Y}\to\nu^\infty$ in $W_p$. As
$F_{M,T}\to F_M$ in $L^2$, $\nu^T\to\mathrm{law}(F_M)$ in $W_2$, so
$\nu^\infty=\mathrm{law}(F_M)$. Convergence in $W_p$ is equivalent to
weak convergence together with convergence of the $p$-th absolute
moments, i.e.\ to convergence of $\int\varphi\,d\nu_{M,Y}$ for all
continuous $\varphi$ of growth $\le C(1+|x|^p)$
\cite[Theorem~6.9]{Vil09}. Finally
$\|F_M\|_{L^p}=\lim_T\|F_{M,T}\|_{L^p}<\infty$, and the mean and the
variance are read off from \eqref{eq:FM}.
\end{proof}

\begin{remark}[what is known here]\label{rem:known}
Existence of $\nu_M$ is \cite[Theorem~1(i)]{WR}, via
Akbary--Ng--Shahabi \cite{ANS14}. The representation as a
push-forward of Haar measure on the closure subtorus, without LI, is
in Rubinstein--Sarnak \cite[\S2]{RS94}, Devin \cite[\S4]{Dev20} and
Bailleul \cite[\S4]{Bai21}. Tails
$\nu(\mathbb R\setminus[-A,A])\ll\exp(-c\sqrt A)$ under GRH alone
are \cite[Theorem~1.2]{RS94} and \cite[Theorem~2.1]{Dev20} for the
unweighted race (\cite[Theorem~1.2]{RS94} bounds the mass outside a
ball; the two one-sided lower bounds are
\cite[Theorem~2.5(1)]{Dev20}). An alternative route to
$F_M\in L^p$: by the Hausdorff--Young inequality on $\cA$,
$\|F_M-1\|_{L^p(\cA)}\le\bigl(2\sum_{\gamma\in\cZ}
(m_\gamma|r_\gamma|/2)^{p'}\bigr)^{1/p'}$ for $1/p+1/p'=1$,
$p\ge2$, finite for every $p<\infty$. We claim no novelty for
Theorem~\ref{thm:limit} beyond writing out the $B^p$ statement for
the weighted races.
\end{remark}

\begin{remark}[without NCZ]\label{rem:noNCZ}
If $n_0>0$ the central zero contributes the constant $2n_0$ to
\eqref{eq:WRexp}, so one expects everything below to hold with mean
$1+2n_0$ in place of $1$ (the Cantelli bound of
Theorem~\ref{thm:cantelli} becoming $V_M/(V_M+(1+2n_0)^2)$), and
$\sigma_M^2=4M(\Lambda'/\Lambda)(M+\tfrac12)-4n_0$. We have not
rewritten the proofs of Lemmas~\ref{lem:WRshift}
and~\ref{lem:WRexp} in that generality; \cite{MC} proves
Theorem~\ref{thm:MCA} for general $n_0$. This remark is not used.
\end{remark}

\begin{remark}[constants, and what Lemma~\ref{lem:moments} cannot
give]\label{rem:noexp}
(i) The constant $C_k$ of Lemma~\ref{lem:moments} is of factorial
size in $k$ (about $(4k)!$); this is irrelevant for the theorems.
(ii) The class of functions covered by Lemma~\ref{lem:moments} does
\emph{not} have exponential moments in general: for frequencies
$n\in(T,X]$ with multiplicity $\lceil\log n\rceil$, coefficients
$1/n$ and aligned phases, $\PP(S>a)\asymp\exp(-\sqrt{2a})$ (a stress
test carried out by the referee of the working notes). This matches
the $\exp(-c\sqrt A)$ tails of Remark~\ref{rem:known} and is why no
statement about $\int e^{tx}\,d\nu_M$ beyond Theorem~\ref{thm:cx} is
made (Corollary~\ref{cor:mono}).
(iii) Dependency: Theorem~\ref{thm:cx} for every convex $\varphi$
needs only $F_M\in L^1(\cA)$, the multiplier identity and the
identification of $\nu_M$ as a weak limit (the case $k=1$ of
Lemma~\ref{lem:moments}); Lemma~\ref{lem:moments} with $k\ge2$ is
used only for the convergence of \emph{time averages} of
$\varphi(E_M(e^y))$ when $\varphi$ grows faster than quadratically.
\end{remark}

\subsection{Convex order}

\begin{theorem}[convex order; GRH and NCZ, no LI]\label{thm:cx}
Assume GRH$(\chi)$ and NCZ$(\chi)$ and let $0<M<M'$. Then
\[
F_M(\omega)=\int F_{M'}(\tau_{-u}\omega)\,\mu_{M,M'}(du)
\qquad\text{for }\hA\text{-a.e.\ }\omega,
\]
and for \emph{every} convex $\varphi:\mathbb R\to\mathbb R$,
\[
\int\varphi\,d\nu_M\;\le\;\int\varphi\,d\nu_{M'}
\qquad(\text{both sides in }(-\infty,+\infty]).
\]
For convex $\varphi$ of polynomial growth both sides are finite and
are the limits of the time averages
$\frac1Y\int_0^Y\varphi(E_M(e^y))\,dy$, resp.\ with $M'$
(Theorem~\ref{thm:limit}(b)).
\end{theorem}

\begin{proof}
Fix a Borel representative of $F_{M'}$. The map
$(\omega,u)\mapsto F_{M'}(\tau_{-u}\omega)$ is jointly measurable and
$\iint|F_{M'}(\tau_{-u}\omega)|\,\hA(d\omega)\,\mu(du)=\|F_{M'}\|_1
<\infty$ by invariance of $\hA$. So
$\Phi(\omega):=\int F_{M'}(\tau_{-u}\omega)\,\mu(du)$ is defined
a.e., lies in $L^2(\cA)$ (Minkowski), and for every character $e$ of
$\cA$, of frequency $\varpi$,
\[
\int\Phi\,\bar e\,d\hA=\int\mu(du)\,e^{-i\varpi u}\int F_{M'}\,\bar e\,d\hA
=\hat\mu(\varpi)\,\hat F_{M'}(e).
\]
By \eqref{eq:FM} and Theorem~\ref{thm:resolvent}(d) these are the
Fourier coefficients of $F_M$ (both families vanish off
$\{1,e_\gamma^{\pm1}\}$, and $\hat\mu(-\xi)=\overline{\hat\mu(\xi)}$).
Characters are complete in $L^2(\cA)$, so $\Phi=F_M$ a.e.

Let $\varphi$ be convex and $x\mapsto ax+b$ a supporting line. For
a.e.\ $\omega$, Jensen's inequality for the probability measure
$\mu$ gives
$\varphi(F_M(\omega))\le\int\varphi(F_{M'}(\tau_{-u}\omega))\,\mu(du)$.
Integrate in $\omega$, apply Tonelli's theorem to the nonnegative
function $\varphi(F_{M'}(\tau_{-u}\omega))-aF_{M'}(\tau_{-u}\omega)-b$,
and use the invariance of $\hA$: the right-hand side becomes
$\int_{\cA}\varphi(F_{M'})\,d\hA$.
\end{proof}

The time-domain version of the same proof (Jensen's inequality in
Theorem~\ref{thm:resolvent}(a), then Theorem~\ref{thm:limit}(b))
gives the inequality for convex $\varphi$ of polynomial growth
without the group; the group version removes the growth condition.

\begin{proposition}[unbounded support]\label{prop:unbounded}
Assume GRH$(\chi)$ and NCZ$(\chi)$. For every $M>0$,
$\operatorname{ess\,inf}F_M=-\infty$ and
$\operatorname{ess\,sup}F_M=+\infty$; that is,
$\nu_M((-\infty,-a))>0$ and $\nu_M((a,\infty))>0$ for every $a$.
\end{proposition}

\begin{proof}
Let $k_T(u)=\frac{T}{2\pi}\bigl(\frac{\sin(Tu/2)}{Tu/2}\bigr)^2$ be
the Fej\'er kernel on $\mathbb R$: $k_T\ge0$, $\int k_T=1$,
$\int k_T(u)e^{-i\xi u}\,du=(1-|\xi|/T)^+$. As in the previous proof,
$\int F_M(\tau_{-u}\omega)k_T(u)\,du=F^{(T)}(\omega)$ a.e., where
\[
F^{(T)}(\omega):=1+\re\sum_{\gamma\in\cZ,\ \gamma<T}
\Bigl(1-\frac\gamma T\Bigr)m_\gamma\,r_\gamma(M)\,\omega_\gamma
\]
is a \emph{continuous} function on $\cA$. Suppose $F_M\ge-a$ a.e.
Then $F^{(T)}\ge-a$ a.e., hence everywhere on $\cA$
(Lemma~\ref{lem:KW}(d)), in particular at $\omega=\iota(-1/T)$:
\begin{align*}
-a\le F^{(T)}(\iota(-1/T))
&=1+\sum_{0<\gamma<T}\Bigl(1-\frac\gamma T\Bigr)\cdot
\frac{4M\bigl[M\cos(\gamma/T)-\gamma\sin(\gamma/T)\bigr]}{M^2+\gamma^2}\\
&\le1+\frac{\sigma_M^2}{2}-4M\sum_{T/4\le\gamma\le T/2}
\frac12\,\sin\bigl(\tfrac14\bigr)\,\frac{\gamma}{M^2+\gamma^2}.
\end{align*}
For $\gamma\ge M$ one has $\gamma/(M^2+\gamma^2)\ge1/(2\gamma)\ge1/T$,
so the last sum is
$\ge\frac{\sin(1/4)}{2T}\bigl(N(T/2)-N(T/4)\bigr)\gg\log T$ by
\eqref{eq:zerocount}. Letting $T\to\infty$ gives a contradiction.
Evaluating at $\iota(+1/T)$ instead gives
$\operatorname{ess\,sup}F_M=+\infty$.
\end{proof}

\begin{lemma}\label{lem:flowunbounded}
Under the hypotheses of Proposition~\ref{prop:unbounded}, for
$\hA$-a.e.\ $\omega$ the function $u\mapsto F_M(\tau_{-u}\omega)$,
$u\ge0$, is essentially unbounded above and below (with respect to
Lebesgue measure on $[0,\infty)$).
\end{lemma}

\begin{proof}
For $a\in\mathbb N$ let
$N_a=\{\omega:F_M(\tau_{-u}\omega)\ge-a\text{ for a.e.\ }u\ge0\}$.
Then $\tau_{-s}N_a\subset N_a$ for $s\ge0$, and $\tau_{-s}$ preserves
$\hA$, so $N_a$ is invariant modulo null sets; by ergodicity
(Lemma~\ref{lem:KW}(c)) $\hA(N_a)\in\{0,1\}$. If it were $1$, Fubini's
theorem and invariance would give
$0=\iint1\{F_M(\tau_{-u}\omega)<-a\}e^{-u}\,du\,\hA(d\omega)
=\hA(F_M<-a)$, contradicting Proposition~\ref{prop:unbounded}. The
upper side is the same; take the countable union.
\end{proof}

\begin{theorem}[strictness]\label{thm:strict}
Assume GRH$(\chi)$ and NCZ$(\chi)$, let $0<M<M'$, and let $\varphi$
be convex, not affine, with $\int\varphi\,d\nu_{M'}<\infty$. Then
$\int\varphi\,d\nu_M<\int\varphi\,d\nu_{M'}$.
\end{theorem}

\begin{proof}
Suppose equality. Then Jensen's inequality in the proof of
Theorem~\ref{thm:cx} is an equality for a.e.\ $\omega$. Fix such an
$\omega$, put $g(u)=F_{M'}(\tau_{-u}\omega)$, $x_0=\int g\,d\mu$, and
let $\alpha$ be a supporting line of $\varphi$ at $x_0$. Then
$\int(\varphi-\alpha)(g)\,d\mu=0$ with $\varphi-\alpha\ge0$
continuous, so $\varphi=\alpha$ on the essential range of $g$ with
respect to $\mu$; as $\mu$ has a positive density on $(0,\infty)$,
this contains the Lebesgue-essential range, whose convex hull is
$\mathbb R$ for a.e.\ $\omega$ (Lemma~\ref{lem:flowunbounded} for
$M'$). The zero set of the convex function $\varphi-\alpha\ge0$ is an
interval, so $\varphi=\alpha$ on $\mathbb R$, a contradiction.
\end{proof}

\begin{corollary}[monotone functionals]\label{cor:mono}
Assume GRH$(\chi)$ and NCZ$(\chi)$. Each of the following is a
strictly increasing function of $M\in(0,\infty)$, and is the limit of
the corresponding logarithmic time average of the race $E_M$ (or
$\Ein_M$) (Theorem~\ref{thm:limit}(b)):
\begin{itemize}
\item the variance
$V_M=\lim_Y\frac1Y\int_0^Y(E_M(e^y)-1)^2\,dy$;
\item every absolute central moment $\int|x-1|^p\,d\nu_M$, $p\ge1$,
and every $\int|x-c|^p\,d\nu_M$;
\item the \emph{mean deficit}
$\int x^-\,d\nu_M=\lim_Y\frac1Y\int_0^Y\max(-E_M(e^y),0)\,dy$, and
more generally $\int(c-x)^+\,d\nu_M$ and $\int(x-c)^+\,d\nu_M$ for
every real $c$.
\end{itemize}
For real $t\ne0$, $M\mapsto\int e^{tx}\,d\nu_M\in(0,\infty]$ is
nondecreasing, and $\int e^{tx}\,d\nu_M<\int e^{tx}\,d\nu_{M'}$ for
$M<M'$ whenever the right-hand side is finite. This is a statement
about $\nu_M$ only: no time-average statement is made, and under GRH
alone we do not know that this integral is finite for any $t\ne0$
(Remark~\ref{rem:noexp}(ii)). Moreover, for convex
$\Psi\ge1_{(-\infty,0]}$ and $M<M'$,
\begin{equation}\label{eq:transport}
\nu_M((-\infty,0])\;\le\;\int\Psi\,d\nu_M\;\le\;\int\Psi\,d\nu_{M'} ;
\end{equation}
so a bound for $\nu_M((-\infty,0])$ obtained by transporting a convex
test function from a heavier weight is never better than the same
test function at $M$.
\end{corollary}

\begin{proof}
Theorems~\ref{thm:cx}, \ref{thm:strict} and~\ref{thm:limit}(b): the
functions $|x-c|^p$, $(c-x)^+$, $(x-c)^+$ are convex, non-affine and
of polynomial growth; $e^{tx}$ is convex and non-affine, and
Theorems~\ref{thm:cx}, \ref{thm:strict} apply to it directly. The
display \eqref{eq:transport} is $1_{(-\infty,0]}\le\Psi$ and
Theorem~\ref{thm:cx}.
\end{proof}

\begin{remark}\label{rem:nogain}
Corollary~\ref{cor:mono} answers the question ``can one improve a
bound at weight $m$ by passing through another weight?'' negatively
for all convex test functions (Cantelli's inequality is one): the
bound computed at $M$ itself is the best in the family.
\end{remark}

\section{Both lead sets have positive lower density}
\label{sec:leadsets}

\begin{theorem}\label{thm:leadsets}
Assume GRH$(\chi)$ and NCZ$(\chi)$ (no LI). For every $m>-\tfrac12$,
\begin{align*}
\liminf_{X\to\infty}\frac1{\log X}\int_2^X1\{D_m(t)<0\}\,\frac{dt}t
&\;\ge\;\nu_M((-\infty,0))\;>\;0,\\
\liminf_{X\to\infty}\frac1{\log X}\int_2^X1\{D_m(t)>0\}\,\frac{dt}t
&\;\ge\;\nu_M((0,\infty))\;\ge\;\frac1{1+V_M}\;>\;0,
\end{align*}
and the same holds for $R_m$. Consequently both upper logarithmic
densities are $<1$. Quantitatively, for every $\gamma\in\cZ$ with
$m_\gamma c_\gamma(M)>2$ (for instance the least ordinate $\gamma_1$
when $M>\gamma_1/\sqrt3$),
\[
\int x^-\,d\nu_M\;\ge\;\frac{m_\gamma c_\gamma(M)-2}{4},
\qquad
\nu_M((-\infty,0))\;\ge\;\frac{(m_\gamma c_\gamma(M)-2)^2}{16\,V_M}.
\]
\end{theorem}

\begin{proof}
For $x>1$, $\{D_m<0\}=\{E_M<0\}$; the set $(-\infty,0)$ is open, so
the portmanteau theorem (Section~\ref{subsec:densities}) and
Proposition~\ref{prop:unbounded} give the first line; the passage
between $\frac1{\log X}\int_2^X\cdots\,dt/t$ and
$\frac1Y\int_0^Y\cdots\,dy$, $Y=\log X$, costs $O(1/Y)$. The second
line is Theorem~\ref{thm:cantelli} below. For the quantitative part:
$m_\gamma c_\gamma/2=|\hat F_M(e_\gamma)|\le\|F_M\|_1
=1+2\int F_M^-\,d\hA$, and
\[
\int F_M^-\,d\hA\le\bigl(\hA(F_M<0)\bigr)^{1/2}
\Bigl(\int(F_M^-)^2\,d\hA\Bigr)^{1/2}
\le\bigl(\hA(F_M<0)\,V_M\bigr)^{1/2},
\]
because $|F_M|\le|F_M-1|$ where $F_M<0$.
\end{proof}

At $m=0$ the positivity is \cite[Theorem~2.5(1)]{Dev20}
($0<\dlow\le\dup<1$, with one-sided tails $\gg\exp(-\exp(ca))$; this
generalises \cite[Theorem~1.2]{RS94}, which bounds the mass outside a
ball) and Kaczorowski \cite{Kac93} (natural lower density); for
$m\ne0$ we did not find it stated. The proof above is qualitative:
the lower bound for $\nu_M((-\infty,0))$ is explicit only for
$M>\gamma_1/\sqrt3$ ($m>2.98$ for $q=4$).

\section{The variance without simple zeros, and explicit control of
multiplicities}\label{sec:variance}

\begin{theorem}\label{thm:var}
Assume GRH$(\chi)$ and NCZ$(\chi)$. For every $M>0$,
\begin{gather*}
\operatorname{Var}(\nu_M)
=\lim_{Y\to\infty}\frac1Y\int_0^Y\bigl(E_M(e^y)-1\bigr)^2\,dy
=V_M=\sum_{\gamma\in\cZ}m_\gamma^2\cdot\frac{8M^2}{M^2+\gamma^2},\\
V_M=\sigma_M^2+\sum_{\gamma\in\cZ}(m_\gamma^2-m_\gamma)\,w_M(\gamma)
\;\ge\;\sigma_M^2=4M\,\frac{\Lambda'}{\Lambda}\bigl(M+\tfrac12,\chi\bigr),
\end{gather*}
with equality if and only if every zero of $L(s,\chi)$ is simple.
The function $M\mapsto V_M$ is strictly increasing, and
$V_M\le C_2(\chi)(2M)^2$ with
$C_2(\chi)=\sum_{\gamma\in\cZ}2m_\gamma^2/\gamma^2<\infty$.
\end{theorem}

\begin{proof}
\eqref{eq:FM} and Theorem~\ref{thm:limit}(b);
Proposition~\ref{prop:WRvar} for $\sigma_M^2$;
Section~\ref{subsec:zeros} (or Lemma~\ref{lem:mult}) for finiteness;
each $w_M(\gamma)$ is strictly increasing in $M$ and
$w_M(\gamma)\le8M^2/\gamma^2$.
\end{proof}

The closed form $\sigma_M^2$ counts the zeros with weight $m_\gamma$;
the variance of the race counts them with weight $m_\gamma^2$. In
general form this is Fiorilli \cite[Lemma~2.5]{Fio14} and Devin
\cite[Theorem~2.1]{Dev20}. Under LI the zeros are simple and the two
quantities coincide; this is the setting of the density statements of
\cite{WR}. Under GRH alone simplicity is not known, so $V_M$, not
$\sigma_M^2$, must be used, and \emph{numerical} statements need an
explicit bound for the excess $V_M-\sigma_M^2$. The rest of this
section provides one.

\begin{lemma}[explicit multiplicity bound]\label{lem:mult}
Assume GRH$(\chi)$ and NCZ$(\chi)$. Let $\chi(-1)=(-1)^a$,
$\eta>\tfrac12$, $w_0=(\tfrac12+\eta+a)/2$ and
$Z(\sigma)=-\zeta'(\sigma)/\zeta(\sigma)$. For every real $t$,
\[
\sum_\rho\frac{\eta^2}{\eta^2+(t-\gamma_\rho)^2}
\;\le\;\mathfrak m_\eta(t):=\frac\eta4\log\bigl(t^2+4w_0^2\bigr)+\eta\,b_\eta,
\qquad
b_\eta:=\tfrac12\log\frac{q}{2\pi}+\frac1{24w_0^2}+Z\bigl(\tfrac12+\eta\bigr),
\]
the sum being over all zeros with multiplicity. In particular
$m_\gamma\le\mathfrak m_\eta(\gamma)$ for every ordinate $\gamma$.
\end{lemma}

\begin{proof}
By \eqref{eq:XiLogDer} with $z=\eta+it$,
$\re(\Lambda'/\Lambda)(\tfrac12+\eta+it)
=\sum_\rho\eta/(\eta^2+(t-\gamma_\rho)^2)$. On the other hand
$\frac{\Lambda'}{\Lambda}(s)=\tfrac12\log\frac q\pi
+\tfrac12\frac{\Gamma'}{\Gamma}\bigl(\frac{s+a}2\bigr)
+\frac{L'}{L}(s,\chi)$, $|\frac{L'}{L}(s,\chi)|\le Z(\re s)$ for
$\re s>1$, and for $\re w=w_0>0$ Binet's formula
\[
\frac{\Gamma'}{\Gamma}(w)=\log w-\frac1{2w}
-\int_0^\infty\Bigl(\frac12-\frac1v+\frac1{e^v-1}\Bigr)e^{-vw}\,dv,
\qquad
0\le\frac12-\frac1v+\frac1{e^v-1}\le\frac v{12},
\]
gives $\re\frac{\Gamma'}{\Gamma}(w)\le\log|w|+1/(12w_0^2)$. With
$w=(\tfrac12+\eta+a+it)/2$,
$\log|w|=\tfrac12\log(t^2+4w_0^2)-\log2$. Multiply by $\eta$.
\end{proof}

For $q=4$ and $\eta=0.9$, $\mathfrak m_\eta(640.17)=4.52$, so zeros
just above $T_4$ have multiplicity at most $4$ (for $q=3$: $4.43$ at
$T_3$) \certA. (The explicit zero-counting bound of
Bennett--Martin--O'Bryant--Rechnitzer \cite{BMOR21},
$|N(T,\chi)-\text{main term}|\le0.22737\,l+2\log(1+l)-0.5$ with
$l=\log(q(T+2)/2\pi)$, where $N(T,\chi)$ counts the zeros with
$|\gamma|\le T$ so that an ordinate of multiplicity $m_\gamma$ makes
it jump by $2m_\gamma$, gives $2m_\gamma\le9.53$ at $T_4$ and
$\le9.32$ at $T_3$, i.e.\ $m_\gamma\le4$, unconditionally, at the
same heights. Lemma~\ref{lem:mult} gives the same integer bound, is
self-contained under GRH, and has the smooth form needed in
Proposition~\ref{prop:Vplus}.)

\begin{lemma}[a log-weighted zero sum in closed form]\label{lem:JA}
Assume GRH$(\chi)$ and NCZ$(\chi)$. For $A>0$,
\[
J_A:=\sum_{\gamma>0}\gamma^{-2}\log\Bigl(1+\frac{\gamma^2}{A^2}\Bigr)
=\int_A^\infty\frac{\Lambda'}{\Lambda}\bigl(s+\tfrac12,\chi\bigr)\,
\frac{ds}{s^2},
\]
and, with $c_a=\tfrac12+a$,
\[
J_A\le J_A^+:=\frac{\log(q/2\pi)}{2A}
+\frac12\Bigl[\frac{\log(A+c_a)}{A}+\frac{\log(1+c_a/A)}{c_a}\Bigr]
-\frac12\Bigl[\frac1{c_aA}-\frac{\log(1+c_a/A)}{c_a^2}\Bigr]
+\frac{Z(A+\tfrac12)}{A},
\]
with $J_A^+-J_A\le1/(18A^3)+2Z(A+\tfrac12)/A$.
\end{lemma}

\begin{proof}
$\int_A^\infty2\,ds/(s(s^2+\gamma^2))
=\gamma^{-2}\log(1+\gamma^2/A^2)$; sum over $\gamma$ (Tonelli) and
use $\sum_{\gamma>0}2s/(s^2+\gamma^2)=(\Xi'/\Xi)(s)$. For the bound:
$\frac{\Gamma'}{\Gamma}(x)\le\log x-1/(2x)$ (and
$\ge\log x-1/(2x)-1/(12x^2)$) for $x>0$,
$|\frac{L'}{L}(s+\tfrac12,\chi)|\le Z(A+\tfrac12)$ for $s\ge A$, and
the elementary integrals
$\int_A^\infty\log(s+c)s^{-2}\,ds=\log(A+c)/A+\log(1+c/A)/c$,
$\int_A^\infty ds/(s^2(s+c))=1/(cA)-\log(1+c/A)/c^2$.
\end{proof}

\begin{proposition}[certified variance bounds]\label{prop:Vplus}
Assume GRH$(\chi)$ and NCZ$(\chi)$, and that every zero with
$0<\gamma\le T$ is simple. Let $A\ge2w_0$, $\eta>\tfrac12$, and
\begin{gather*}
\mathrm{tail}_2:=\frac{C(\chi)}2-\sum_{\gamma\le T}\gamma^{-2},
\qquad
\mathrm{tail}_{\log}:=J_A^+-\sum_{\gamma\le T}\gamma^{-2}
\log\Bigl(1+\frac{\gamma^2}{A^2}\Bigr),\\
\tau_M:=\sigma_M^2-\sum_{\gamma\le T}w_M(\gamma),
\qquad
c':=\eta\Bigl(\frac{\log A}2+b_\eta\Bigr)-1 .
\end{gather*}
Then
\[
V_M\le V_M^+:=\sigma_M^2+8M^2\,\frac\eta4\,\mathrm{tail}_{\log}
+c'\,\tau_M,
\qquad
C_2(\chi)\le C^+(\chi):=C(\chi)+2\Bigl[\frac\eta4\,\mathrm{tail}_{\log}
+c'\,\mathrm{tail}_2\Bigr].
\]
(For $A=30$, $\eta=0.9$: $c'=2.1447$ for $\chi_4$ and $2.0152$ for
$\chi_3$.)
\end{proposition}

\begin{proof}
$V_M-\sigma_M^2=\sum_{\gamma>T}(m_\gamma-1)\,w_M(\gamma)$ (sum with
multiplicity; the zeros up to $T$ are simple). By
Lemma~\ref{lem:mult} and
$\log(\gamma^2+4w_0^2)\le\log A^2+\log(1+\gamma^2/A^2)$,
$m_\gamma-1\le\frac\eta4\log(1+\gamma^2/A^2)+c'$. Use
$w_M(\gamma)\le8M^2/\gamma^2$ in the first term; the sums over
$\gamma>T$ of $\gamma^{-2}\log(1+\gamma^2/A^2)$, of $\gamma^{-2}$ and
of $w_M(\gamma)$ are $\le\mathrm{tail}_{\log}$, $=\mathrm{tail}_2$
and $=\tau_M$ respectively, by Lemma~\ref{lem:JA} and
Proposition~\ref{prop:WRvar}. The bound for $C_2$ is the same
computation with $2/\gamma^2$ in place of $w_M(\gamma)$.
\end{proof}

\noindent\textbf{Numerical values} \certA{}
(\texttt{grhonly/certified\_bounds.py}, Arb, $A=30$, $\eta=0.9$;
output in \texttt{certified\_bounds\_results.txt}):

\begin{center}
\begin{tabular}{lll}
\toprule
 & $\chi_4$ ($T=640.1657$) & $\chi_3$ ($T=700.0247$)\\
\midrule
$C(\chi)$ & $0.15602964988964488461$ & $0.11340215685275857166$\\
$\mathrm{tail}_2=\sum_{\gamma>T}\gamma^{-2}$ & $0.00174181939635$ &
$0.00154713475655$\\
$\mathrm{tail}_{\log}\le$ & $0.0146477305$ & $0.0133016045$\\
excess $\sum_{\gamma>T}(m_\gamma^2-m_\gamma)/\gamma^2\le$ &
$0.0070314280$ & $0.0061107212$\\
$C^+(\chi)$ & $0.170092506$ & $0.125623600$\\
$\sigma_{1/2}^2$ (closed form) & $0.15556797992358593$ &
$0.11322996985747234$\\
$\tau_{1/2}(T)$ (exact tail) & $0.0034836382$ & $0.0030942690$\\
$V_{1/2}^+$ & $\mathbf{0.1696308345}$ & $\mathbf{0.1254514113}$\\
\bottomrule
\end{tabular}
\end{center}

\noindent(In the rows $\mathrm{tail}_2$ and ``excess'' the sum
$\sum_{\gamma>T}$ is over ordinates with multiplicity in the first
and over distinct ordinates in the second.) So
$\sigma_{1/2}^2\le V_{1/2}\le1.0904\,\sigma_{1/2}^2$ for $\chi_4$ and
$V_{1/2}\le1.1080\,\sigma_{1/2}^2$ for $\chi_3$, under GRH alone; and
$V_M\le C^+(\chi)(2M)^2$ for all $M>0$.

\begin{proposition}[tier B]\label{prop:tierB}
Assume GRH$(\chi)$ and NCZ$(\chi)$, and that all zeros with
$0<\gamma\le X$ are simple, where $X\ge2w_0$; let $\eta>\tfrac12$
with $\eta\,b_\eta\ge1$. Then $V_M\le\sigma_M^2+8M^2E_X$, where
\begin{gather*}
E_X:=\frac\eta4\bigl(J_X^++2\log X\cdot U_X\bigr)+(\eta\,b_\eta-1)\,U_X,\\
U_X:=\frac1X\Bigl[\tfrac12\log\frac q\pi
+\tfrac12\log\frac{X+c_a}2+Z\bigl(X+\tfrac12\bigr)\Bigr]
\;\ge\;\sum_{\gamma>X}\gamma^{-2}.
\end{gather*}
\end{proposition}

\begin{proof}
For $\gamma>X$: $\gamma^{-2}\le2/(X^2+\gamma^2)$ and
$\sum_{\gamma>0}2/(X^2+\gamma^2)
=(\Lambda'/\Lambda)(X+\tfrac12)/X\le U_X$;
$\log(\gamma^2+4w_0^2)\le\log X^2+\log(1+\gamma^2/X^2)$; then
Lemmas~\ref{lem:mult} and~\ref{lem:JA}. (For the parameters used,
$\eta\,b_\eta-1=0.614$ for $\chi_4$ and $0.485$ for $\chi_3$.)
\end{proof}

With $X=10^8/q$ and $\eta=0.9$: $E_X\le2.83\cdot10^{-6}$ ($\chi_4$),
$2.12\cdot10^{-6}$ ($\chi_3$), giving
$V_{1/2}(\chi_4)\le0.1555736302$ and
$V_{1/2}(\chi_3)\le0.1132342076$ \certB.

\begin{remark}[status of tier B]\label{rem:tierB}
Platt \cite{Pla16} verifies GRH for all primitive characters of
modulus $q\le400\,000$ to height $\max(10^8/q,\dots)$; Theorem~7.1
there states only GRH. The method (\S3 there) locates sign changes of
the completed $L$-function and matches their number against Turing's
count \emph{with multiplicity}, which forces every zero in the range
to be simple. The same inference is made explicitly in
\cite[\S2.1]{PT26} (``\dots lie on the $1/2$ line (and are simple)'')
for a different computation. We therefore state the \certB{}
constants as conditional on the hypothesis ``all zeros of
$L(s,\chi_q)$ up to height $10^8/q$ are simple'', which we believe is
a by-product of Platt's computation but which we have not seen stated
as a theorem. Tier~A does not depend on it.
\end{remark}

\begin{remark}[how to close the gap between the tiers]
\label{rem:closegap}
The tier-A excess $0.0141$ (for $\chi_4$, $M=\tfrac12$) is about
$4.04\,\tau_{1/2}$ and decays like $\log^2T/T$; certifying the zero
list (with simplicity) to $T=10^4$ would reduce it to about $0.002$.
An explicit pair-correlation bound would do the same without more
zeros; neither was attempted.
\end{remark}

\section{Bias bounds without LI}\label{sec:biasbounds}

\subsection{Cantelli: a one-sided Chebyshev bound}

\begin{theorem}\label{thm:cantelli}
Assume GRH$(\chi)$ and NCZ$(\chi)$ (no LI). For every $m>-\tfrac12$,
\begin{align*}
\liminf_{X\to\infty}\frac1{\log X}\int_2^X1\{D_m(t)>0\}\,\frac{dt}t
&\;\ge\;\frac1{1+V_M},\\
\limsup_{X\to\infty}\frac1{\log X}\int_2^X1\{D_m(t)\le0\}\,\frac{dt}t
&\;\le\;\nu_M((-\infty,0])\;\le\;\frac{V_M}{1+V_M},
\end{align*}
and the same holds for $R_m$. In particular the lower logarithmic
density of $\{D_m>0\}$ exceeds $\tfrac12$ whenever $V_M<1$, and
tends to $1$ as $m\to-\tfrac12^+$, at least as fast as
$1-C_2(\chi)(2m+1)^2$.
\end{theorem}

\begin{proof}
For $x>1$, $D_m(x)\le0\iff E_M(x)\le0$. The set $(-\infty,0]$ is
closed, so by Theorem~\ref{thm:limit}(b) and the portmanteau theorem
$\limsup_Y\nu_{M,Y}((-\infty,0])\le\nu_M((-\infty,0])$; the passage
between $\frac1{\log X}\int_2^X\cdots\,dt/t$ and
$\frac1Y\int_0^Y\cdots\,dy$, $Y=\log X$, costs $O(1/Y)$. The measure
$\nu_M$ has mean $1$ and variance $V_M$ (Theorem~\ref{thm:var}), and
Cantelli's one-sided inequality
$\PP(Z-\EE Z\le-t)\le\operatorname{Var}Z/(\operatorname{Var}Z+t^2)$
at $t=1$ gives the bound. Ties $D_m=0$ are counted on the losing
side throughout. The last assertion uses $V_M\le C_2(\chi)(2M)^2$.
\end{proof}

The two-sided (Chebyshev) form of this bound, $\dlow\ge1-V$ (resp.\
$1-2V/\mathrm{mean}^2$), with the multiplicity-squared variance, is
Devin \cite[Corollary~5.8]{Dev20} and Fiorilli
\cite[Lemma~2.7]{Fio14}; the gain from the one-sided inequality is
$1-V\to1/(1+V)$. What is added here is the weighted family, the
explicit control of $V_M$ under GRH alone, and the evaluation.

\begin{theorem}[Chebyshev's race, explicit; GRH only]
\label{thm:cantelli-q}
\begin{enumerate}
\item[(a)] Assume GRH$(\chi_4)$. The set
$\{x:\pi(x;4,3)>\pi(x;4,1)\}$ has lower logarithmic density
\begin{itemize}
\item[] $\ge1/(1+0.1696308345)>0.8549$ \certA;
\item[] $\ge0.8653$ if moreover all zeros of $L(s,\chi_4)$ up to
height $2.5\cdot10^7$ are simple \certB;
\item[] $\ge1/(1+\sigma_{1/2}^2)=0.865375\ldots$ if all zeros of
$L(s,\chi_4)$ are simple.
\end{itemize}
\item[(b)] Assume GRH$(\chi_3)$. The set
$\{x:\pi(x;3,2)>\pi(x;3,1)\}$ has lower logarithmic density
$>0.8885$ \certA; $\ge0.8982$ \certB; $\ge0.898286\ldots$ if all
zeros of $L(s,\chi_3)$ are simple.
\item[(c)] For $q\in\{4,3\}$, under GRH$(\chi_q)$, and all
$m>-\tfrac12$:
$\dlow\{D_m>0\}\ge1/\bigl(1+C^+(\chi_q)(2m+1)^2\bigr)$, with
$C^+(\chi_4)=0.170093$ and $C^+(\chi_3)=0.125624$ \certA.
\end{enumerate}
\end{theorem}

\begin{proof}
Theorem~\ref{thm:cantelli} and
Propositions~\ref{prop:Vplus}, \ref{prop:tierB}; NCZ is unconditional
for $q=3,4$ by \eqref{eq:Lpositive}; the zero lists are complete and
simple up to $T_4$, $T_3$ (Section~\ref{subsec:tags}).
\end{proof}

Certified values of $1/(1+V_M^+)$ \certA{} and, in brackets, the
soft-threshold refinement of Theorem~\ref{thm:peel} (all values
rounded down):

\begin{center}
\small
\begin{tabular}{lllll}
\toprule
$m$ & $-0.45$ & $-0.4$ & $-0.3$ & $-0.25$\\
\midrule
$\chi_4$ & $0.998302$ $[0.999497]$ & $0.993243$ $[0.996494]$ &
$0.973517$ $[0.978593]$ & $0.959238$ $[0.963911]$\\
$\chi_3$ & $0.998745$ $[0.999593]$ & $0.995000$ $[0.997270]$ &
$0.980300$ $[0.983745]$ & $0.969560$ $[0.972670]$\\
\midrule
$m$ & $-0.1$ & $0$ & $0.25$ & $0.5$\\
\midrule
$\chi_4$ & $0.901981$ $[0.903337]$ & $0.854970$ & $0.724433$ &
$0.597688$\\
$\chi_3$ & $0.925644$ $[0.926460]$ & $0.888532$ & $0.780162$ &
$0.666771$\\
\bottomrule
\end{tabular}
\end{center}

\noindent The bound $1/(1+V_M^+)$ stays above $\tfrac12$ for
$m<0.722$ ($\chi_4$) and $m<0.918$ ($\chi_3$) \numtag{} (from the
\certA{} excess); if all zeros are simple it does so for $m<0.7778$,
resp.\ $m<0.9947$ (where $\sigma_M^2=1$).

\begin{theorem}[peeling off an absolutely convergent head]
\label{thm:peel}
Assume GRH$(\chi)$ and NCZ$(\chi)$. Let $S\subset\cZ$ be finite and
$0\le s_\gamma\le m_\gamma c_\gamma(M)$ for $\gamma\in S$, with
$B:=\sum_{\gamma\in S}s_\gamma<1$. Then
\[
\nu_M((-\infty,0])\;\le\;\frac{V'}{V'+(1-B)^2},
\qquad
V':=V_M-\frac12\sum_{\gamma\in S}\Bigl[(m_\gamma c_\gamma)^2
-(m_\gamma c_\gamma-s_\gamma)^2\Bigr].
\]
\end{theorem}

\begin{proof}
Write $F_M=1+H_S+U$ with
$H_S:=\re\sum_{\gamma\in S}s_\gamma u_\gamma e_\gamma$ and
$u_\gamma=r_\gamma/|r_\gamma|$. Then $|H_S|\le B$ everywhere, $U$ has
mean $0$ and variance $V'$ by orthogonality
(Lemma~\ref{lem:KW}(b)), and $\{F_M\le0\}\subset\{U\le-(1-B)\}$.
Apply Cantelli's inequality to $U$.
\end{proof}

The optimal choice is soft thresholding,
$s_\gamma=(m_\gamma c_\gamma-\tau)^+$. To first order there is a gain
if and only if the largest amplitude exceeds $2V$ (the derivative of
the bound in $s$ at $s=0$ is $(2V-m_\gamma c_\gamma)/(1+V)^2$). At
$m=0$ in tier~A it gains nothing ($c_{\gamma_1}(\tfrac12)=0.3310
<2V^+=0.3393$ for $\chi_4$); with simple zeros
$2\sigma_{1/2}^2=0.3111<0.3310$ and there is a slight gain. It wins
near $m=-\tfrac12$ (bracketed values in the table above), and
asymptotically it changes the order of magnitude:

\begin{proposition}\label{prop:endpoint}
Assume GRH$(\chi)$ and NCZ$(\chi)$. There are $c,C>0$ depending on
$\chi$ such that for $0<M\le\tfrac12$
\[
\dup\{x:D_m(x)\le0\}\;\le\;\nu_M((-\infty,0])\;\le\;
C\exp\bigl(-c\,M^{-1/2}\bigr).
\]
\end{proposition}

\begin{proof}
Take $S=\{\gamma\in\cZ:\gamma\le G\}$ and $s_\gamma=m_\gamma c_\gamma$
in Theorem~\ref{thm:peel}. By \eqref{eq:zerocount},
$B\le4M\sum_{0<\gamma\le G}1/\gamma\le C_3M\log^2(G+3)$, and by
Lemma~\ref{lem:mult} and \eqref{eq:zerocount},
\[
V'=\sum_{\gamma\in\cZ,\ \gamma>G}m_\gamma^2\,w_M(\gamma)
\le8M^2\sum_{\gamma>G}\frac{\mathfrak m_\eta(\gamma)}{\gamma^2}
\le C_4\,M^2\,\frac{\log^2(G+3)}{G}.
\]
Choose $G=\exp((2C_3M)^{-1/2})-3$ (for $M$ small; adjust $C$
otherwise): then $B\le\tfrac12$ and the bound is
$\le4V'\le(2C_4/C_3)\,M\exp(-(2C_3M)^{-1/2})\,(1+o(1))$. The first
inequality of the statement is Theorem~\ref{thm:cantelli}.
\end{proof}

Cantelli's inequality alone gives only $O(M^2)$. Under LI one has the
upper bound $\exp(-1/(8C(\chi)M^2))$ \cite[Theorem~1(iii)]{WR}; the
true order under LI is expected to be much smaller still, doubly
exponential in $M^{-1/2}$ (a heuristic, not proved here, based on
$\sum_\gamma c_\gamma(M)=\infty$ with
$c_\gamma\approx4M/\gamma$). At the endpoint
$m=-\tfrac12$ Hayani \cite{Hay25} and Sheth \cite{Sheth24} prove
density one under GRH.

\subsection{Beyond Cantelli without LI: finite non-resonance}
\label{subsec:NR}

LI is an infinite statement, but \emph{the absence of integer
relations of bounded order among finitely many certified zeros} is a
finite statement, verifiable (when true) by comparing enclosures; an
exact relation, if one existed, could not be detected this way. The
following theorem converts such a certificate into a density bound.

\emph{Setting.} Assume GRH$(\chi)$ and NCZ$(\chi)$. Let
$\gamma_1<\dots<\gamma_N$ be the $N$ smallest ordinates and
$T\ge3\gamma_N$ a height such that the list of zeros with
$0<\gamma\le T$ is \emph{complete} (this is used in (NR2), (NR3)
below) and all of them are simple. Put $T'=T-2\gamma_N$,
$\cK=\cZ\cap(\gamma_N,T']$, $a_n=c_{\gamma_n}(M)$, and
\[
V_{\mathrm h}=\tfrac12\sum_{n\le N}a_n^2,\qquad
V_{\cK}=\tfrac12\sum_{\gamma\in\cK}c_\gamma(M)^2,\qquad
V_{\mathrm t}=V_M-V_{\mathrm h}-V_{\cK}.
\]
For an integer vector $\mathbf n=(n_1,\dots,n_N)$ write
$\mathbf n\cdot\gamma=\sum_in_i\gamma_i$ and
$|\mathbf n|_1=\sum_i|n_i|$. The \emph{non-resonance conditions}
are:
\begin{itemize}
\item[(NR1)] $\mathbf n\cdot\gamma\ne0$ for all $\mathbf n\ne0$ with
$|\mathbf n|_1\le4$;
\item[(NR2)] for $|\mathbf n|_1\le3$: $|\mathbf n\cdot\gamma|$ is not
an ordinate outside $\{\gamma_1,\dots,\gamma_N\}$;
\item[(NR3)] for $1\le|\mathbf n|_1\le2$:
$\mathbf n\cdot\gamma\ne\pm\gamma'\pm\gamma''$ for all ordinates
$\gamma',\gamma''\in(\gamma_N,T]$.
\end{itemize}

\begin{theorem}[finite non-resonance bound; GRH and NCZ, no LI]
\label{thm:NR}
In the setting above assume \textup{(NR1)--(NR3)}. Then for all
$\varkappa_0>0$ and $b_1,\dots,b_N\ge0$ with $\sum_nb_n\le1$,
\begin{multline*}
\nu_M((-\infty,0])\;\le\;J(\varkappa_0,b):=(1-\varkappa_0)^2
+(1-\varkappa_0)\varkappa_0\sum_na_nb_n\\
+\varkappa_0^2\Bigl[\tfrac12\sum_n(a_n-b_n)^2
+\tfrac14\Bigl(\sum_na_n^2\Bigr)\Bigl(\sum_nb_n^2\Bigr)
+\tfrac12\Bigl(\sum_na_nb_n\Bigr)^2-\tfrac38\sum_na_n^2b_n^2\\
+\Bigl(1+\tfrac12\sum_nb_n^2\Bigr)V_{\cK}
+\Bigl(1+\sum_nb_n\Bigr)^2V_{\mathrm t}\Bigr],
\end{multline*}
and therefore $\dlow\{D_m>0\}\ge1-\inf J$. (For $b=0$,
$\inf_{\varkappa_0}J=V_M/(1+V_M)$: Cantelli.)
\end{theorem}

\begin{proof}
Let $\vartheta_n:=\arg\bigl(r_{\gamma_n}e_{\gamma_n}(\omega)\bigr)$,
so that
\begin{gather*}
F_M=1+H+W_{\cK}+W_{\mathrm t},\qquad
H=\sum_{n\le N}a_n\cos\vartheta_n,\\
W_{\cK}=\re\sum_{\gamma\in\cK}r_\gamma e_\gamma,\qquad
W_{\mathrm t}=\re\sum_{\gamma\in\cZ,\ \gamma>T'}m_\gamma r_\gamma e_\gamma
\end{gather*}
(the zeros in $(T',T]$ are put into $W_{\mathrm t}$), with
$\|W_{\cK}\|_2^2=V_{\cK}$ and $\|W_{\mathrm t}\|_2^2=V_{\mathrm t}$.
Let $\beta=\sum_nb_n\cos\vartheta_n$ and
$\varkappa=\varkappa_0(1-\beta)$; then
$0\le\varkappa\le\varkappa_0(1+\sum_nb_n)$. Since $\varkappa\ge0$,
$1-\varkappa F_M\ge1$ wherever $F_M\le0$, so
\[
\nu_M((-\infty,0])=\hA(F_M\le0)\le\int_{\cA}(1-\varkappa F_M)^2\,d\hA .
\]
Write $1-\varkappa F_M=Q-\varkappa W_{\cK}-\varkappa W_{\mathrm t}$
with $Q:=1-\varkappa(1+H)$, and expand:
\[
\int(1-\varkappa F_M)^2=\int Q^2-2\int Q\varkappa W_{\cK}
-2\int Q\varkappa W_{\mathrm t}+\int\varkappa^2W_{\cK}^2
+2\int\varkappa^2W_{\cK}W_{\mathrm t}+\int\varkappa^2W_{\mathrm t}^2 .
\]
The integral over $\cA$ of a monomial in the characters
$e_\gamma^{\pm1}$ is $1$ if its total frequency is $0$ and $0$
otherwise (Lemma~\ref{lem:KW}(b)). The functions $Q$, $Q\varkappa$,
$\varkappa^2$ are trigonometric polynomials in
$e_{\gamma_1},\dots,e_{\gamma_N}$ of degrees $2$, $3$, $2$.
\begin{enumerate}
\item[1.] $\int Q^2$: degree $4$ in the head; by (NR1) only the
formally constant monomials contribute, so the integral equals its
value for independent uniform phases, denoted $\EE$. With
$\mathcal G:=H-\beta-\beta H$ we have
$Q=(1-\varkappa_0)-\varkappa_0\mathcal G$,
$\EE\,\mathcal G=-\tfrac12\sum a_nb_n$,
$\EE\,\mathcal G^2=\tfrac12\sum(a_n-b_n)^2+\EE[\beta^2H^2]$ (the
degree-$3$ term has mean $0$), and
\[
\EE[\beta^2H^2]=\tfrac14\sum a_n^2\sum b_n^2
+\tfrac12\Bigl(\sum a_nb_n\Bigr)^2-\tfrac38\sum a_n^2b_n^2
\]
(from $\EE\cos^2=\tfrac12$, $\EE\cos^4=\tfrac38$). This gives the
first line of $J$ and the first four terms of the bracket.
\item[2.] $\int Q\varkappa W_{\cK}=0$ and
$\int Q\varkappa W_{\mathrm t}=0$: a head frequency
$\mathbf n\cdot\gamma$, $|\mathbf n|_1\le3$, would have to equal
$\pm\gamma'$ for an ordinate $\gamma'$ outside the head. For
$\gamma'\le T$ this is excluded by (NR2); for $\gamma'>T$ it is
impossible because $|\mathbf n\cdot\gamma|\le3\gamma_N\le T$.
\item[3.] $\int\varkappa^2W_{\cK}^2=(\int\varkappa^2)\,V_{\cK}
=\varkappa_0^2(1+\tfrac12\sum b_n^2)\,V_{\cK}$: a nonconstant head
monomial of $\varkappa^2$ has frequency $\mathbf n\cdot\gamma\ne0$,
$1\le|\mathbf n|_1\le2$ (NR1), and would need
$\mathbf n\cdot\gamma=\pm\gamma'\pm\gamma''$ with
$\gamma',\gamma''\in\cK$; this is excluded by (NR3). The constant
term of $\varkappa^2$ is $\varkappa_0^2(1+\tfrac12\sum b_n^2)$ by
(NR1).
\item[4.] $\int\varkappa^2W_{\cK}W_{\mathrm t}=0$: one needs
$\mathbf n\cdot\gamma\pm\gamma'\pm\gamma''=0$ with
$|\mathbf n|_1\le2$, $\gamma'\in\cK$, and $\gamma''>T'$ an ordinate
outside the head. If $\mathbf n=0$ formally then $\gamma'=\gamma''$,
impossible. If $\gamma''\le T$ this is excluded by (NR3). If
$\gamma''>T$: $\gamma''\le|\mathbf n\cdot\gamma|+\gamma'
\le2\gamma_N+T'=T$, a contradiction. (This is the reason for the
buffer $T'=T-2\gamma_N$.)
\item[5.] $\int\varkappa^2W_{\mathrm t}^2
\le\|\varkappa\|_\infty^2\,V_{\mathrm t}
\le\varkappa_0^2(1+\sum b_n)^2\,V_{\mathrm t}$.
\end{enumerate}
Summing gives $J$. The density statement follows as in
Theorem~\ref{thm:cantelli}.
\end{proof}

The function $J$ is increasing in $V_{\mathrm t}$, so an upper bound
for $V_M$ (Proposition~\ref{prop:Vplus} or~\ref{prop:tierB}) may be
used in place of the unknown exact value. Only items 1--4 use
arithmetic information, and each is a finite check against the
certified list: the unknown zeros beyond $T$ enter only through the
single number $V_{\mathrm t}$, in a term where no orthogonality is
used. They can resonate with nothing that enters items 1--4, by the
size argument in item~4 and by $3\gamma_N\le T$ in item~2; they may
resonate with the head at order $\ge4$, with $\cK$, or among
themselves, with any multiplicity, and the proof is indifferent to
this.

\medskip\noindent\textbf{Computation.} The script
\texttt{grhonly/nonresonance.py} searches for a good
$(\varkappa_0,b)$ in floating point; any feasible point gives a valid
bound, so the search carries no rigor. The certificate is produced by
\texttt{grhonly/nonresonance\_certified.py}:
\begin{itemize}
\item (NR1)--(NR3) are verified in \emph{exact integer arithmetic} on
the stored ordinates multiplied by $10^{24}$. Every linear form
involved has coefficient mass $\le4$, so the enclosure radius
$10^{-20}$ moves it by at most $4\cdot10^{-20}$; the smallest exact
absolute value over all head sizes $N\le20$ and both characters is
$7.9412\cdot10^{-7}$ ($\chi_3$, (NR3)).
\item $J$ is evaluated in \emph{Arb} ($200$ bits, ordinates as balls
of radius $10^{-20}$) at an exactly rational feasible point (each
$b_n$ rounded down to $6$ decimals, $\sum b_n\le1$ checked on
rationals), with
$V_{\mathrm t}=V^++10^{-10}-V_{\mathrm h}-V_{\cK}$, where $V^+$ is
the certified bound of Proposition~\ref{prop:Vplus} (tier~A) or
Proposition~\ref{prop:tierB} (tier~B) rounded up; the set $\cK$ is
decided on balls.
\end{itemize}

\begin{theorem}[Chebyshev's race, improved; GRH only]\label{thm:NRq}
Under GRH$(\chi_4)$, the lower logarithmic density of
$\{x:\pi(x;4,3)>\pi(x;4,1)\}$ is
\begin{itemize}
\item[] $>0.884276434$ \quad\certA{} ($N=8$,
$J=0.1157235655252294\ldots$, ball radius $2.4\cdot10^{-19}$);
\item[] $>0.922718033$ \quad\certB{} ($N=16$).
\end{itemize}
Under GRH$(\chi_3)$, that of $\{x:\pi(x;3,2)>\pi(x;3,1)\}$ is
$>0.911240824$ \certA{} ($N=10$, $J=0.0887591750953148\ldots$, ball
radius $6.3\cdot10^{-21}$) and $>0.944555790$ \certB{} ($N=20$).
\end{theorem}

\begin{proof}
Theorem~\ref{thm:NR} with $M=\tfrac12$, the verification of
(NR1)--(NR3) and the enclosures of $J$ just described; NCZ holds by
\eqref{eq:Lpositive}; $3\gamma_N\le T$ is checked by the script. The
printed values are outward-rounded: the comparison $1-J>d$ is
certified by Arb.
\end{proof}

(An independent Arb evaluation by the referee of the working notes
gave for $1-J$ the values $0.884276434\ldots$, $0.922718034\ldots$,
$0.911240824\ldots$ and $0.944555174\ldots$; the tier-A values
coincide,
the tier-B values differ in the ninth and seventh decimals because a
different rational point and a slightly different rounding of $V^+$
were used. Every one of these numbers is a valid lower bound.)

Table~\ref{tab:bias} collects lower bounds for $\dlow\{D_m>0\}$ at
other weights, in the form tier~A / tier~B; all entries are rounded
down, and the non-resonance rows are certified by the same script.

\begin{table}[ht]
\centering
\small
\begin{tabular}{lccccc}
\toprule
$m$ & $-0.4$ & $-0.25$ & $0$ & $0.25$ & $0.5$\\
\midrule
$\chi_4$, Cantelli & $0.9932/0.9937$ &
$0.9592/0.9624$ & $0.8549/0.8653$ & $0.7244/0.7414$ &
$0.5976/0.6184$\\
$\chi_4$, non-res. & $0.9976/0.9986$ &
$0.9784/0.9883$ & $0.8842/0.9227$ & $0.7378/0.7850$ &
$0.5983/0.6298$\\
$\chi_3$, non-res. & & $0.9830/0.9908$ &
$0.9112/0.9445$ & & $0.6695/0.7121$\\
\bottomrule
\end{tabular}
\caption{Lower bounds for the lower logarithmic density of
$\{D_m>0\}$ under GRH$(\chi_q)$ alone; ``Cantelli'' is
Theorem~\ref{thm:cantelli-q} with Propositions~\ref{prop:Vplus},
\ref{prop:tierB}, ``non-res.'' is the non-resonance bound of
Theorem~\ref{thm:NR}. First number: \certA{}
(unconditional given GRH); second number: \certB{} (assuming in
addition that all zeros up to height $10^8/q$ are simple). Under GRH
and LI the density at $m=0$ exists and equals $0.9959$ $(q=4)$,
$0.9991$ $(q=3)$ \cite{RS94}.}
\label{tab:bias}
\end{table}

\begin{remark}[limits of the method, and what was tried and failed]
\label{rem:NRlimits}
(1) \emph{Degree and positivity of the multiplier.}
Theorem~\ref{thm:NR} uses a multiplier $\varkappa$ of degree one in
the head phases. \emph{Heuristic}: with all information (every zero
in the head, full LI) a quadratic test function
$(1-\varkappa F)^2$ with $\varkappa=\varkappa(\text{head})$ can reach
$\EE[v/(v+(1+H)^2)]$ ($v$ the variance not in the head), which tends
to the LI value $\PP(1+H_\infty\le0)$ as the head grows; so in
principle the method converges to the LI density as the certified
order and height grow, provided LI is true (otherwise some
non-resonance check eventually fails). In practice the gain is
limited by (i) the sup-norm penalty
$\|\varkappa\|_\infty^2V_{\mathrm t}$, (ii) the constraint
$\sum b_n\le1$ that keeps $\varkappa\ge0$ (it binds from
$N\approx16$), and (iii) the variance $V_{\cK}$ of the zeros that are
certified but not in the head. A numerical exploration with
degree-two multipliers
$\varkappa=\varkappa_0(1-\beta+s(\beta^2-\EE\beta^2))$
(\texttt{grhonly/explore\_degree2.py}, \numtag, non-resonance of
order $6$ not verified) gave no measurable gain over degree one for
$N\le5$ (optimal $s=0$ in tier~A; gain $<0.001$ in tier~B). A
positive multiplier of Riesz-product type
$\prod(1-b_n\cos\vartheta_n)$ would remove (ii) at the price of
non-resonance checks of order about $2N$; this was not attempted.

(2) \emph{The tail.} In tier~A the term
$(1+\sum b_n)^2V_{\mathrm t}$ with $V_{\mathrm t}\approx0.0177$ (of
which $0.0141$ is the multiplicity excess) is what stops the gain at
$N=8$. Extending the certified list is the cheapest improvement
(Remark~\ref{rem:closegap}).

(3) \emph{Fourth moments.} A Paley--Zygmund argument is not
available: $\int(F_M-1)^4$ depends on all relations
$\gamma_1+\gamma_2=\gamma_3+\gamma_4$. The Hausdorff--Young bound of
Remark~\ref{rem:known} is rigorous but gives only
$\|F_{1/2}-1\|_4\lesssim1.9$ for $\chi_4$ \numtag{} against the LI
value $\approx0.51$; it is useless at $m=0$.

(4) \emph{Union bound ``head $+$ tail''.}
$\PP(F\le0)\le\PP(1+H\le t)+\PP(W\le-t)$ with Cantelli on the tail
cannot beat Cantelli at $m=0$ for $T=640$, whatever is known about
the head (the two terms sum to more than $0.15$). The quadratic test
function of Theorem~\ref{thm:NR} is what makes head information
usable.

(5) \emph{Transfer through other weights.} Using
Corollary~\ref{cor:thetaA} to bound $\{E_M\le0\}$ by an event at
$M'>M$ and Cantelli there is strictly worse than Cantelli at $M$ for
all small $M$ (computed: the derivative in $M'$ at $M'=M$ has the
wrong sign whenever
$\sigma_M^2/2>M^2\sum_\gamma(M^2+\gamma^2)^{-2}/
\sum_\gamma(M^2+\gamma^2)^{-1}$); and no convex test function can
gain (Remark~\ref{rem:nogain}).

(6) \emph{Smoothed versions.} Mean squares of smoothed races are
smaller, but the convex order $K_Mh\le_{\mathrm{cx}}h$ goes the wrong
way for bounding $\PP(h\le0)$; smoothing helps only through
Theorem~\ref{thm:peel} and Corollary~\ref{cor:thetaA}.
\end{remark}

\section{Explicit theorems for every \texorpdfstring{$x$}{x}:
\texorpdfstring{$q=4$ and $q=3$}{q=4 and q=3}}\label{sec:explicit}

In this section $q\in\{4,3\}$, $\chi=\chi_q$ is the nonprincipal
(real, primitive, odd) character mod $q$, and the only hypothesis is
GRH$(\chi_q)$: no hypothesis on $\zeta$, no linear independence, no
simplicity of zeros. The exponents $m,m'$ are arbitrary real numbers.

\subsection{Notation and statements}\label{subsec:explicitstatements}

Let $\ell$ be the prime dividing $q$ ($2$ resp.\ $3$) and $p_0$ the
least prime not dividing $q$ ($3$ resp.\ $2$). Thus
$\chi^*(p)=+1$ for $p\equiv3$ mod $4$, $-1$ for $p\equiv1$ mod $4$,
$0$ for $p=2$ if $q=4$; and $\chi^*(p)=+1$ for $p\equiv2$ mod $3$,
$-1$ for $p\equiv1$ mod $3$, $0$ for $p=3$ if $q=3$. Put
\begin{gather*}
\Theta(t):=R_0(t)=\sum_{p\le t}\chi^*(p)\log p=-\theta(t,\chi),\\
\Theta_1(t):=\int_0^t\Theta(u)\,du=\sum_{p\le t}\chi^*(p)\log p\cdot(t-p),
\qquad
G(t):=\frac{\Theta_1(t)}{t^{3/2}},\\
\cM(x):=\sum_{p\le x}\chi^*(p)\Bigl(1-\frac px\Bigr)
=D_0(x)-\frac{D_1(x)}x=\frac1x\int_0^xD_0(t)\,dt,
\qquad
\cN(x):=\frac{\cM(x)\log x}{\sqrt x},\\
\theta(y):=\sum_{p\le y}\log p,\qquad
\theta(y,\chi):=\sum_{p\le y}\chi(p)\log p,\\
\vth(y):=\sum_{p\le y,\ p\nmid q}\log p=\theta(y)-(\log\ell)\,1_{y\ge\ell}.
\end{gather*}
The function $\cM$ is continuous (the term of $p$ vanishes at $x=p$)
and equals $A-B/x$ between consecutive primes; $\Theta_1$ is
continuous and piecewise linear, and both vanish on $[0,p_0]$.

\medskip\noindent\textbf{Constants.} All digits shown are correct;
their provenance is given in Lemmas~\ref{lem:zerosq}, \ref{lem:theta}
and Section~\ref{sec:comp}.
\begin{center}
\begin{tabular}{lll}
\toprule
 & $\chi_4$ & $\chi_3$\\
\midrule
$\sigma_0^2(\chi)$ & $0.155567979923586\ldots$ &
$0.113229969857472\ldots$\\
$c_0=\tfrac23-\sigma_0^2$ & $0.511098686743\ldots$ &
$0.553436696809\ldots$\\
$c_1$ & $2.071193$ & $2.081942$\\
$Y_1$ & $10^{10}$ & $10^{10}$\\
$\eps_1=1/\log^3Y_1$ & $8.191\cdot10^{-5}$ & $8.191\cdot10^{-5}$\\
$c_\theta=1+\eps_1$ & $1.0000820$ & $1.0000820$\\
\bottomrule
\end{tabular}
\end{center}
Here $\sigma_0^2(\chi):=\log(q/\pi)-\gamma_{\mathrm E}
+2\frac{L'}{L}(1,\chi)$ is a closed form ($\gamma_{\mathrm E}$ is
Euler's constant); under GRH$(\chi)$ it equals $\sum_\rho|\rho|^{-2}
=\sigma_{1/2}^2$ (Lemma~\ref{lem:zerosq}(b)). For $\chi_4$ it is
$\gamma_{\mathrm E}+6\log2+5\log\pi-8\log\Gamma(\tfrac14)$.

Define, for $t\ge p_0$,
\begin{gather}
E_{\mathrm{odd}}(t):=\tfrac34\,t^{-1/6}+\tfrac56\,t^{-3/10}
+t^{-5/14}\cdot\max\Bigl(0,\ \frac12\Bigl(\frac{\log t}{\log p_0}-5\Bigr)\Bigr),
\notag\\
\cL_\chi(t):=c_0-\tfrac45\,c_1\,t^{-1/4}-c_\theta\,E_{\mathrm{odd}}(t)
-1_{t>Y_1^2}\cdot\Bigl[\tfrac23\,\eps_1+(\log\ell)\,t^{-1/2}\Bigr].
\label{eq:defL}
\end{gather}
The function $\cL_\chi$ is nondecreasing on $[4000,Y_1^2]$ and on
$(Y_1^2,\infty)$; $\cL_{\chi_4}(10^4)=0.0681$,
$\cL_{\chi_3}(10^4)=0.0182$, $\cL_\chi(Y_1^2+)=0.5107$ resp.\
$0.5530$, and $\cL_\chi(t)\to c_0-\tfrac23\eps_1=0.51104$ resp.\
$0.55338$ as $t\to\infty$. We write
$\cL_*(t_1,x):=\inf_{t_1\le t\le x}\cL_\chi(t)$; for $t_1\ge10^4$
this equals $\cL_\chi(t_1)$ if $x\le10^{20}$ and
$\min(\cL_\chi(t_1),\cL_\chi(10^{20}+))$ otherwise.

\begin{theorem}[integrated $\theta$-race]\label{thm:Theta1}
Let $q\in\{4,3\}$ and assume GRH$(\chi_q)$. Then
\[
\Theta_1(t)\ge\cL_{\chi_q}(t)\cdot t^{3/2}\quad\text{for every }t\ge p_0,
\qquad\text{and}\qquad
\Theta_1(t)>0\quad\text{for every }t>p_0 .
\]
\end{theorem}

\begin{theorem}[Ces\`aro mean of Chebyshev's race]\label{thm:cesaro}
Let $q\in\{4,3\}$ and assume GRH$(\chi_q)$. Then for every $x>p_0$
\[
\cM(x)=\frac1x\int_0^xD_0(t)\,dt\;\ge\;\frac{\Theta_1(x)}{x\log x}\;>\;0,
\]
and for every $x\ge10^4$ and every $t_1\in[10^4,x]$
\[
\cM(x)\;\ge\;\frac{\sqrt x}{\log x}\Bigl[\cL_\chi(x)
+\frac{\cL_*(t_1,x)}{\log x}\Bigl(\frac83-2\Bigl(\frac{t_1}x\Bigr)^{1/2}
-\frac23\Bigl(\frac{t_1}x\Bigr)^{3/2}\Bigr)\Bigr].
\]
In particular $\cN(x)\ge G(x)\ge\cL_\chi(x)$ for every $x>p_0$.
\end{theorem}

For $q=4$ this reads
$\int_0^x(\pi(t;4,3)-\pi(t;4,1))\,dt>0$ for every $x>3$, and for
$q=3$, $\int_0^x(\pi(t;3,2)-\pi(t;3,1))\,dt>0$ for every $x>2$.

\begin{corollary}\label{cor:half}
Assume GRH$(\chi_4)$. Then $\cM(x)>\sqrt x/(2\log x)$ for every
$x\ge23$. Assume GRH$(\chi_3)$. Then $\cM(x)>\sqrt x/(2\log x)$ for
every $x\ge6$.
\end{corollary}

The ranges are essentially best possible at the lower end: for $q=4$,
$\cN(19)=0.4977$ and on $[19,23)$, where $\cM(x)=1-5/x$, the value
$\tfrac12$ is crossed at $x\approx19.46$; for $q=3$,
$\cN(5)=0.4319$ and on $[5,7)$ the value $\tfrac12$ is crossed at
$x\approx5.343$.

\begin{remark}[the constant $\tfrac12$]\label{rem:halfLI}
Under GRH$(\chi_4)$ \emph{and LI in the usual sense} (the positive
ordinates are linearly independent over $\mathbb Q$; in particular
the zeros are simple),
$\liminf\cN=\liminf G=\tfrac23-\sum_{\gamma>0}2/|\rho(\rho+1)|
\approx0.5129$, so for all large $x$ the constant $\tfrac12$ of
Corollary~\ref{cor:half} cannot be replaced by anything larger than
that. Here $\liminf\cN=\liminf G$ because, by \eqref{eq:keyid}
below, $\cN-G=(\log x/x^{3/2})\int\Theta_1h_x''$ is $\ge0$ and
$O(1/\log x)$. The limiting constant of Theorem~\ref{thm:Theta1}
itself is $c_0-\tfrac23\eps_1=0.51104$. This remark is the only place
in this section where LI is mentioned, and it is not used.
\end{remark}

\begin{corollary}[inheritance of leads, $m=0$ against $m'=1$]
\label{cor:inherit}
Let $q\in\{4,3\}$ and assume GRH$(\chi_q)$. For every $x>p_0$,
\[
D_1(x)\ge0\ \Longrightarrow\ D_0(x)>0,
\qquad
D_0(x)\le0\ \Longrightarrow\ D_1(x)<0,
\]
and more precisely $D_0(x)>D_1(x)/x+\sqrt x/(2\log x)$ for $x\ge23$
$(q=4)$, resp.\ $x\ge6$ $(q=3)$. So, for the pair $(m,m')=(0,1)$,
the two implications of Theorem~\ref{thm:MCA}(ii) hold for every
$x>p_0$, not just for large $x$.
\end{corollary}

\begin{theorem}[the cone $m\le0$, $m<m'\le1$]\label{thm:cone}
Let $q\in\{4,3\}$ and assume GRH$(\chi_q)$. For all real $m\le0$ and
$m<m'\le1$ ($m'$ may be negative), with $\kappa=m'-m$, and every
$x>p_0$:
\[
D_m(x)-x^{-\kappa}D_{m'}(x)\;\ge\;
\frac{\kappa\,x^{m-1}\,\Theta_1(x)}{\log x}\;\ge\;
\frac{\kappa\,x^{m+1/2}\,\cL_\chi(x)}{\log x},
\]
the first lower bound being $>0$. Hence
$D_{m'}(x)\ge0\Longrightarrow D_m(x)>0$ for every $x>p_0$.
\end{theorem}

The equivalence with GRH is Theorem~\ref{thm:equiv} in
Section~\ref{sec:landau}.

\medskip\noindent\textbf{Numerical truth} (sieve,
Section~\ref{sec:comp}; not used in the proofs except where stated).
For $q=4$: $\cM(x)>0$ on $3<x\le10^{10}$ (exact integer test) with
$\min_{x\ge5}\cM=2/7$ at $x=7$; $\inf\cN$ over $[10^4,10^{10}]$ is
$0.682694$, at the prime $x=1\,070\,765\,639$; $\inf\cN$ over
$[11,10^{10}]$ is $\cN(11)=(6/11)\log11/\sqrt{11}=0.394360$;
$\inf G$ over $[2^{13},10^{10}]$ is $0.5196$. For $q=3$:
$\min_{x\ge5}\cM=0.6$ at $x=5$; $\inf\cN$ over $[10^4,10^{10}]$ is
$0.720493$, at $x=323\,405\,729$; $\inf G$ over $[2^{13},10^{10}]$ is
$0.5411$. Checkpoints: $D_0(10^{10})=5960$,
$D_1(10^{10})=27\,201\,020\,910\,942$ for $q=4$, and
$D_0(10^{10})=6264$, $D_1(10^{10})=27\,948\,685\,258\,617$ for $q=3$.

\subsection{The summation identity}

\begin{lemma}[second-order partial summation; no hypothesis]
\label{lem:ingham}
Let $x>p_0$ and $F\in C^2[p_0,x]$. Then
\[
\sum_{p\le x}\chi^*(p)\log p\cdot F(p)
=F(x)\,\Theta(x)-F'(x)\,\Theta_1(x)+\int_{p_0}^x\Theta_1(t)\,F''(t)\,dt .
\]
\end{lemma}

\begin{proof}
The left-hand side is
$\int_{p_0^-}^xF\,d\Theta=F(x)\Theta(x)-\int_{p_0}^x\Theta F'\,dt$.
Since $\Theta_1'=\Theta$ a.e., $\Theta_1$ is absolutely continuous
and $\Theta_1(p_0)=0$, a second integration by parts gives
$\int_{p_0}^x\Theta F'=\Theta_1(x)F'(x)-\int_{p_0}^x\Theta_1F''$.
\end{proof}

Take $F(t)=h_x(t):=(x-t)/\log t$. Then $h_x(x)=0$,
$h_x'(x)=-1/\log x$, and
\begin{equation}\label{eq:hsecond}
h_x''(t)=\frac{2}{t\log^2t}+(x-t)\Bigl[\frac1{t^2\log^2t}
+\frac2{t^2\log^3t}\Bigr]
=\frac{x+t}{t^2\log^2t}+\frac{2(x-t)}{t^2\log^3t}\;>\;0
\qquad(1<t\le x).
\end{equation}
Since $\sum_{p\le x}\chi^*(p)\log p\cdot h_x(p)
=\sum_{p\le x}\chi^*(p)(x-p)=x\,\cM(x)$, we get the \emph{exact}
identity
\begin{equation}\label{eq:keyid}
x\,\cM(x)=\frac{\Theta_1(x)}{\log x}
+\int_{p_0}^x\Theta_1(t)\,h_x''(t)\,dt\qquad(x>p_0).
\end{equation}
(Checked numerically to $22$ digits at $x=200.5$ and $x=1000.25$ by
the script \texttt{verify\_small.py} of Section~\ref{sec:comp}.)

\medskip\noindent\textbf{Consequence.} If $\Theta_1\ge0$ on
$[p_0,x]$, then $\cM(x)\ge\Theta_1(x)/(x\log x)$. So the positivity
of $\cM$ --- a sum with weights $1-p/x$ and no $\log p$ --- follows
from the positivity of the \emph{one} function $\Theta_1$, with no
loss; in particular the factor $1/\log x$ that an estimate of the
kernel term would cost is not lost.

\medskip\noindent\textbf{Attribution.} Lemma~\ref{lem:ingham} is not
new. It is the second-order partial summation of Ingham
\cite[pp.~103--104]{Ing32} (as cited by Johnston), and
\cite[Lemma~2.7]{Joh23} is the same identity, with the same positive
kernel $1/(t^2\log^2t)+2/(t^2\log^3t)$, acting on the integrated
error $\int_2^x(\psi(t)-t)\,dt$. In that setting the integrated error
is not one-signed, so the kernel term has to be estimated (this is
done in \cite{Joh23} for $x>10^{19}$, a finite verification being
used below). The only new observation here is that the function
integrated against the kernel, $\Theta_1$, \emph{is} one-signed, so
the kernel term can simply be dropped; this is why the analytic
argument can start at $10^4$.

\subsection{Zeros: an explicit formula for
\texorpdfstring{$\psi_1(x,\chi)$}{psi\_1(x,chi)}}

\begin{lemma}\label{lem:zerosq}
Let $\chi$ be a primitive odd character mod $q$ (so $L(0,\chi)\ne0$
and the trivial zeros are simple, at $s=-1,-3,-5,\dots$). Put
$\mathfrak b:=\frac{L'}{L}(0,\chi)$ and let $\mathfrak b_0$ be the
constant term of the Laurent expansion
$\frac{L'}{L}(s,\chi)=\frac1{s+1}+\mathfrak b_0+O(s+1)$.
\begin{enumerate}
\item[(a)] \emph{(Explicit formula for $\psi_1$; no hypothesis.)} For
$x\ge1$,
\[
\psi_1(x,\chi)=-\sum_\rho\frac{x^{\rho+1}}{\rho(\rho+1)}
-\mathfrak b\,x+\log x+1+\mathfrak b_0
-\sum_{n\ge1}\frac{x^{-2n}}{2n(2n+1)},
\]
the sum over all nontrivial zeros (with multiplicity) being
absolutely convergent.
\item[(b)] If $\chi$ is real and GRH$(\chi)$ holds, then
\[
\sum_\rho\frac1{|\rho(\rho+1)|}\;<\;\sum_\rho\frac1{|\rho|^2}
=\log\frac q\pi-\gamma_{\mathrm E}+2\,\frac{L'}{L}(1,\chi)=\sigma_0^2(\chi).
\]
\end{enumerate}
\end{lemma}

This sharpens, for odd $\chi$, the form with remainder $O(t\log2t)$
used in \cite[the lemma ``explicit formula for $\psi_1$'']{MC}; the case of $\zeta$ is classical
(\cite[Theorem~27]{Ing32}, as cited in \cite{Joh23}).

\begin{proof}[Proof of \textup{(a)}]
Let $x>1$ and $c=2$. Then
\[
\psi_1(x,\chi)=\frac1{2\pi i}\int_{(c)}\Bigl(-\frac{L'}{L}\Bigr)(s,\chi)\,
\frac{x^{s+1}}{s(s+1)}\,ds
\]
(Mellin inversion of $(x-n)_+$, term by term; the integral converges
absolutely). Integrate over the rectangle with vertices $c\pm iT$,
$-U\pm iT$, where $U$ is an even integer (so the left side passes
midway between the trivial zeros $-U\pm1$) and $T=T_k\to\infty$ is
chosen with $|T_k-\gamma|\gg1/\log T_k$ for every ordinate $\gamma$
(possible because there are $\ll\log(qT)$ ordinates in $[T,T+1]$).
Two standard bounds are used (\cite[Ch.~19]{Dav00}, where the harder
kernel $x^s/s$ is treated; cf.\ Ch.~17 there for $\zeta$):
\begin{itemize}
\item[(i)] $\frac{L'}{L}(\sigma\pm iT_k,\chi)\ll\log^2(qT_k)$
uniformly for $-1\le\sigma\le2$;
\item[(ii)] $\frac{L'}{L}(s,\chi)\ll\log(q(|s|+2))$ for $\re s\le-1$,
$|s+(2n+1)|\ge\tfrac12$ for all $n\ge0$
\end{itemize}
((ii) follows from the functional equation). \emph{Horizontal
sides.} On $\im s=\pm T_k$ we have $|s(s+1)|\ge T_k^2$, so by (i) and
(ii) the two integrals are
\[
\ll T_k^{-2}\int_{-\infty}^{2}\bigl[\log^2(qT_k)+\log(|\sigma|+2)\bigr]
x^{\sigma+1}\,d\sigma\;\ll_x\;\frac{\log^2(qT_k)}{T_k^2},
\]
uniformly in $U$ (here $x>1$ makes
$\int_{-\infty}^2\log(|\sigma|+2)\,x^{\sigma+1}\,d\sigma$ finite);
this tends to $0$ as $k\to\infty$. \emph{Left side.} On $\re s=-U$,
by (ii),
\[
\Bigl|\int\Bigr|\ll x^{1-U}\int_{-\infty}^{\infty}
\frac{\log(q(U+|t|+2))}{U^2+t^2}\,dt\ll\frac{x^{1-U}\log(qU)}{U}
\to0\qquad(U\to\infty).
\]
Letting first $U\to\infty$ and then $k\to\infty$, $\psi_1(x,\chi)$
equals the sum of all residues to the left of $\re s=c$; the sum over
$\rho$ needs no grouping, since
$\sum_\rho|x^{\rho+1}/(\rho(\rho+1))|\le x^2\sum_\rho|\rho|^{-2}
<\infty$. The residues are: at $s=\rho$,
$-x^{\rho+1}/(\rho(\rho+1))$; at $s=0$, $-\mathfrak b\,x$; at
$s=-(2n+1)$, $n\ge1$, $-x^{-2n}/(2n(2n+1))$; at the double pole
$s=-1$, with $-\frac{L'}{L}=-\frac1{s+1}-\mathfrak b_0+\cdots$ and
$\frac{x^{s+1}}{s}=-1-(s+1)(\log x+1)+\cdots$, the residue is
$\log x+1+\mathfrak b_0$. Both sides of the formula are continuous on
$[1,2]$ (the two series converge uniformly there), so it also holds
at $x=1$.
\end{proof}

\begin{proof}[Proof of \textup{(b)}]
$|\rho+1|>|\rho|$ since $\re\rho>0$. By \cite[Ch.~12, eqs.~(17),
(18)]{Dav00},
$\frac{L'}{L}(s,\chi)=-\tfrac12\log\frac q\pi
-\tfrac12\frac{\Gamma'}{\Gamma}\bigl(\frac{s+1}2\bigr)+B(\chi)
+\sum_\rho\bigl(\frac1{s-\rho}+\frac1\rho\bigr)$ and
$\re B(\chi)=-\sum_\rho\re(1/\rho)$. At $s=1$, taking real parts
($\chi$ is real, so the left-hand side is real):
$\sum_\rho\re\frac1{1-\rho}=\frac{L'}{L}(1,\chi)
+\tfrac12\log\frac q\pi-\tfrac12\gamma_{\mathrm E}$. Under GRH,
$\re\frac1{1-\rho}=\tfrac12/|\rho|^2$.
\end{proof}

(The equation numbers in \cite{Dav00} in the two proofs above are
quoted from memory and were not checked against the book; see
Section~\ref{sec:gaps}.)

\medskip\noindent\textbf{Values and checks.} For $\chi_4$:
$\mathfrak b=0.78318878541$, $\mathfrak b_0=-0.06026691418$; for
$\chi_3$: $\mathfrak b=0.94819882667$, $\mathfrak b_0=0.14383164456$;
these come from the functional equation
($\mathfrak b=-\log(q/\pi)+\gamma_{\mathrm E}+\log2
-\frac{L'}{L}(1,\chi)$,
$\mathfrak b_0=-\log(q/\pi)+\tfrac12\gamma_{\mathrm E}
-\tfrac12\frac{\Gamma'}{\Gamma}(\tfrac32)-\frac{L'}{L}(2,\chi)$,
$\chi$ real) and were confirmed by direct evaluation of $L'/L$ near
$0$ and $-1$. Three independent checks of the formula of (a) were
made (\texttt{explicit/constants.py}): (i) its $x$-derivative is the
classical formula $\psi(x,\chi)=-\sum_\rho x^\rho/\rho
-\mathfrak b+\sum_{n\ge1}x^{1-2n}/(2n-1)$ for odd $\chi$
\cite[Ch.~19]{Dav00}; (ii) evaluating at $x=1$, where
$\psi_1(1,\chi)=0$, the constant must equal
$\sum_\rho1/(\rho(\rho+1))+\mathfrak b+1-\log2$, and the Hadamard
product gives $\sum_\rho1/(\rho(\rho+1))=\mathfrak b_0-\mathfrak b
+\log2$, so the constant is $1+\mathfrak b_0$, the same value;
numerically $\sum_\rho1/(\rho(\rho+1))$ from the certified zero lists
is $-0.150312$ ($\chi_4$), $-0.111224$ ($\chi_3$) against
$\mathfrak b_0-\mathfrak b+\log2=-0.150309$, $-0.111220$; (iii)
direct comparison with the prime-power sum, using $511$ resp.\ $537$
zeros, at small $x$, where the truncation of the zero sum is small
enough for the test to say something (``bound'' is
$x^{3/2}\sum_{\gamma>T}2/\gamma^2$ from the Riemann--von Mangoldt
density; a wrong constant term would show up as a difference of
order $1$):
\begin{center}
\small
\begin{tabular}{rrrrr}
\toprule
$x$ & $\psi_1$ direct $-$ formula ($\chi_4$) & bound &
$\psi_1$ direct $-$ formula ($\chi_3$) & bound\\
\midrule
$1$ & $+3.48\cdot10^{-3}$ & $0.0035$ & $+3.09\cdot10^{-3}$ & $0.0031$\\
$3$ & $+1.64\cdot10^{-3}$ & $0.018$ & $-6.5\cdot10^{-7}$ & $0.016$\\
$5.5$ & $-1.2\cdot10^{-6}$ & $0.045$ & $+1.6\cdot10^{-5}$ & $0.040$\\
$10$ & $+7.5\cdot10^{-5}$ & $0.11$ & $+1.7\cdot10^{-4}$ & $0.098$\\
$20$ & $-1.6\cdot10^{-3}$ & $0.31$ & $-8.9\cdot10^{-4}$ & $0.28$\\
$50$ & $-3.2\cdot10^{-4}$ & $1.2$ & $+3.1\cdot10^{-3}$ & $1.1$\\
$200$ & $-8.3\cdot10^{-2}$ & $9.9$ & $+3.6\cdot10^{-3}$ & $8.8$\\
\bottomrule
\end{tabular}
\end{center}
(At $x=1$ the difference is the non-oscillating tail
$\sum_{\gamma>T}1/(\rho(\rho+1))$.) For (b): $511$ zeros of $\chi_4$
plus the Riemann--von Mangoldt tail give $0.1555712$ against the
closed form $0.1555680$; for $\chi_3$, $0.1132341$ against
$0.1132300$. The true value of $S:=\sum_\rho1/|\rho(\rho+1)|$ is
about $0.15380$ ($\chi_4$) and $0.11256$ ($\chi_3$); the proof below
uses the closed-form majorant $\sigma_0^2$, losing $0.0018$.

\subsection{Inputs on \texorpdfstring{$\theta$}{theta}}

\begin{lemma}\label{lem:theta}
\begin{enumerate}
\item[(a)] \emph{(Computation.)} For $0<y\le Y_1=10^{10}$:
$\ y-c_1\sqrt y\le\vth(y)<y$, with $c_1=2.071193$ $(q=4)$,
$2.081942$ $(q=3)$.
\item[(b)] \emph{(Cited.)} For $y\ge Y_1$:
$|\theta(y)-y|<y/\log^3y\le\eps_1y$, $\eps_1=1/\log^3(10^{10})$.
\item[(b$'$)] \emph{(Fallback, classical; not used in the main
line.)} $x-x/\log x<\theta(x)<x+x/(2\log x)$ for $x\ge41$.
\end{enumerate}
Consequently $\vth(y)\le c_\theta\,y$ for all $y>0$, and
$\vth(y)\ge(1-\eps_1)y-\log\ell$ for $y>Y_1$.
\end{lemma}

\noindent\emph{Basis of \textup{(a)}.} The supremum of
$(y-\vth(y))/\sqrt y$ is attained as $y\to1423^-$ (its value is
$2.0711926180$ resp.\ $2.0819411967$). For $y\le2^{22}$ both
inequalities are \emph{certified in ball arithmetic}
(\texttt{explicit/cert\_small.py}, part~C). For
$2^{22}<y\le10^{10}$ they rest on a \texttt{long double} sieve
(\texttt{explicit/cesaro\_scan.c}): there the supremum of
$(y-\vth(y))/\sqrt y$ over $y\ge1427$ is $1.98197$ resp.\ $1.99271$
(as $y\to1427^-$), so the lower inequality holds with margin
$\ge0.089\sqrt y\ge180$, and $\vth(y)-y\le-0.26\sqrt y$; the
worst-case accumulated rounding error of $\theta$ in that
computation is $0.21$. This part of (a) is therefore safe but
\emph{not certified in interval arithmetic};
Remark~\ref{rem:variantC} (variant~C) shows how to do without it. No
cited bound is used for (a). (Subtracting the $\log\ell/\sqrt{1423}$
contributed by the missing prime gives $2.05282$ for $\theta$ itself,
a constant we recall from Rosser--Schoenfeld \cite{RS62}; not checked
against the source. An independent source for the upper inequality
$\vth(y)<y$ is B\"uthe \cite[Theorem~2, (1.7)]{But18}:
$x-\vartheta(x)>0.05\sqrt x$ for $1\le x\le10^{19}$; it is not
needed.)

\noindent\emph{Source of \textup{(b)}.} Broadbent, Kadiri, Lumley, Ng
and Wilk \cite[Corollary~11.2, eq.~(3.39)]{BKLNW21}:
$|\vartheta(x)-x|<M_kx/\log^kx$ for $x\ge X_0$, with the row $k=3$,
$M_3=1$, $X_0=89\,967\,803$. On its input: their Appendix~A.1.1 says
that the computations use the partial verification of RH up to
$H=2\,445\,999\,556\,030$ announced by Gourdon and Wedeniwski, and
that Platt and Trudgian \cite{PT21} verify
$H=3\,000\,175\,332\,800$ with more rigorous techniques and thus
confirm the earlier calculations. Since
$3\,000\,175\,332\,800>2\,445\,999\,556\,030$, the input of
\cite{BKLNW21} is a consequence of the verification of \cite{PT21}.
No hypothesis on $\zeta$ is involved: (b) is unconditional.

\noindent\emph{Source of \textup{(b$'$)}.} This is the form in which
\cite[Lemma~2.2, (2.2)]{Joh23} quotes Rosser--Schoenfeld
\cite[(3.15)--(3.16)]{RS62}; the original was not consulted. With
(b$'$) in place of (b) one has $\eps_{\mathrm{low}}=1/\log Y_1
=0.0434$, $\eps_{\mathrm{up}}=0.0217$, $c_\theta=1.0217$, and
$\tfrac23\eps_1$ is replaced by $\tfrac23\eps_{\mathrm{low}}=0.0290$;
then the positivity thresholds of $\cL$ become $5229$ / $9148$,
$\cL(10^4)=0.0622$ / $0.0103$, and the limit of $\cL$ is $0.4821$ /
$0.5245$. Theorems~\ref{thm:Theta1}, \ref{thm:cesaro},
\ref{thm:cone}, \ref{thm:equiv} and the implications of
Corollary~\ref{cor:inherit} are unaffected;
Corollary~\ref{cor:half} for $\chi_4$ beyond $x=10^{20}$ (and with it
the quantitative clause of Corollary~\ref{cor:inherit} in that range)
is the only statement that needs (b) (for $\chi_3$ the fallback limit
$0.5245$ still exceeds $\tfrac12$).

\begin{remark}[variant C: certified inputs only]\label{rem:variantC}
The uncertified part of Lemma~\ref{lem:theta}(a) can be avoided
altogether. Take $Y_1^{C}=2^{22}=4\,194\,304$, where (a) is certified
in ball arithmetic, and for $y\ge2^{22}$ use another row of the same
\cite[Corollary~11.2]{BKLNW21}, $k=2$, $M_2=0.2$,
$X_0=3\,594\,641$: $|\theta(y)-y|<0.2\,y/\log^2y\le\eps_1^{C}y$ with
$\eps_1^{C}=0.2/\log^2(2^{22})=8.601\cdot10^{-4}$,
$c_\theta^{C}=1.000860$. Let $\cL^{C}$ be the function
\eqref{eq:defL} with $(Y_1,\eps_1,c_\theta)$ replaced by
$(Y_1^{C},\eps_1^{C},c_\theta^{C})$, so that the indicator becomes
$1_{t>2^{44}}$, $2^{44}=1.759\cdot10^{13}$. The proofs below go
through word for word, so \emph{every statement of
Section~\ref{subsec:explicitstatements} holds with $\cL^{C}$ in place
of $\cL$} (and $2^{44}$ in place of $10^{20}$ in the description of
$\cL_*$). Numbers ($\chi_4$ / $\chi_3$): $\cL^{C}(10^4)=0.0679$ /
$0.0179$; positivity thresholds $4904$ / $8548$;
$\cL^{C}(2^{44}+)=0.5048$ / $0.5469$; limit $0.51053$ / $0.55286$;
the analytic part of Corollary~\ref{cor:half} starts at the same
points $1.514\cdot10^8$ / $2.512\cdot10^7$. One has
$\cL^{C}\le\cL$, and $\cL-\cL^{C}$ is below $3\cdot10^{-4}$ for
$10^4\le t\le2^{44}$ and below $8\cdot10^{-4}$ for all $t\ge10^4$.
\emph{In this variant Theorems~\ref{thm:Theta1}, \ref{thm:cesaro},
\ref{thm:cone}, \ref{thm:equiv} and
Corollaries~\ref{cor:half}, \ref{cor:inherit} use only: the exact and
ball-certified finite computations of
\textup{\texttt{cert\_small.py}}, \cite[Corollary~11.2]{BKLNW21}, the
double-precision evaluations of the closed-form function $\cL^{C}$
described in the proof of Corollary~\ref{cor:half}, and the
analysis.} The main text keeps $Y_1=10^{10}$ only because the
constants are slightly better.
\end{remark}

\subsection{Proof of Theorem~\ref{thm:Theta1}}

\emph{Step 1: decomposition.} Since
$\psi(u,\chi)=\sum_{k\ge1}\theta(u^{1/k},\chi^k)$ and $\chi^k$ is
$\chi$ for odd $k$ and the principal character for even $k$,
\[
\Theta(u)=-\psi(u,\chi)+\vth(u^{1/2})
+\sum_{k\ge3\ \mathrm{odd}}\theta(u^{1/k},\chi)
+\sum_{k\ge4\ \mathrm{even}}\vth(u^{1/k}).
\]
Integrate over $[0,t]$, $t\ge p_0$. The last sum is $\ge0$ and is
dropped.

\emph{Step 2: zeros.} By Lemma~\ref{lem:zerosq}, under GRH
($|t^{\rho+1}|=t^{3/2}$):
\[
-\psi_1(t,\chi)\;\ge\;-\sigma_0^2\,t^{3/2}
+\bigl(\mathfrak b\,t-\log t-1-\mathfrak b_0\bigr),
\]
the sum over the trivial zeros entering $-\psi_1$ with a plus sign.
The bracket is increasing for $t>1/\mathfrak b$ and positive at
$t=p_0$ ($0.311$ for $\chi_4$ at $t=3$; $0.059$ for $\chi_3$ at
$t=2$), hence $-\psi_1(t,\chi)>-\sigma_0^2\,t^{3/2}$ for $t\ge p_0$.

\emph{Step 3: prime squares (the bias).} Let
$B_2(t):=\int_0^t\vth(u^{1/2})\,du=\int_0^{\sqrt t}\vth(y)\,2y\,dy$.
If $\sqrt t\le Y_1$, Lemma~\ref{lem:theta}(a) gives
$B_2(t)\ge\int_0^{\sqrt t}(2y^2-2c_1y^{3/2})\,dy
=\tfrac23t^{3/2}-\tfrac45c_1t^{5/4}$. If $\sqrt t>Y_1$, split at
$Y_1$ and use Lemma~\ref{lem:theta}(b) on $[Y_1,\sqrt t]$:
\begin{align*}
B_2(t)&\ge\tfrac23Y_1^3-\tfrac45c_1Y_1^{5/2}
+\tfrac23(1-\eps_1)\bigl(t^{3/2}-Y_1^3\bigr)-(\log\ell)\bigl(t-Y_1^2\bigr)\\
&\ge\tfrac23(1-\eps_1)\,t^{3/2}-\tfrac45c_1\,t^{5/4}-(\log\ell)\,t .
\end{align*}

\emph{Step 4: odd prime powers $k\ge3$.}
$|\theta(y,\chi)|\le\vth(y)\le c_\theta\,y$, and $\theta(y,\chi)=0$
for $y<p_0$. So
$|\int_0^t\theta(u^{1/k},\chi)\,du|\le c_\theta\frac{k}{k+1}t^{1+1/k}$,
and the term vanishes unless $k\le\log t/\log p_0$. For $k=3$ and
$k=5$ keep the constants $\tfrac34$ and $\tfrac56$; for odd $k\ge7$
use $\frac{k}{k+1}t^{1+1/k}\le t^{8/7}$; there are at most
$\max(0,(\log t/\log p_0-5)/2)$ such $k$. Hence the total is at most
$c_\theta\,E_{\mathrm{odd}}(t)\,t^{3/2}$.

\emph{Step 5.} Adding up:
$\Theta_1(t)\ge\cL_\chi(t)\,t^{3/2}$ for all $t\ge p_0$ (with strict
inequality, by Step~2).

\emph{Step 6: monotonicity and positivity.} The functions $t^{-1/4}$,
$t^{-1/6}$, $t^{-3/10}$ decrease, and
$t^{-5/14}(\log t/\log p_0-5)$ decreases once
$\log t>5\log p_0+14/5$, i.e.\ for $t>3996$ $(q=4)$, $t>527$
$(q=3)$. So $\cL_\chi$ is nondecreasing on $[4000,Y_1^2]$ and on
$(Y_1^2,\infty)$, and
$\inf_{t\ge10^4}\cL_\chi(t)=\min(\cL_\chi(10^4),\cL_\chi(Y_1^2+))
=0.0681$ $(\chi_4)$, $0.0182$ $(\chi_3)$. Hence $\Theta_1(t)>0$ for
$t\ge10^4$. (The exact thresholds of the analytic bound are
$t\ge4892$ for $\chi_4$ and $t\ge8526$ for $\chi_3$.)

\emph{Step 7: small $t$.} $\Theta_1$ is piecewise linear with nodes
at the primes $p\nmid q$, $\Theta_1(p_0)=0$, and its first slope is
$\log p_0>0$; so $\Theta_1>0$ on $(p_0,10^4]$ as soon as
$\Theta_1(p)>0$ at every prime $p\nmid q$ with $p_0<p\le10^4$. This is
checked \emph{exactly} (\texttt{explicit/cert\_small.py}, part~A):
$\Theta_1(p)>0$ is the inequality between integers
\[
\prod_{r\le p,\ \chi^*(r)=+1}r^{\,p-r}\;>\;
\prod_{r\le p,\ \chi^*(r)=-1}r^{\,p-r}
\qquad(r\text{ prime}),
\]
verified with integers of up to $3.6\cdot10^7$ bits at all $1227$
primes $p\nmid q$ with $p_0<p\le10^4$, for both characters: no
failure. The smallest values are
$\Theta_1(7)=4\log3-2\log5=1.17557$ ($\chi_4$; $G(7)=0.0635$) and
$\Theta_1(5)=3\log2=2.07944$ ($\chi_3$; $G(5)=0.1860$). (Beyond what
the proof needs: $\Theta_1(p)>0$ is certified in $256$-bit ball
arithmetic for $p\le2^{22}$, part~B of the same script, and found in
\texttt{long double} for $p\le10^{10}$.) \qed

\subsection{Proofs of Theorem~\ref{thm:cesaro} and
Corollaries~\ref{cor:half}, \ref{cor:inherit}}

\begin{proof}[Proof of Theorem~\ref{thm:cesaro}]
By Theorem~\ref{thm:Theta1} and \eqref{eq:keyid}, for $x>p_0$:
$x\,\cM(x)\ge\Theta_1(x)/\log x>0$, i.e.\
$\cN(x)\ge G(x)\ge\cL_\chi(x)$. For the refinement let
$10^4\le t_1\le x$. On $[t_1,x]$,
$\Theta_1(t)\ge\cL_*(t_1,x)\,t^{3/2}$ with $\cL_*>0$, and
$h_x''(t)\ge(x+t)/(t^2\log^2x)$ by \eqref{eq:hsecond}. Therefore
\[
\int_{p_0}^x\Theta_1h_x''\;\ge\;\frac{\cL_*}{\log^2x}\int_{t_1}^x
t^{-1/2}(x+t)\,dt
=\frac{\cL_*}{\log^2x}\Bigl[\tfrac83x^{3/2}-2x\,t_1^{1/2}
-\tfrac23t_1^{3/2}\Bigr].\qedhere
\]
\end{proof}

(The asymptotic expansion of the mean of $\cN$ is
$\tfrac23(1+\tfrac83/\log x+\cdots)$, which is why the refinement
matters numerically: at $x=10^{10}$ it is worth $0.048$.)

\begin{proof}[Proof of Corollary~\ref{cor:half} for $\chi_4$]
\emph{The range $x\ge1.52\cdot10^8$ (analytic).} Put
$\varpi(u)=\tfrac83-2\sqrt u-\tfrac23u^{3/2}$ and
$\varrho=10^{1/50}$. For $x\in[a,\varrho a]$ with $10^6\le a$ and
$\varrho a\le10^{20}$, Theorem~\ref{thm:cesaro} with $t_1=a/100$
gives
\[
\cN(x)\;>\;\beta(a):=\cL_\chi(a)
+\cL_*(a/100,\infty)\,\frac{\varpi(1/100)}{\log(\varrho a)},
\]
because $\cL_\chi$ is nondecreasing on $[4000,10^{20}]$,
$t_1/x\le1/100$, and $\log x\le\log(\varrho a)$; the inequality is
strict by Step~5 above. The script \texttt{explicit/bounds.py}
evaluates $\beta(a)$ on the grid $a=10^6\varrho^j$:
$\beta(a)\ge\tfrac12$ for every grid point in
$[1.514\cdot10^8,\,5.495\cdot10^{11}]$, and
$\inf_{t\ge a}\cL_\chi(t)\ge\tfrac12$ from $a=5.495\cdot10^{11}$ on.
Hence $\cN(x)>\tfrac12$ for all $x\ge1.52\cdot10^8$. For
$x>10^{20}$ this step uses Lemma~\ref{lem:theta}(b); with the
fallback (b$'$) the limit of $\cL$ is $0.4821<\tfrac12$ and the
corollary is only obtained for $x\le10^{20}$. The grid evaluation is
in double precision. Its margin $\beta(a)-\tfrac12$ is $0.0346$ or
more from $a=10^{10}$ on, but it is small at the lower end of the
grid: $1.5\cdot10^{-4}$ at the first grid point
$a=1.514\cdot10^8$ and $7.9\cdot10^{-4}$ at the next one,
$a=1.585\cdot10^8$, which is the first one needed in view of the next
paragraph (these two values were recomputed for this article; see
Section~\ref{sec:gaps}, item~\ref{gap:margin}).

\emph{The range $29\le x\le1.6\cdot10^8$ (certified computation).} On
a gap $[p,p']$ between consecutive primes not dividing $q$ with
$p\ge11$, $\cM$ is monotone and $\log x/\sqrt x$ decreases, so
$\cN(x)\ge\min(\cM(p),\cM(p'))\cdot\log p'/\sqrt{p'}$. The script
\texttt{cert\_small.py}, part~E, certifies in ball arithmetic (exact
integers $D_0$, $D_1$; Arb for the division, $\log$ and
$\sqrt{\cdot}$) that this gap bound exceeds $\tfrac12$ on all
$8\,974\,448$ gaps with $29\le p<p'\le1.6\cdot10^8$; the smallest gap
bound is $0.603379$ (the gap starting at $47$), a margin of $0.103$.
(The \texttt{long double} sieve gives the same minimum over all gaps
up to $10^{10}$; those values are formed from the exact $D_0$, $D_1$
by a fixed handful of operations, with no accumulation.)

\emph{The range $23\le x\le29$ (by hand; also certified in
\texttt{cert\_small.py}, part~D).} On $[23,29]$,
$\cM(x)=2-28/x$. On $[23,25]$: $\cM\ge0.7826$,
$\log x/\sqrt x\ge\log25/5=0.6438$, product $\ge0.5038$. On
$[25,29]$: $\cM\ge0.88$, $\log x/\sqrt x\ge0.6253$, product
$\ge0.550$.
\end{proof}

\begin{proof}[Proof of Corollary~\ref{cor:half} for $\chi_3$]
The analytic argument applies from $2.52\cdot10^7$ on (grid points
from $a=2.512\cdot10^7$; margin $3.9\cdot10^{-4}$ at that first
grid point, recomputed for this article), and
$\cL_{\chi_3}(t)\ge\tfrac12$ alone from $3.47\cdot10^8$ on. Certified
computation (part~E): the gap bound exceeds $\tfrac12$ on all
$1\,624\,522$ gaps with $11\le p<p'\le2.6\cdot10^7$, the smallest
being $0.711389$ (the gap starting at $11$; the same minimum holds up
to $10^{10}$ in the \texttt{long double} sieve). On $[7,11]$,
$\cM\equiv1$ and $\cN=\log x/\sqrt x\ge0.7230$; on $[6,7]$,
$\cM=2-7/x\ge5/6$ and $\log x/\sqrt x\ge\log6/\sqrt6$ (it increases
up to $e^2$), product $\ge0.6096$. Lemma~\ref{lem:theta}(b) is not
needed here: with (b$'$) the limit of $\cL_{\chi_3}$ is
$0.5245>\tfrac12$.
\end{proof}

\begin{proof}[Proof of Corollary~\ref{cor:inherit}]
$D_0(x)=D_1(x)/x+\cM(x)$, $\cM(x)>0$ for $x>p_0$
(Theorem~\ref{thm:cesaro}), and Corollary~\ref{cor:half}.
\end{proof}

\subsection{Proof of Theorem~\ref{thm:cone}}

Fix $x>p_0$ and put $f(t)=t^m-x^{-\kappa}t^{m'}$, $F=f/\log t$ on
$[p_0,x]$. Then
$\sum_{p\le x}\chi^*(p)\log p\cdot F(p)=D_m(x)-x^{-\kappa}D_{m'}(x)$,
$F(x)=0$, and $F'(x)=f'(x)/\log x=-\kappa x^{m-1}/\log x$.
\begin{itemize}
\item $f\ge0$ on $[p_0,x]$.
\item $f'(t)=t^{m-1}\bigl[m-m'(t/x)^\kappa\bigr]\le0$: if $m'\ge0$
because $m\le0$; if $m'<0$ because $|m'|(t/x)^\kappa\le|m'|<|m|$.
\item $f''(t)=t^{m-2}\bigl[m(m-1)-m'(m'-1)(t/x)^\kappa\bigr]\ge0$:
$m(m-1)\ge0$; if $0\le m'\le1$ then $m'(m'-1)\le0$; if $m'<0$ then
$m'(m'-1)<m(m-1)$ since $u(u-1)$ is decreasing on $u<\tfrac12$.
\item $1/\log t$ is positive, decreasing and convex on $t>1$.
\end{itemize}
A product of two nonnegative nonincreasing convex functions is
convex ($(fg)''=f''g+2f'g'+fg''$), so $F''\ge0$.
Lemma~\ref{lem:ingham} and Theorem~\ref{thm:Theta1} give
\[
D_m(x)-x^{-\kappa}D_{m'}(x)=-F'(x)\,\Theta_1(x)+\int_{p_0}^x\Theta_1F''
\;\ge\;\frac{\kappa\,x^{m-1}\,\Theta_1(x)}{\log x}\;>\;0 ,
\]
and $\Theta_1(x)\ge\cL_\chi(x)x^{3/2}$. \qed

\smallskip
(Checked numerically at all primes up to $2^{21}$ for ten pairs,
including $(-1,0)$, $(-2,1)$, $(-3,-1)$, $(0,0.01)$, by
\texttt{verify\_small.py}.)

\subsection{General pairs of weights: the same reduction}

For real $m<m'$ and $\kappa=m'-m$ put
\begin{gather*}
I_{m,m'}(x):=D_m(x)-x^{-\kappa}D_{m'}(x)
=\sum_{p\le x}\chi^*(p)\,p^m\Bigl(1-\Bigl(\frac px\Bigr)^\kappa\Bigr),\\
\Phi_{m,m'}(t):=\sum_{p\le t}\chi^*(p)\log p\cdot
p^m\Bigl(1-\Bigl(\frac pt\Bigr)^\kappa\Bigr),
\end{gather*}
the second being the $\log$-weighted Riesz mean.

\begin{lemma}[no hypothesis]\label{lem:pairs}
For $x>p_0$,
\[
I_{m,m'}(x)=\frac{\Phi_{m,m'}(x)}{\log x}
+\int_{p_0}^x\Phi_{m,m'}(t)\,w_x(t)\,dt,
\qquad
w_x(t)=\frac{1+(t/x)^\kappa}{t\log^2t}
+\frac{2\,(1-(t/x)^\kappa)}{\kappa\,t\log^3t}\;>\;0 .
\]
In particular, if $\Phi_{m,m'}\ge0$ on $[p_0,x]$ then
$I_{m,m'}(x)\ge\Phi_{m,m'}(x)/\log x$.
\end{lemma}

\begin{proof}
With $A(t)=\sum_{p\le t}\chi^*(p)p^{m'}\log p$:
$I(x)=\int(t^{-\kappa}-x^{-\kappa})/\log t\;dA(t)$ and
$d\Phi(t)=\kappa t^{-\kappa-1}A(t)\,dt$, $\Phi(p_0)=0$. Two
integrations by parts give
$I(x)=\int g\,d\Phi=g(x)\Phi(x)-\int\Phi g'$ with
$g(t)=1/\log t+(1-(t/x)^\kappa)/(\kappa\log^2t)$, $g(x)=1/\log x$,
$-g'=w_x$.
\end{proof}

(Checked numerically to $23$ digits for the pairs $(0,1)$, $(1,2)$,
$(0.5,3)$, $(-1,0.5)$ at $x=700.3$.) For $(m,m')=(0,1)$ this is
\eqref{eq:keyid}, since $\Phi_{0,1}(t)=\Theta_1(t)/t$. So for
\emph{any} pair the problem reduces, without loss, to the sign of
$\Phi_{m,m'}$, which has an absolutely convergent explicit formula.
For pairs in the cone of Theorem~\ref{thm:cone} nothing more is
needed: by Lemma~\ref{lem:ingham}, applied with
$F(u)=u^m(1-(u/t)^\kappa)$ for each $t$,
$\Phi_{m,m'}(t)\ge\kappa\,t^{m-1}\Theta_1(t)\ge0$. For pairs outside
the cone, making $x_0$ explicit requires a certified positive margin
$\lambda_\chi(M,M')$, explicit lower-order terms, and a finite
computation of the sign of $\Phi_{m,m'}$; this was not carried out
(Section~\ref{sec:gaps}).

\subsection{Cross-check of the analytic bound against the sieve}

In the following table (from the output file \texttt{bounds.out}),
``$B(x)$'' is
the bound of Theorem~\ref{thm:cesaro} for $\cN(x)$ with $t_1$
optimised, and the last column is the resulting lower bound for
$\cM(x)$; the ``exact'' columns come from the sieve.

\begin{center}
\small
\begin{tabular}{llllll}
\toprule
$x$ & $\cL(x)$ & $B(x)$ & $\cN(x)$ exact & $\cM(x)$ exact &
bound for $\cM$\\
\midrule
\multicolumn{6}{l}{$\chi_4$}\\
$10^4$ & $0.06815$ & $0.06815$ & $0.830478$ & $9.0168$ & $0.74$\\
$10^5$ & $0.23659$ & $0.25661$ & $0.825252$ & $22.66735$ & $7.05$\\
$10^6$ & $0.34322$ & $0.37807$ & $0.878495$ & $63.587553$ & $27.37$\\
$10^7$ & $0.40862$ & $0.45215$ & $0.867385$ & $170.175890$ & $88.71$\\
$10^8$ & $0.44822$ & $0.49582$ & $0.777441$ & $422.047828$ & $269.17$\\
$10^9$ & $0.47217$ & $0.52089$ & $0.710167$ & $1083.682422$ &
$794.86$\\
$10^{10}$ & $0.48673$ & $0.53484$ & $0.746014$ & $3239.897909$ &
$2322.78$\\
\midrule
\multicolumn{6}{l}{$\chi_3$}\\
$10^6$ & $0.35882$ & $0.39320$ & $0.853793$ & $61.799590$ & $28.46$\\
$10^8$ & $0.48366$ & $0.53415$ & $0.796002$ & $432.124105$ &
$289.97$\\
$10^{10}$ & $0.52739$ & $0.57926$ & $0.798797$ & $3469.131474$ &
$2515.72$\\
\bottomrule
\end{tabular}
\end{center}

The bound lies below the truth everywhere, at $43\%$ ($10^6$), $64\%$
($10^8$), $72\%$ ($10^{10}$) of the exact value of $\cM$ for
$\chi_4$. Against the \emph{minimum} of $\cN$ in each dyadic range
the ratio rises to $0.764$ ($\chi_4$, at $x=1\,070\,765\,639$) and
$0.800$ ($\chi_3$). Also $\cL(2^k)\le G_{\mathrm{low}}(k)$, a lower
bound for $G$ on $[2^k,2^{k+1})$ from the sieve, in every dyadic
range $k=14,\dots,33$, for both characters and both variants (a
consistency test; all passed). The loss at $x=10^{10}$ ($\chi_4$)
comes from: $\sigma_0^2$ being a worst-case bound for an oscillating
sum ($0.156$ is subtracted from $\tfrac23$, whereas the true value is
$G(10^{10})=0.6485$); the odd prime powers being bounded trivially
($0.019$); the prime squares ($0.005$). The analytic argument proves
$\Theta_1(t)>0$ from $t=4892$ ($\chi_4$) and $t=8526$ ($\chi_3$); the
theorem needs the computation only up to $10^4$, where it is done in
exact integer arithmetic; the floating-point computation was carried
to $10^{10}$, so the two ranges overlap by six orders of magnitude.

\begin{remark}[one class against $\mathrm{li}$]\label{rem:oneclass}
The problem list of Hamieh, Kadiri, Martin and Ng \cite{HKMN24}
contains (as Problem~23, attributed there to Johnston) the question,
for $A_1^\pi(x;q,a):=\int_0^x(\varphi(q)\pi(t;q,a)-\mathrm{li}(t))\,dt$,
of classifying under GRH those $q$ and $a$ for which
$A_1^\pi(x;q,a)<0$ for all sufficiently large $x$. That problem
concerns one class against $\mathrm{li}$, not the difference of two
classes; the link is the identity
$A_1^\pi(x;4,3)-A_1^\pi(x;4,1)=2x\,\cM(x)$. A consequence of
Theorem~\ref{thm:cesaro} together with \cite[Theorem~1.1]{Joh23}:
\emph{under RH and GRH$(\chi_4)$},
\[
\int_2^x\bigl(2\pi(t;4,1)-\mathrm{li}(t)\bigr)\,dt
=\int_2^x\bigl(\pi(t)-\mathrm{li}(t)\bigr)\,dt-(x-2)-x\,\cM(x)\;<\;0
\qquad\text{for every }x>2,
\]
and the same holds for the class $1$ mod $3$ under RH and
GRH$(\chi_3)$, with $-(x-3)^+$ in place of $-(x-2)$. The classes
$3$ mod $4$ and $2$ mod $3$ are not decided this way (the two biases
cancel to leading order).
\end{remark}

\section{The Landau converse: equivalence with GRH}\label{sec:landau}

\begin{theorem}[equivalence]\label{thm:equiv}
Let $q\in\{4,3\}$. The following are equivalent:
\begin{enumerate}
\item[(a)] GRH$(\chi_q)$;
\item[(b)] $\int_0^xD_0(t)\,dt>0$ for every $x>p_0$ (for $q=4$:
$\int_0^x(\pi(t;4,3)-\pi(t;4,1))\,dt>0$ for every $x>3$);
\item[(c)] $\int_0^xD_0(t)\,dt\ge0$ for all sufficiently large $x$.
\end{enumerate}
\end{theorem}

The implication (c)$\Rightarrow$(a) is Landau's argument and needs no
hypothesis; it is the case $(m,m')=(0,1)$ of
\cite[the corollary ``equivalence with GRH for $q\in\{3,4\}$'', part~(a)]{MC}, and its classical prototype for this
very character is Hardy--Littlewood and Landau
\cite{HL18,Lan18}: Chebyshev's assertion
$\sum_{p>2}(-1)^{(p-1)/2}e^{-xp}\to-\infty$ as $x\to0$ is equivalent
to GRH$(\chi_4)$ (Abel mean and limit $-\infty$ there; Ces\`aro mean
and sign for every $x$ here). We include the proof for completeness.

\begin{lemma}[Landau; {\cite[Lemma~15.1]{MV07}}]\label{lem:landau}
Let $A:[1,\infty)\to\mathbb R$ be bounded and Riemann integrable on
every finite interval, with $A(x)\ge0$ for all $x>X_0$. Let
$\sigma_c$ be the infimum of the real $\sigma$ for which
$\int_{X_0}^\infty A(x)x^{-\sigma-1}\,dx<\infty$. Then
$\mathcal F(s):=\int_1^\infty A(x)x^{-s-1}\,dx$ is analytic in the
half-plane $\re s>\sigma_c$, and, if $\sigma_c>-\infty$,
$\mathcal F$ has no analytic continuation to a neighbourhood of the
point $s=\sigma_c$.
\end{lemma}

(This is the form in which the lemma is stated in
\cite[the lemma quoting Landau's theorem]{MC}; in \cite{MV07} the kernel is $x^{-s}$.)

\begin{proof}[Proof of Theorem~\ref{thm:equiv}]
(a)$\Rightarrow$(b) is Theorem~\ref{thm:cesaro}; (b)$\Rightarrow$(c)
is trivial. Assume (c), i.e.\ $\cM(x)\ge0$ for $x\ge x_0$.

\emph{Step 1 (Mellin transform).} The function $\cM$ is continuous
and $|\cM(x)|\le\pi(x)\le x$. For $\re s>1$, integrating term by
term,
\[
\mathcal F(s):=\int_1^\infty\cM(x)\,x^{-s-1}\,dx
=\sum_p\chi^*(p)\int_p^\infty(x-p)\,x^{-s-2}\,dx
=\frac{P(s)}{s(s+1)},
\]
where $P(s):=\sum_p\chi^*(p)\,p^{-s}$.
From the Euler product,
$\log L(s,\chi)=\sum_p\sum_{k\ge1}\chi(p)^k/(kp^{ks})$, so for
$\re s>1$
\begin{equation}\label{eq:Pstructure}
P(s)=-\log L(s,\chi)+\tfrac12\log L(2s,\chi_0)+H(s),
\qquad
L(2s,\chi_0)=\zeta(2s)\bigl(1-\ell^{-2s}\bigr),
\end{equation}
where $\chi_0$ is the principal character mod $q$ and
$H(s)=\sum_p\sum_{k\ge3}c_k(p)p^{-ks}$ (bounded coefficients) is
analytic in $\re s>\tfrac13$.

\emph{Step 2 (the real axis).} On the real segment $\sigma>\tfrac12$,
$\zeta(2\sigma)(1-\ell^{-2\sigma})>0$, and $L(\sigma,\chi)>0$ by
\eqref{eq:Lpositive}. The zeros of $L(\cdot,\chi)$ form a closed
discrete set disjoint from the ray $(\tfrac12,\infty)$, so there is a
continuous $r:(\tfrac12,\infty)\to(0,\infty)$ such that
$L(w,\chi)\ne0$ on
$U:=\{\sigma+it:\sigma>\tfrac12,\ |t|<r(\sigma)\}$. The region $U$ is
simply connected, contains part of the half-plane $\re w>1$, and its
intersection with any half-plane $\re w>c$ is connected. On $U$,
$L(\cdot,\chi)$ has an analytic logarithm which is real on the real
ray and therefore extends $\log L(w,\chi)$ from $U\cap\{\re w>1\}$;
$L(2w,\chi_0)$ is analytic and zero-free on $\re w>\tfrac12$; and
$H$ is analytic there. By \eqref{eq:Pstructure}, $P$ extends
analytically to $U$, and $\mathcal F_1(s):=P(s)/(s(s+1))$ is analytic
on $U$.

\emph{This is where the absence of real zeros is used, and it cannot
be dropped:} a real zero $\beta>\tfrac12$ of $L(s,\chi)$ would give
$P$ a logarithmic singularity at $\beta$ with $P\to+\infty$ as
$s\to\beta^+$, which is perfectly consistent with $\cM\ge0$;
Landau's theorem would then only give $\sigma_c=\beta$. For $\chi_4$
and $\chi_3$ the elementary positivity \eqref{eq:Lpositive} settles
it; for a general real character the same proof needs ``no real zero
in $(\tfrac12,1)$'' as a hypothesis (compare
\cite[the Landau-type converse and the remark on the exceptional-zero caveat]{MC} and \cite[Remark~1]{Zhao25}).

\emph{Step 3 (Landau).} Let $\sigma_c$ be as in
Lemma~\ref{lem:landau} with $A=\cM$. By Step~1, $\sigma_c\le1$ and
$\mathcal F=\mathcal F_1$ on $\re s>1$. Suppose
$\sigma_c>\tfrac12$. Then $\mathcal F$ is analytic on
$\re s>\sigma_c$, $\mathcal F_1$ is analytic on $U$, and they agree
on the connected set $U\cap\{\re s>\sigma_c\}$, because they agree on
its nonempty open subset $U\cap\{\re s>1\}$. Hence $\mathcal F_1$ is
an analytic continuation of $\mathcal F$ to a neighbourhood of
$s=\sigma_c\in U$, contradicting Lemma~\ref{lem:landau}. Therefore
$\sigma_c\le\tfrac12$: $\mathcal F$, and with it
$P(s)=s(s+1)\mathcal F(s)$, is analytic in $\re s>\tfrac12$.

\emph{Step 4 (conclusion).} In $\re s>\tfrac12$, $L(2s,\chi_0)$ is
analytic and zero-free and $H$ is analytic, so
$\mathcal G(s):=-P(s)+\tfrac12\log L(2s,\chi_0)+H(s)$ is analytic
there, and $e^{\mathcal G(s)}=L(s,\chi)$ for $\re s>1$. By the
identity theorem $e^{\mathcal G}=L(\cdot,\chi)$ on $\re s>\tfrac12$,
so $L(s,\chi)$ has no zero with $\re s>\tfrac12$. By the functional
equation of the real primitive character $\chi$ (the zeros are
symmetric about $\re s=\tfrac12$) there is none with
$0<\re s<\tfrac12$ either.
\end{proof}

\section{LI-free monotonicity of the densities, and ties}
\label{sec:ties}

For $m>-\tfrac12$ write
\[
\delta^{>}_{-}(m)=\dlow\{D_m>0\},\quad
\delta^{>}_{+}(m)=\dup\{D_m>0\},\quad
\delta^{\ge}_{-}(m)=\dlow\{D_m\ge0\},\quad
\delta^{\ge}_{+}(m)=\dup\{D_m\ge0\}.
\]

\begin{theorem}\label{thm:ties}
Let $\chi$ be a real primitive character. Parts \textup{(a)} and
\textup{(b)} assume GRH$(\chi)$ only; part \textup{(c)} assumes
GRH$(\chi)$ and NCZ$(\chi)$. No LI is assumed. Let $-\tfrac12<m<m'$
with $\lambda_\chi(M,M')>0$, where $\lambda_\chi$ is defined by
\eqref{eq:deflambda} (with the term $1+2n_0$). Sufficient:
$\sigma_M\sigma_{M'}<2$; in particular $M'\le M_\sigma(\chi)$.
\begin{enumerate}
\item[(a)] (\cite[the corollary on densities without LI]{MC}.) For $*\in\{-,+\}$:
$\ \delta^{>}_{*}(m')\le\delta^{\ge}_{*}(m')\le\delta^{>}_{*}(m)
\le\delta^{\ge}_{*}(m)$. The same holds for natural densities.
\item[(b)] \emph{Ties.} Let $I\subset(-\tfrac12,\infty)$ be an
interval of weights such that $\lambda_\chi(M,M')>0$ for all
$m<m'$ in $I$ (for instance
$I=(-\tfrac12,M_\sigma(\chi)-\tfrac12]$). Then each of the four
functions $\delta^{>}_{\pm}$, $\delta^{\ge}_{\pm}$ is nonincreasing on
$I$, and $\delta^{>}_{*}(m)=\delta^{\ge}_{*}(m)$ for all but countably
many $m\in I$ $(*=\pm)$: outside a countable subset of $I$, counting
ties for or against the leader does not change the upper or the lower
logarithmic density. Moreover the tie sets
$\{x\ge x_0:D_m(x)=0\}$, $m\in I$, are pairwise disjoint beyond a
point depending on the pair, so for any finitely many
$m_1,\dots,m_k\in I$, $\sum_j\dlow\{D_{m_j}=0\}\le1$; in particular
at most countably many $m\in I$ have a tie set of positive lower
logarithmic density.
\item[(c)] For $m<m'$ as above,
$\delta^{>}_{-}(m)\ge\delta^{\ge}_{-}(m')\ge1/(1+V_{M'})$.
\end{enumerate}
\end{theorem}

\begin{proof}
(a) By Theorem~\ref{thm:MCA}(ii),
$\{x\ge x_0:D_{m'}(x)\ge0\}\subset\{D_m>0\}$; densities are monotone
under inclusion and ignore bounded sets. (b) Monotonicity is (a). By
(a), for $m<m'$ in $I$:
$\delta^{\ge}_{*}(m')\le\delta^{>}_{*}(m)\le\delta^{\ge}_{*}(m)$, so
the monotone functions $\delta^{>}_{*}$ and $\delta^{\ge}_{*}$
interlace; hence they have the same one-sided limits at every
interior point and coincide wherever one of them is continuous, i.e.\
off a countable set. For $m<m'$ in $I$ and $x\ge x_0(m,m')$,
$D_{m'}(x)=0$ forces $D_m(x)>0$, so the tie sets are eventually
disjoint; lower densities are superadditive on disjoint sets. (c) is
(a) and Theorem~\ref{thm:cantelli}, since
$\dlow\{D_{m'}\ge0\}\ge\dlow\{D_{m'}>0\}\ge1/(1+V_{M'})$.
\end{proof}

We cannot exclude under GRH alone that $\nu_M$ has an atom at $0$ for
some $M$; the theorem shows that this can affect at most countably
many weights in $I$. For $q\in\{4,3\}$ and the pairs of the cone of
Theorem~\ref{thm:cone} intersected with $m>-\tfrac12$, i.e.\
$-\tfrac12<m\le0$, $m<m'\le1$, the inclusions used in the proof hold
with $x_0$ replaced by any number $>p_0$.

\begin{definition}\label{def:selfsuff}
An ordinate $\gamma_0\in\cZ$ is \emph{self-sufficient} if $\gamma_0$
is not a rational linear combination of $\cZ\setminus\{\gamma_0\}$ ---
a far weaker hypothesis than LI, though equally unproved. (This is
\cite[Definition~5.3]{Dev20}: it is formulated on the \emph{set} of
distinct ordinates, so $\gamma_0$ may be a multiple zero.
\cite[Definition~1.3]{MN20}, from which the name comes, is on the
\emph{multiset} of ordinates and forces the zero to be simple.)
\end{definition}

\begin{proposition}[one self-sufficient ordinate]\label{prop:selfsuff}
Assume GRH$(\chi)$ and NCZ$(\chi)$, and suppose that $\cZ$ contains a
self-sufficient ordinate $\gamma_0$. Then for every $M>0$:
\begin{enumerate}
\item[(a)] $\nu_M$ is absolutely continuous; hence
$\delta(m):=\nu_M((0,\infty))$ is the logarithmic density of
$\{D_m>0\}$, of $\{D_m\ge0\}$ and of $\{R_m>0\}$; ties have
logarithmic density zero.
\item[(b)] $m\mapsto\delta(m)$ is continuous on $(-\tfrac12,\infty)$,
$0<\delta(m)<1$, and $\delta$ is nonincreasing on every interval $I$
as in Theorem~\ref{thm:ties}(b).
\item[(c)] $\delta(m)\ge1/(1+V_M)$, and all bounds of
Section~\ref{sec:biasbounds} hold for $\delta(m)$ itself (in
particular, by Theorem~\ref{thm:NRq}, $\delta(0)>0.8842$ for $q=4$
and $\delta(0)>0.9112$ for $q=3$, when $\chi_q$ has a self-sufficient
ordinate).
\end{enumerate}
\end{proposition}

\begin{proof}
(a) Self-sufficiency means that every integer relation
$\sum_{\gamma\in\cZ}n_\gamma\gamma=0$ has $n_{\gamma_0}=0$. Hence the
annihilator of $\cA$ in the dual of $\mathbb T\times\cA'$ ($\cA'$ the
closure of the one-parameter group in
$\mathbb T^{\cZ\setminus\{\gamma_0\}}$) is trivial, so
$\cA=\mathbb T\times\cA'$ and $\hA$ is the product measure. Thus
$F_M=\re(m_{\gamma_0}r_{\gamma_0}\omega_{\gamma_0})+F'(\omega')$ with
independent summands, the first having an (absolutely continuous)
arcsine law; a convolution with an absolutely continuous law is
absolutely continuous. Then $\nu_M(\partial(0,\infty))=0$ and the
portmanteau theorem give the densities. (b) $M\mapsto F_M\in L^2(\cA)$
is continuous (dominated convergence in \eqref{eq:FM}), so
$M\mapsto\nu_M$ is weakly continuous, and $\nu_M(\{0\})=0$;
positivity is Theorem~\ref{thm:leadsets}; monotonicity is
Theorem~\ref{thm:ties}. (c) Theorems~\ref{thm:cantelli}
and~\ref{thm:NR}.
\end{proof}

Devin \cite[Theorem~2.2]{Dev20} (no atoms, and the density exists)
and Bailleul \cite[Theorem~1.7]{Bai21} treat the unweighted race
under this type of hypothesis; Devin proves continuity of the
limiting measure, the absolute continuity in (a) is our (easy)
addition. Martin--Ng \cite[Theorem~1.5(a)]{MN20} is a two-way
statement: three self-sufficient ordinates (in their multiset sense)
imply that the race is weakly inclusive, i.e.\ that the logarithmic
densities exist. Symmetry of $\nu_M$ about $1$ (hence
$\delta\ge\tfrac12$) would follow from Devin's parity condition
\cite[Theorem~2.3]{Dev20}. The \emph{strict} decrease of the density
(\cite[the coupling monotonicity theorem]{MC}) uses the freedom to rotate each phase
separately and genuinely needs LI; under GRH alone only the weak
monotonicity above survives.

\section{The zeta analogue (statements; proofs parallel)}
\label{sec:zeta}

\noindent\textbf{Status of this section.} The statements below are
the analogues, for $\sum_{p\le x}p^m$ against
$\mathrm{li}(x^{m+1})$, of results proved above or in \cite{MC} for
characters. They are taken from a working note that has \emph{not}
been separately refereed. We include only statements that follow by
the same proofs, indicate which proof is meant, and do not write the
proofs out. The numerical values are \emph{computed values}, with the
accuracy stated; none is certified in interval arithmetic. The
hypotheses here are RH (for $\zeta$) and, where stated, LI for the
ordinates of $\zeta$.

\subsection{Set-up}
For $m>-1$ let
\[
\pi_m(x):=\sum_{p\le x}p^m,\qquad
I_m(x):=\int_2^x\frac{t^m}{\log t}\,dt
=\mathrm{li}(x^{m+1})-\mathrm{li}(2^{m+1})
\]
(substitute $u=t^{m+1}$), and, with $M=m+\tfrac12$,
\[
\Delta_m:=I_m-\pi_m,\qquad
G^\zeta_m(x):=x^{-M}(\log x)\,\Delta_m(x).
\]
Replacing $I_m$ by $\mathrm{li}(x^{m+1})$ changes $G^\zeta_m$ by
$O(x^{-M}\log x)=o(1)$ for $M>0$, so the statements below hold
verbatim with $\mathrm{li}(x^{m+1})$. Under RH let
$\rho=\tfrac12+i\gamma$ run over the nontrivial zeros of $\zeta$ and
\[
\lambda_\zeta(M,M'):=\frac1{MM'}
-\sum_{\gamma>0}\frac4{|M+i\gamma|\,|M'+i\gamma|},
\qquad
\sigma_{\zeta,M}^2:=\sum_{\gamma>0}\frac{8M^2}{M^2+\gamma^2}
=4M\,\frac{\xi'}{\xi}\bigl(\tfrac12+M\bigr),
\]
the last identity by the Hadamard product for $\xi$ (parallel to
Proposition~\ref{prop:WRvar}).

\subsection{Statements under RH alone}

\noindent\textbf{(Z1) Exact identity (unconditional).} For
$-1<m<m'$, $\kappa=m'-m$ and $x\ge2$,
\[
\Delta_m(x)=x^{-\kappa}\Delta_{m'}(x)
+\kappa\int_2^xt^{-\kappa-1}\Delta_{m'}(t)\,dt .
\]
\emph{Parallel proof:} Theorem~\ref{thm:resolvent}(a) and
Theorem~\ref{thm:MCA}(i) (partial summation;
$d\pi_m=t^{-\kappa}d\pi_{m'}$ and $dI_m=t^{-\kappa}dI_{m'}$).

\smallskip\noindent\textbf{(Z2) Inheritance (RH; no LI, no
simplicity).} Let $-\tfrac12<m<m'$. As $x\to\infty$,
\[
G^\zeta_m(x)-G^\zeta_{m'}(x)=\frac{\kappa}{2MM'}
+\kappa\sum_\rho\frac{x^{i\gamma}}{(M+i\gamma)(M'+i\gamma)}
+O\Bigl(\frac1{\log x}\Bigr),
\]
the series over all nontrivial zeros converging absolutely;
consequently
$G^\zeta_m-G^\zeta_{m'}\ge\tfrac\kappa2\lambda_\zeta(M,M')-O(1/\log x)$.
If $\lambda_\zeta(M,M')>0$, there is $x_0$ such that for all
$x\ge x_0$
\[
\pi_{m'}(x)\le\mathrm{li}(x^{m'+1})\ \Longrightarrow\
\pi_m(x)<\mathrm{li}(x^{m+1}),
\]
equivalently
$\pi_m(x)\ge\mathrm{li}(x^{m+1})\Rightarrow
\pi_{m'}(x)>\mathrm{li}(x^{m'+1})$. In words: ``$\mathrm{li}$ ahead''
is inherited downward in the weight, and ``$\pi$ ahead'' (the
Skewes-type event) upward. \emph{Parallel proof:}
\cite[Theorem~A]{MC} (Theorem~\ref{thm:MCA}), with the truncated
explicit formula for $\theta(t)$ under RH in place of that for
$\psi(t,\chi)$; the term $-\sqrt t$ of $\theta(t)$ produces the bias
$\kappa/(2MM')$, and the pole of $\zeta$ at $s=1$ produces $I_m$
exactly. The threshold $x_0$ is not effective as stated.

\smallskip\noindent\textbf{(Z3) Interval criterion (RH).}
(a) Since $MM'\lambda_\zeta(M,M')
=1-\sum_{\gamma>0}4\frac{M}{|M+i\gamma|}\frac{M'}{|M'+i\gamma|}$ is
strictly decreasing in each variable,
$MM'\lambda_\zeta(M,M')>M'^2\lambda_\zeta(M',M')
=1-\tfrac12\sigma_{\zeta,M'}^2$ for $M<M'$; so
$\lambda_\zeta(M,M')>0$ whenever $\sigma_{\zeta,M'}^2\le2$ (compare
Theorem~\ref{thm:MCA}(iii)). (b) Every pair in $(0,b]$ satisfies
$\lambda_\zeta>0$ if and only if $\sigma_{\zeta,b}^2\le2$. (c)
Consequently, under RH alone, the upper and the lower logarithmic
densities of $\{x:\pi_m(x)>\mathrm{li}(x^{m+1})\}$ are nondecreasing
along any chain of pairs with $\lambda_\zeta>0$, in particular on
$m\in(-\tfrac12,M_\sigma(\zeta)-\tfrac12]$, where
$\sigma^2_{\zeta,M_\sigma(\zeta)}=2$ (parallel to
Theorem~\ref{thm:ties}(a)). \emph{Computed values:}
$M_\sigma(\zeta)=3.317706$, i.e.\ $m\le2.817706$ (closed form of
$\sigma_{\zeta,M}^2$; stated error $<10^{-6}$); and
\[
\lambda_\zeta(\tfrac12,\tfrac32)=+1.2411,\qquad
\lambda_\zeta(\tfrac32,\tfrac52)=+0.1749,\qquad
\lambda_\zeta(\tfrac52,\tfrac72)=+0.0232
\]
(from the first $10^4$ zeros of $\zeta$, enclosed by Arb with radius
$<10^{-20}$, plus a tail estimate from the Riemann--von Mangoldt
formula with an explicit bound for $S(T)$ quoted from Trudgian
\cite{Tru14}; stated tail radius $2\cdot10^{-7}$; the evaluation is
floating point). So the pairs $(m,m')=(0,1)$, $(1,2)$, $(2,3)$ are
covered by (Z2), and, under RH alone,
\[
\dlow^\zeta(0)\le\dlow^\zeta(1)\le\dlow^\zeta(2)\le\dlow^\zeta(3),
\qquad
\dup^\zeta(0)\le\dup^\zeta(1)\le\dup^\zeta(2)\le\dup^\zeta(3),
\]
where $\dlow^\zeta(m)$, $\dup^\zeta(m)$ are the lower and upper
logarithmic densities of $\{x:\pi_m(x)>\mathrm{li}(x^{m+1})\}$.

\begin{remark}\label{rem:zetaJohnston}
For $(m,m')=(0,1)$ the left-hand side of (Z1), rearranged, is
$\Delta_0(x)-\Delta_1(x)/x
=x^{-1}\int_2^x\bigl(\mathrm{li}(t)-\mathrm{li}(2)-\pi(t)\bigr)\,dt$,
which differs from the quantity of Johnston's theorem
\cite[Theorem~1.1]{Joh23} only by the term
$\mathrm{li}(2)(x-2)/x>0$. So for this pair (Z2) gives, for all
sufficiently large $x$, the inequality
$\int_2^x(\pi(t)-\mathrm{li}(t))\,dt<0$ that
\cite[Theorem~1.1]{Joh23} establishes under RH for every $x>2$.
\end{remark}

\subsection{Densities under RH and LI (computed values)}

\noindent\textbf{(Z4) Limiting distribution and densities (RH; LI
where stated).} Under RH, $2M\,G^\zeta_m(e^u)$ has a
limiting distribution of mean $1$ (parallel to
\cite[Theorem~1(i)]{WR} and Theorem~\ref{thm:limit}). Under RH and LI
for $\zeta$ it is the law of
$X^\zeta_M=1+\sum_{\gamma>0}c_\gamma(M)\cos\theta_\gamma$ with
independent uniform phases, it is absolutely continuous with support
$\mathbb R$, and
\[
\delta^\zeta(m):=\text{logarithmic density of }
\{x:\pi_m(x)>\mathrm{li}(x^{m+1})\}=\PP(X^\zeta_M<0)
\]
exists (parallel to \cite[Theorem~1(ii)]{WR}; for $m=0$ this is
Rubinstein--Sarnak \cite{RS94}). Moreover, under RH and LI,
$m\mapsto\delta^\zeta(m)$ is strictly \emph{increasing} on
$(-\tfrac12,M_0(\zeta)-\tfrac12]$, where $M_0(\zeta)$ is the root of
$\sum_{\gamma>0}4u^3(u^2+\gamma^2)^{-3/2}=1$ (parallel to
\cite[the coupling monotonicity theorem]{MC}, with the event $\{X<0\}$ in place of
$\{X>0\}$); computed value $M_0(\zeta)=7.536853\pm10^{-8}$.

\smallskip\noindent\emph{Computed values of the densities} (RH and
LI):
\begin{center}
\begin{tabular}{lcccc}
\toprule
$m$ & $0$ & $1$ & $2$ & $3$\\
\midrule
$\delta^\zeta(m)$ & $2.62997\cdot10^{-7}$ & $0.0605587$ &
$0.177311$ & $0.253681$\\
rounded & $2.63\cdot10^{-7}$ & $0.0606$ & $0.1773$ & $0.2537$\\
\bottomrule
\end{tabular}
\end{center}
These were obtained by numerical inversion of the characteristic
function $\prod_{\gamma\le T}J_0(c_\gamma t)\cdot
e^{-v_{\mathrm{tail}}t^2/2}$ (Gil--Pelaez formula, trapezoid rule with
step $0.05$, the first $10^4$ zeros, and a Gaussian factor carrying
the exact variance of the remaining zeros; the script is
\texttt{zeta\_density.py} in \texttt{monotonicity/extensions/}).
\emph{Stated accuracy:} the
working note's error budget (Gaussian replacement of the zeros beyond
the $10^4$-th: $\le7\cdot10^{-11}$; truncation in $t$: $<10^{-18}$;
aliasing: negligible) gives an error below $10^{-9}$ on every entry;
halving the step changes nothing to $10^{-15}$; a Monte Carlo run
with $10^6$ samples agrees within $1.4$ standard deviations at
$m=1,2,3$. This budget has not been refereed and the computation is
not in interval arithmetic. The entry at $m=0$ reproduces the value
$2.6\cdot10^{-7}$ of \cite{RS94}. So, under RH and LI, the sum of the
primes up to $x$ exceeds $\mathrm{li}(x^2)$ on about $6.06\%$ of the
logarithmic scale.

Sums of powers of primes have been compared with $\pi(x^{k+1})$,
rather than with $\mathrm{li}(x^{k+1})$, in \cite{GW18,Was24} and, in
arithmetic progressions, in \cite{BBLMR24}; we do not discuss that
comparison here.

\section{Computations and their rigor class}\label{sec:comp}

All scripts are in the repository
\begin{center}
\url{https://github.com/raghavs1729/weighted-prime-races};
\end{center}
paths are relative to its root. Each runs on one core. We distinguish five
rigor classes: \textbf{[exact]} exact integer or rational arithmetic;
\textbf{[ball]} ball arithmetic (Arb), outward rounded;
\textbf{[ld]} \texttt{long double} with margins stated;
\textbf{[double]} double precision or multiprecision floating point
without error control; \textbf{[search]} a search whose output is
verified separately.

\medskip\noindent\textbf{Zero lists (used in
Sections~\ref{sec:variance}, \ref{sec:biasbounds}).}
\begin{itemize}
\item \texttt{prototype/explore/chi4\_zeros.txt},
\texttt{chi3\_zeros.txt}: $511$ and $537$ ordinates with radius
$10^{-20}$. \texttt{prototype/explore/certify\_arb.py --K 0}
\textbf{[ball]}: completeness and simplicity up to
$T_4=640.165707$, $T_3=700.024747$ by sign changes and a winding
count. (Rerun by the referee of the working notes: winding balls
$[511\pm3\cdot10^{-12}]$, $[537\pm3\cdot10^{-11}]$.) These lists are
what the tag \certA{} refers to.
\end{itemize}

\noindent\textbf{Directory \texttt{monotonicity/grhonly/}
(Sections~\ref{sec:resolvent}--\ref{sec:biasbounds}).}
\begin{itemize}
\item \texttt{certified\_bounds.py} $\to$
\texttt{certified\_bounds\_results.txt} \textbf{[ball]}: $C(\chi)$,
$\sigma_M^2$, the exact tails, the bounds of
Lemmas~\ref{lem:mult}, \ref{lem:JA}, $V_M^+$, the Cantelli and
soft-threshold constants (\certA), and tier~B (\certB). It
establishes the numerical table of Section~\ref{sec:variance}, the
constants of Theorem~\ref{thm:cantelli-q} and the table after it.
\item \texttt{nonresonance.py} $\to$
\texttt{nonresonance\_results.txt} \textbf{[search]}: floating-point
search for a good $(\varkappa_0,b)$; margins of (NR1)--(NR3) in
double precision. No rigor is claimed for it.
\item \texttt{nonresonance\_certified.py} $\to$
\texttt{nonresonance\_certified\_results.txt}: (NR1)--(NR3) in exact
integer arithmetic \textbf{[exact]}; Arb enclosure of $J$ at a
rational feasible point \textbf{[ball]}; a grid self-test of the
closed formulas for $\int Q^2$ and $\int\varkappa^2$ for $N=1,2,3$
(maximal deviation $1.1\cdot10^{-16}$) \textbf{[double]}. It
establishes Theorem~\ref{thm:NRq} and the non-resonance rows of
Table~\ref{tab:bias} (\certA/\certB).
\item \texttt{sieve\_checks.py} $\to$
\texttt{sieve\_checks\_results.txt} \textbf{[double]}, sanity only:
sieve to $10^7$; the identities of Theorem~\ref{thm:resolvent}(a)
and Remark~\ref{rem:exactunlogged} (relative residuals
$\le1.3\cdot10^{-13}$ for $q=4,3$ and three weight pairs);
boundedness of $y\cdot(E_M-\int E_{M'}d\mu)$ and of
$y\cdot(E_M-\Ein_M)$ on $[10^3,10^7]$ (about $2.0$--$2.4$ and
$2.8$--$3.4$ for $(m,m')=(0,1)$; for $(-\tfrac14,\tfrac12)$ about
$6$--$7$ and $7.5$--$8.3$ and, for $q=4$, still slowly increasing at
$10^7$, consistent with a constant of size $1/M$ but not a
verification of the rate); Theorem~\ref{thm:resolvent}(c) (prime
side minus zero side: $-0.103$, $-0.048$, $-0.027$, $-0.010$ at
$x=10^4,\dots,10^7$ for $q=4$); empirical variances over
$y\in[\log10^3,\log10^7]$ ($0.0489$, $0.1598$, $0.3508$, $0.6169$,
$1.3631$ at $M=0.25,0.5,0.75,1,1.5$ against
$\sigma_M^2=0.0390$, $0.1556$, $0.3487$, $0.6169$, $1.3686$). At
$M=0.1$ the finite range breaks the ordering of some empirical
moments (the mean over this window is $0.57$, not $1$); this is a
finite-$x$ effect, not a counterexample: the theorems concern
$Y\to\infty$.
\item \texttt{explore\_degree2.py} $\to$
\texttt{explore\_degree2\_results.txt} \textbf{[double]}:
exploration of degree-two multipliers
(Remark~\ref{rem:NRlimits}(1)); establishes nothing.
\end{itemize}

\noindent\textbf{Directory \texttt{monotonicity/explicit/}
(Sections~\ref{sec:explicit}, \ref{sec:landau}).}
\begin{itemize}
\item \texttt{cert\_small.py} $\to$ \texttt{cert\_small.out}
(about $140$~s). This is the certified companion of the sieve and
contains \emph{everything finite that the proofs use}:
\begin{itemize}
\item[(A)] \textbf{[exact]} $\Theta_1(p)>0$ at all $1227$ primes
$p\nmid q$ with $p_0<p\le10^4$, as an inequality between integers,
for both characters (used in Step~7 of the proof of
Theorem~\ref{thm:Theta1});
\item[(B)] \textbf{[ball]} $\Theta_1(p)>0$ for $p\le2^{22}$ in
$256$-bit Arb (not needed);
\item[(C)] \textbf{[ball]} Lemma~\ref{lem:theta}(a) for
$y\le2^{22}$, including the extremum at $1423$ (needed in
variant~C; part of the main-text $Y_1=10^{10}$);
\item[(D)] \textbf{[ball]}, with exact rational $\cM$:
$\cN>\tfrac12$ on $[23,29]$ and on all gaps with
$29\le p<p'\le2^{22}$ ($\chi_4$); on $[6,11]$ and all gaps with
$11\le p<p'\le2^{22}$ ($\chi_3$);
\item[(E)] \textbf{[ball]} ($64$-bit balls; $D_0$, $D_1$ exact) the
same gap bounds up to $1.6\cdot10^8$ ($\chi_4$) and $2.6\cdot10^7$
($\chi_3$) (used in Corollary~\ref{cor:half}).
\end{itemize}
All five parts pass for both characters.
\item \texttt{cesaro\_scan.c} $\to$ \texttt{out\_q4\_1e10.txt},
\texttt{out\_q3\_1e10.txt} (segmented sieve to $10^{10}$, about
$70$~s per character). $D_0$ (64-bit integers) and $D_1$ (128-bit
integers) are exact, and the test $\cM(p)>0$ is the exact integer
test $p\,D_0(p)-D_1(p)>0$ \textbf{[exact]}; since $\cM$ is
continuous and monotone between consecutive primes, positivity at
all primes is positivity everywhere. $\cN$ and its gap bounds are
formed from the exact $D_0$, $D_1$ by a fixed handful of
\texttt{long double} operations, with no accumulation
\textbf{[ld]}. The quantities $\Theta$, $\theta$ and
$\sum\chi^*(p)p\log p$ are \emph{accumulated} in \texttt{long
double} ($64$-bit mantissa) \textbf{[ld]}: worst-case accumulated
rounding at $10^{10}$ is $\le0.21$ in $\theta$ and $\le6\cdot10^4$
in $\Theta_1$ (which is about $6\cdot10^{14}$ there). True margins of
the quantities tested: $\Theta_1(p)\ge1.1756$ ($\chi_4$, at $p=7$)
resp.\ $2.0794$ ($\chi_3$, at $p=5$), $G\ge0.13$ on $[11,10^{10}]$
and $G\ge0.52$ on $[2^{13},10^{10}]$;
$(y-\vth(y))/\sqrt y\le1.98197$ resp.\ $1.99271$ for $y\ge1427$
against $c_1$; $\vth(y)-y\le-0.26\sqrt y$. The sign of
$\Theta_1(p)$ in this program is a floating-point decision and is
\emph{not} what the proof uses. What the proofs take from this
program is only Lemma~\ref{lem:theta}(a) on $(2^{22},10^{10}]$, and
not even that in variant~C.
\item \texttt{constants.py} $\to$ \texttt{constants.out}
\textbf{[double]} (\texttt{mpmath}, $40$ digits): the closed-form
constants $\sigma_0^2$, $\mathfrak b$, $\mathfrak b_0$ (using
$\frac{L'}{L}(1,\chi_4)=\gamma_{\mathrm E}+2\log2+3\log\pi
-4\log\Gamma(\tfrac14)$) and the checks of Lemma~\ref{lem:zerosq}
against the zero lists.
\item \texttt{bounds.py} $\to$ \texttt{bounds.out}
\textbf{[double]}: $\cL(t)$, the thresholds, the bound of
Theorem~\ref{thm:cesaro}, the comparisons with the sieve, and the
grid of Corollary~\ref{cor:half}, for the main text, for variant~C
and for the fallback (b$'$).
\item \texttt{verify\_small.py} $\to$
\texttt{verify\_small\_q4.out}, \texttt{verify\_small\_q3.out}
\textbf{[double]}: independent recomputation up to $2^{21}$ (a
different code path), the identities \eqref{eq:keyid} and
Lemma~\ref{lem:pairs} to $22$--$23$ digits, the cone check.
\item \texttt{general\_pairs.py} $\to$ \texttt{general\_pairs.out}
\textbf{[double]}: numerics for pairs outside the cone (not used).
\end{itemize}

\noindent\textbf{Directory \texttt{monotonicity/extensions/}
(Section~\ref{sec:zeta} only).} \texttt{zeta\_zeros.py}
\textbf{[ball]} ($10^4$ zeros of $\zeta$ from Arb);
\texttt{common.py}, \texttt{zeta\_thresholds.py},
\texttt{zeta\_density.py} \textbf{[double]}: only their outputs for
$M_\sigma(\zeta)$, $M_0(\zeta)$, the three values of $\lambda_\zeta$
and the four densities quoted in Section~\ref{sec:zeta} are used.

\medskip\noindent\textbf{Summary of what each theorem takes from a
computation.}
\begin{center}
\small
\begin{tabular}{p{0.34\textwidth}p{0.58\textwidth}}
\toprule
statement & computational input and its class\\
\midrule
Theorems~\ref{thm:Theta1}, \ref{thm:cesaro}, \ref{thm:cone},
\ref{thm:equiv}, Cor.~\ref{cor:inherit} (first part) &
\texttt{cert\_small.py}~(A) \textbf{[exact]};
Lemma~\ref{lem:theta}(a): (C) \textbf{[ball]} to $2^{22}$, sieve
\textbf{[ld]} beyond (not needed in variant~C); closed-form constants
and evaluations of $\cL$ \textbf{[double]}, margins $\ge0.018$\\
Corollary~\ref{cor:half} & in addition (D), (E) \textbf{[ball]};
grid of $\beta$ \textbf{[double]} with the margins stated in its
proof; \cite{BKLNW21} for $x>10^{20}$ ($\chi_4$)\\
Theorems~\ref{thm:cantelli-q}, \ref{thm:NRq},
Table~\ref{tab:bias} & zero lists \textbf{[ball]};
\texttt{certified\_bounds.py} \textbf{[ball]};
\texttt{nonresonance\_certified.py} \textbf{[exact]}/\textbf{[ball]}\\
all other theorems of Sections~\ref{sec:resolvent}--\ref{sec:ties} &
none\\
Section~\ref{sec:zeta}, numerical values & \textbf{[double]}; not
refereed\\
\bottomrule
\end{tabular}
\end{center}

\section{Formal verification}\label{sec:formal}

\noindent\textbf{None of the theorems of this article has been
formalized.} What exists is a Lean~4 development of a \emph{related}
statement, in the directory \texttt{monotonicity/lean-race/} of the
repository; the following description is taken from its
\texttt{README.md} (status dated 24 November 2026), and we have not
re-run the build for this article.

\begin{sloppypar}
\medskip\noindent\textbf{What is formalized.} The development is a
small layer (about $1250$ lines) on top of the public library
\texttt{anthropics/formal-math}, folder \texttt{zeta23/}, whose
literature-form explicit formula for a primitive Dirichlet
$L$-function carries the analytic weight. It is reported to compile
with no \texttt{sorry}, no \texttt{axiom} and no
\texttt{native\_decide}, and \texttt{\#print axioms} reports only
\texttt{propext}, \texttt{Classical.choice}, \texttt{Quot.sound} for
all $16$ headline statements. For a window $f$ put
$g(x)=\int f(t)\overline{f(t-x)}\,dt$, so that $\hat g=|\hat f|^2\ge0$
on the real line. The headline theorem
(\texttt{grh\_chi4\_smoothed\_race\_real}) reads, in words: assume
the Riemann hypothesis for the single $L$-function $L(s,\chi_4)$, as
an explicit hypothesis of the theorem
(\texttt{GRHFor chi4}: every zero $s$ with $0<\re s<1$ has
$\re s=\tfrac12$). Let $f$ be a real $C^2$ window supported in
$[-a/2,a/2]$ with $a<\log3$, and suppose that the explicit number
\[
\Gamma_0=\frac1{2\pi}\int|\hat f(r)|^2\Bigl[\re\frac{\Gamma'}{\Gamma}
\Bigl(\frac34+\frac{ir}2\Bigr)+\log\frac4\pi\Bigr]dr
\quad\text{satisfies}\quad
\Gamma_0<\tfrac12\Bigl(\int f\Bigr)^2-\mathrm{margin}.
\]
Then there is $L_0$ such that for every $L\ge L_0$
\[
\sum_p\chi_4(p)\,\frac{\log p}{\sqrt p}\,g(\log p-L)<-\mathrm{margin}.
\]
With $\mathrm{margin}\ge0$ and $f\ge0$ this says that the primes
$\equiv3$ mod $4$ are ahead in this \emph{smoothed} race for all
large $L$, with no exceptional set, no linear-independence
hypothesis and no simplicity hypothesis. A two-sided form (for every
$\eps>0$, eventually
$|\sum_p\chi_4(p)(\log p)p^{-1/2}g(\log p-L)+\tfrac12(\int f)^2|
\le\Gamma_0+\eps$) is proved with no numerical hypothesis; an
underlying inequality for the smoothed sum over prime powers is
proved for a general primitive character under the corresponding
hypothesis \texttt{GRHFor}~$\chi$; and the prime number theorem in
the form $\theta(x)/x\to1$, which the argument needs for the prime
squares, is proved (from the library's port of a quantitative prime
number theorem), not assumed. (That module of the library is outside
the library's default build target; the README reports that it was
compiled for this development and that its axiom report shows only
the three standard axioms.)
\end{sloppypar}

\medskip\noindent\textbf{Hypotheses that remain in the formal
statements.} Exactly two, both explicit hypotheses of the theorems
and never axioms: (1) GRH for the one $L$-function; (2) for the
one-sided race form only, the real inequality on $\Gamma_0$ displayed
above. \emph{No window is certified}: the inequality must come from a
certified computation, and the only numerics available are
floating-point values that prove nothing (they indicate room for a
positive margin for a plateau window with quintic ramps).

\medskip\noindent\textbf{What is not formalized.}
\begin{itemize}
\item Any statement of the present article: not the unsmoothed race;
not the Ces\`aro mean $\cM$ (the first Riesz mean) and hence none of
Theorem~\ref{thm:allx}; not \cite[Theorem~A]{MC}
(Theorem~\ref{thm:MCA}); not the limiting distribution, the convex
order, the variance or the density bounds of
Theorems~\ref{thm:order}--\ref{thm:tiesintro}; nothing involving LI.
\item GRH itself, in any form. The development neither imports,
checks nor depends on the announced proof mentioned in
Remark~\ref{rem:GRHstatus}; its hypothesis only has the same shape as
the statement of that announcement, instantiated at one character.
\item A certified value of $\Gamma_0$ for any window.
\item The race statements for characters other than $\chi_4$ (only
the underlying inequality for prime powers is general).
\item An effective $L_0$: the theorems give existence only.
\item An independent comparison of statement and proof with the
\texttt{comparator} tool was not run; as weaker substitutes the
README reports a textual check that the challenge statements are
copies of the solution statements and a check that the $16$
statements elaborate to identical types on both sides.
\end{itemize}
The formalized statement and Theorem~\ref{thm:allx} are of the same
kind --- under GRH$(\chi_4)$ alone, a positive bias term beats a sum
over zeros after smoothing, for all large arguments --- but they
concern different smoothings (a compactly supported multiplicative
window on the scale $\log p$ there; the Ces\`aro mean here), and the
formal one is asymptotic and conditional on a numerical inequality
that has not been certified. It should not be read as a formal
verification of any result above.

\section{Gaps and caveats}\label{sec:gaps}

\begin{enumerate}
\item\label{gap:margin} \textbf{A discrepancy found while this
article was being written (margins in the proof of
Corollary~\ref{cor:half}).} The working notes say that the
double-precision grid evaluation of $\beta(a)$ in the proof of
Corollary~\ref{cor:half} has ``margins $\ge0.03$''. That is true
only from $a=10^{10}$ on (smallest margin $0.0346$, as the output
file \texttt{bounds.out} says). A direct re-evaluation of the
closed-form function for this article gives
$\beta-\tfrac12=1.5\cdot10^{-4}$ at the first grid point
$a=1.514\cdot10^8$ and $7.9\cdot10^{-4}$ at $a=1.585\cdot10^8$ for
$\chi_4$ ($1.1\cdot10^{-4}$ and $7.4\cdot10^{-4}$ in variant~C), and
$3.9\cdot10^{-4}$ at $a=2.512\cdot10^7$ for $\chi_3$
($3.1\cdot10^{-4}$ in variant~C). The statement of
Corollary~\ref{cor:half} is not affected --- a margin of $10^{-4}$ is
far above the rounding error of a double-precision evaluation of an
explicit elementary expression --- but the evaluation is not in
interval arithmetic and its margin near the junction with the
certified range is three orders of magnitude smaller than the notes
state. A reader who wants a comfortable margin can extend part~(E) of
\texttt{cert\_small.py} (ball arithmetic) from $1.6\cdot10^8$ to
$10^{10}$; this was not done. The re-evaluation reproduced every
other number quoted for $\cL$ (thresholds, values at $10^4$ and at
$Y_1^2$, grid end points).
\item \textbf{Other differences between the sources, all immaterial.}
(i) The working note on Section~\ref{sec:explicit} says in one place
that Lemma~\ref{lem:theta}(a) is ball-certified up to $2^{21}$ and
elsewhere up to $2^{22}$; the script output says $2^{22}$, which is
what we state. (ii) $M_\sigma(\chi_3)$ appears as $2.1277\ldots$ and
as $2.1278$ (truncation versus rounding of $2.12775$). (iii) The
tier-B variance bounds are $0.1555736302$ and $0.1132342076$ in
Proposition~\ref{prop:tierB}, while the non-resonance script uses the
slightly larger valid upper bounds $0.1555736304$ and
$0.1132342078$. (iv) The headers of
Propositions~\ref{prop:Vplus}, \ref{prop:tierB} in the working notes
do not list NCZ$(\chi)$, which their proofs use through
Lemmas~\ref{lem:mult}, \ref{lem:JA}; it is listed here.
\item \emph{Section~\ref{sec:zeta} has not been separately refereed.}
It contains statements only; its numerical values are computed, not
certified; the tail bounds it uses are quoted from \cite{Tru14} as
they appear in the working note and were not re-checked against the
published statement.
\item \emph{Quoted inputs.} Lemmas~\ref{lem:WRclassical},
\ref{lem:WRshift}, \ref{lem:WRexp} and
Proposition~\ref{prop:WRvar} are quoted from \cite{WR} (through
\cite{MC}) and not re-derived; Theorem~\ref{thm:MCA} is quoted from
\cite{MC}, which is itself an unsubmitted draft with its own list of
gaps. The referee of the working notes checked that the quotations
agree with the statements of \cite{WR} and with a sieve, not the
proofs in \cite{WR}.
\item \emph{Lemma~\ref{lem:zerosq}(a).} The contour estimate is
written out, but the two bounds (i), (ii) for $L'/L$ that it uses are
the standard ones of \cite[Ch.~19]{Dav00}, quoted from memory; we
have no numbered citation for the formula itself in the $L$-function
case (\cite[Theorem~27]{Ing32} is the $\zeta$ case). The formula is
confirmed numerically in three ways. Only the size $O(x)$ of the
non-zero terms and the sign of
$\mathfrak b\,t-\log t-1-\mathfrak b_0$ matter for the theorems.
\item \emph{Citations checked only partially.}
\cite[Corollary~11.2 and Appendix~A.1.1]{BKLNW21},
\cite[Theorem~2]{But18}, \cite{Joh23}, \cite{Zhao25} and the
introduction of \cite{PT19} were read in the arXiv versions; the
numbering in the published versions was not checked. \cite{PT21} was
not read; the height $3\,000\,175\,332\,800$ is as quoted in
\cite{BKLNW21}. \cite[(3.15)--(3.16)]{RS62} is quoted from
\cite[Lemma~2.2]{Joh23}, not from the original. The originals of
\cite{HL18,Lan18,Pin91,Ing32} were not read (they are known to us
through \cite{PT19} and \cite{Joh23}). The equation numbers
\cite[Ch.~12, (17)--(18)]{Dav00} and the lemma number
\cite[Lemma~15.1]{MV07} are quoted from memory.
\cite[Theorem~1.2]{RS94} and \cite{Kac93} are quoted through
\cite{Dev20} and through surveys; \cite{BT15} was not read;
\cite[Theorem~6.9]{Vil09} was not checked against the book. For
\cite{Fio14,Dev20,Bai21,MN20,Pla16,PT26,Hay25} the statements used
were read in the arXiv versions (by the author of the working notes
or by its referee). Every entry of the reference list whose details
were not verified against the published source carries a visible
\textsf{TODO-cite} mark.
\item \emph{Floating point in Section~\ref{sec:explicit}.} In
variant~C (Remark~\ref{rem:variantC}) every finite computation that
the proofs of Theorems~\ref{thm:Theta1}, \ref{thm:cesaro},
\ref{thm:cone}, \ref{thm:equiv} and
Corollaries~\ref{cor:half}, \ref{cor:inherit} use is exact or
ball-certified: $\Theta_1(p)>0$ for $p\le10^4$ (exact integers),
Lemma~\ref{lem:theta}(a) for $y\le2^{22}$ (Arb), the gap bounds for
$\cN$ up to $1.6\cdot10^8$ resp.\ $2.6\cdot10^7$ (Arb). The main text
($Y_1=10^{10}$) uses in addition Lemma~\ref{lem:theta}(a) on
$(2^{22},10^{10}]$ from the \texttt{long double} sieve, with margins
$0.089\sqrt y$ and $0.26\sqrt y$ against a worst-case accumulated
error of $0.21$; that range is not interval-certified (it would need
directed rounding in the C code or a much slower ball-arithmetic
sieve). In both variants the evaluations of the closed-form function
$\cL$ are ordinary double precision (margins $\ge0.018$ for the
positivity at $10^4$), and Corollary~\ref{cor:half} uses the grid
argument discussed in item~\ref{gap:margin} and, for $x>10^{20}$ and
$\chi_4$ only, Lemma~\ref{lem:theta}(b).
\item \emph{Tier B} (\certB) rests on reading Platt's computation
\cite{Pla16} as a proof of simplicity (Remark~\ref{rem:tierB}). The
tier-A constants rest on the correctness of the zero lists, of
\texttt{certify\_arb.py} and of the Arb library; we did not audit the
latter.
\item \emph{Atoms.} Under GRH alone we cannot exclude
$\nu_M(\{0\})>0$, so the existence of the logarithmic density of
$\{D_m>0\}$ is not proved (it is open even at $m=0$:
\cite[Remark~1.3]{Fio14}). Theorem~\ref{thm:ties}(b) is the best we
have; Proposition~\ref{prop:selfsuff} needs one self-sufficient
ordinate, an unproved hypothesis.
\item \emph{Strict monotonicity of the densities without LI} is not
obtained (only weak monotonicity, and only on a range where
$\lambda_\chi>0$). The convex order does not imply monotonicity of
$\nu_M((0,\infty))$.
\item \emph{The gap to the truth.} Under GRH alone we prove $0.8842$
(tier~A) / $0.9227$ (tier~B) at $q=4$, $m=0$; under GRH and LI the
density is $0.9959$. Higher-degree multipliers
(Remark~\ref{rem:NRlimits}(1)) and longer certified zero lists would
narrow the gap; neither was done.
\item \emph{NCZ-free versions} of
Sections~\ref{sec:resolvent}--\ref{sec:biasbounds} are only sketched
(Remark~\ref{rem:noNCZ}). For $q=3,4$ NCZ is unconditional.
\item \emph{Natural densities.} The Cantelli-type bounds are for
logarithmic density only. Whether the natural density of
$\{D_m>0\}$ fails to exist for every $m>-\tfrac12$ under GRH (as is
attributed to Kaczorowski at $m=0$) was not examined.
\item \emph{Theorem~\ref{thm:leadsets} is qualitative} for
$M\le\gamma_1/\sqrt3$: the proof of
Proposition~\ref{prop:unbounded} gives no explicit lower bound for
$\nu_M((-\infty,0))$ there.
\item \emph{Effectivity and scope of the all-$x$ results.}
Theorem~\ref{thm:allx} is proved for $q\in\{4,3\}$ and the cone
$m\le0$, $m<m'\le1$ only. For other pairs with
$\lambda_\chi(M,M')>0$ the threshold $x_0$ of
Theorem~\ref{thm:MCA}(ii) remains inexplicit; Lemma~\ref{lem:pairs}
reduces the problem to the sign of $\Phi_{m,m'}$, but the certified
margin, the explicit lower-order terms and the finite computation
that an explicit $x_0$ would need were not carried out, and the
numerical values of $\lambda_\chi$ for such pairs in the repository
are not certified. For pairs with $\lambda_\chi(M,M')<0$ the method
gives nothing.
\item \emph{Novelty.} The statements ``not found'' in
Section~\ref{subsec:new} are limited by the reading recorded there
and in the repository; in particular \cite{LWS26} was seen at
abstract level only, and \cite{Suz25,Fuj88} and the 1918 originals
were not read.
\item \emph{The announcement mentioned in
Remark~\ref{rem:GRHstatus}} is not used anywhere, and we examined
neither the note nor the formalization.
\item \emph{Relation to \cite{WR}.} The density statements of
\cite{WR} are made under LI, where the zeros are simple and
$V_M=\sigma_M^2$. Any statement about mean squares under GRH alone
must be read with $V_M$ (Theorem~\ref{thm:var}).
\item \emph{No claim is made} for non-real characters or for races
with more than two competitors; outside
Section~\ref{sec:explicit} no claim is made for $m\le-\tfrac12$.
\end{enumerate}

\appendix
\section{Review history}\label{sec:review}

This article is assembled from three working notes of the repository
(directory \texttt{monotonicity/}). Two of them were each read by an
adversarial AI referee on 24 November 2026 and repaired; the third
(the source of Section~\ref{sec:zeta}) was not refereed. Neither
referee found a fatal error. Their findings and the resolution of
each are summarised below; ``R'' / ``REP'' items were classified as
repairable, the others as minor. The companion manuscript \cite{MC},
whose Theorem~A is quoted, has its own review history.

\subsection*{Referee of the note behind
Sections~\ref{sec:explicit}--\ref{sec:landau}
(\texttt{explicit/REFEREE-REPORT.md})}
The referee re-derived every identity, sign and constant with
independent code (the summation identity and the cone to $30$ digits
at several points; the residues of Lemma~\ref{lem:zerosq}; the
constants $\sigma_0^2$, $c_1$, $\cL$; the grid of
Corollary~\ref{cor:half}; an independent sieve to $10^7$).
\begin{itemize}
\item[REP-1.] \emph{The input of Lemma~\ref{lem:theta}(b) was an
arXiv preprint (arXiv:1002.0442) whose own input is an unpublished
zero verification.} Resolved: replaced by
\cite[Corollary~11.2]{BKLNW21} (the same inequality, so no constant
changes), whose input is covered by \cite{PT21}; \cite{PT21} itself
was not read.
\item[REP-2.] \emph{The contour estimate in Lemma~\ref{lem:zerosq}(a)
was only cited.} Resolved: written out (choice of $T_k$, the bounds
(i), (ii), the horizontal and left integrals, the limit, the case
$x=1$). Equation numbers in \cite{Dav00} remain unchecked.
\item[REP-3.] \emph{The sign of $\Theta_1(p)$ was decided in
\texttt{long double}; a stated margin was wrong at $p=7$; the values
of $\cN$ used in Corollary~\ref{cor:half} were uncertified.}
Resolved: new script \texttt{cert\_small.py} (parts A--E of
Section~\ref{sec:comp}); the margin statement corrected
($\min\Theta_1=1.1756$ at $7$, $G(7)=0.0635$); variant~C added.
\item[REP-4.] \emph{Novelty of the summation identity; priority for
the Landau mechanism.} Resolved: Lemma~\ref{lem:ingham} is attributed
to Ingham and to \cite[Lemma~2.7]{Joh23}; \cite{HL18,Lan18,Pin91} are
cited; what is claimed is only the all-$x$ statement, its reading as
lead inheritance, and the one-signedness observation with its uses.
\item[Minor.] The cone is written ``$m\le0$, $m<m'\le1$'' everywhere.
The quotation of \cite[Lemma~2.2]{Joh23} has the range $x\ge41$; the
constant $2.05282$ is marked as recalled; Lemma~\ref{lem:theta}(a)
says that it rests on the computation. $\inf\cN$ over
$[11,10^{10}]$ is stated exactly. Vacuous checks of the explicit
formula at $x=10^3,10^4$ were replaced by the checks at small $x$.
The parenthesis on LI in Remark~\ref{rem:halfLI} includes simplicity
and explains $\liminf\cN=\liminf G$. Corollary~\ref{cor:half} for
$\chi_3$ is stated from $x\ge6$. The Landau steps and the role of
real zeros are written out. One cosmetic remark on an output file
could not be reproduced.
\item[Not fixed.] Interval certification of the \texttt{long double}
sieve on $(2^{22},10^{10}]$ (variant~C avoids that range);
verification of equation numbers in \cite{Dav00,MV07,RS62}, of the
published versions of the arXiv papers, and of the originals of
\cite{HL18,Lan18,Pin91}.
\end{itemize}

\subsection*{Referee of the note behind
Sections~\ref{sec:resolvent}--\ref{sec:biasbounds}
and~\ref{sec:ties} (\texttt{REFEREE-4.md})}
The referee re-derived every theorem, proposition and lemma,
reproduced every certified constant (independently in
\texttt{mpmath}, and by rerunning \texttt{certify\_arb.py}), tested
Lemma~\ref{lem:moments} on a maximally resonant configuration, and
tested Theorem~\ref{thm:NR} on toy models, including adversarial
resonant tails and controls with a hypothesis dropped (the bound was
never violated, and was violated when (NR2) or (NR3) was dropped). No
number in the summary changed at the displayed precision.
\begin{itemize}
\item[R1.] \emph{Corollary~\ref{cor:mono} over-claimed for the
exponential functional, and its last sentence was vacuous.}
Resolved: the time-average claim is restricted to functionals of
polynomial growth; the exponential functional is a statement about
$\nu_M$ only, conditional on finiteness; the last sentence was
replaced by \eqref{eq:transport}. Remark~\ref{rem:noexp} was added.
\item[R2.] \emph{The comparison with \cite{BMOR21} after
Lemma~\ref{lem:mult} was wrong by a factor $2$ in the working note}
($N(T,\chi)$ counts $|\gamma|\le T$). Resolved: the bound gives
$2m_\gamma\le9.53$, i.e.\ $m_\gamma\le4$; the claim that
Lemma~\ref{lem:mult} is sharper was dropped. Nothing downstream used
it.
\item[R3.] \emph{After Proposition~\ref{prop:endpoint}:}
$\exp(-1/(8C(\chi)M^2))$ is an upper bound under LI, not the true
order. Resolved.
\item[R4.] \emph{Hypotheses of Theorem~\ref{thm:ties}.} Resolved: (a)
and (b) under GRH via \cite[Theorem~A]{MC} with the general
$\lambda_\chi$; (c) under GRH and NCZ; the statement on ties is
restricted to an interval on which $\lambda_\chi>0$.
\item[R5.] \emph{Theorem~\ref{thm:NRq} was a rigorous inequality
evaluated in floating point.} Resolved: new script
\texttt{nonresonance\_certified.py} (exact-integer verification of
(NR1)--(NR3); Arb enclosure of $J$ at a rational feasible point).
\item[R6.] \emph{Numbering inconsistencies between the note and its
scripts.} Resolved in the note; the numbering of the present article
is its own.
\item[R7.] \emph{Citations of Martin--Ng.} Resolved:
\cite[Theorem~1.5(a)]{MN20} is a two-way statement; their definition
is on the multiset; the set-based definition used here is
\cite[Definition~5.3]{Dev20}.
\item[Minor.] The split at $u=y/2$ in
Theorem~\ref{thm:resolvent}(b) is for $y\ge4$. The constant $c'$ is
written unambiguously, with its values. ``Soft thresholding gains
nothing at $m=0$'' is stated for tier~A only. The soft threshold is
no longer restricted in the script (two table values at $m=-0.45$
improved). ``Decidable'' became ``verifiable when true''. The grid
check of the closed formulas is in a script; the LI value of
$\|F-1\|_4$ is $\approx0.51$. Theorem~\ref{thm:NR} states that the
zero list must be complete. \cite[Theorem~1.2]{RS94} is quoted as a
bound for the mass outside a ball, the one-sided statements being
credited to \cite[Theorem~2.5(1)]{Dev20}. ``If odd'' became ``if
$\chi$ is odd''. One further item of the report concerned text
addressed to the earlier preprint \cite{WR} rather than to the
present results; it is not reproduced here (see
Section~\ref{sec:gaps}, the item on the relation to \cite{WR}).
\item[Not done.] The referee's observation that
$y\cdot(E_M-\int E_{M'}d\mu)$ is still growing at $10^7$ for
$(m,m')=(-0.4,0)$ is consistent with Section~\ref{sec:comp} and
needed no change.
\end{itemize}

\subsection*{Changes made in assembling this article}
Notation was unified (in particular $\cM$, $\cN$, $\cL$ for the
Ces\`aro mean, its normalisation and the explicit lower-bound
function; $\omega$ for points of the group $\cA$; $\varkappa$,
$\beta$ for the multiplier of Theorem~\ref{thm:NR}; $\eta$ for the
parameter of Lemma~\ref{lem:mult}). The proof of
Theorem~\ref{thm:equiv} uses the more careful formulation of the
analytic continuation from \cite[the Landau-type converse]{MC}. NCZ$(\chi)$ is
listed explicitly in
Propositions~\ref{prop:Vplus}, \ref{prop:tierB}. The margins in the
proof of Corollary~\ref{cor:half} were recomputed
(Section~\ref{sec:gaps}, item~\ref{gap:margin}). No statement was
otherwise weakened or strengthened.

\begin{thebibliography}{99}
\bibitem{WR} R. Sharma, \emph{Rubinstein--Sarnak densities for
weighted prime number races}, preprint, August 2026,
\href{https://doi.org/10.5281/zenodo.21947101}{doi:10.5281/zenodo.21947101};
code and data: \url{https://github.com/raghavs1729/weighted-prime-races}.
\bibitem{MC} R. Sharma, \emph{Monotonicity of Chebyshev's bias in the
weight: a coupling argument, inheritance of leads under GRH, and
logarithmic weights}, draft of November 2026, not yet submitted; file
\texttt{monotonicity/monotone-coupling.tex} of
\url{https://github.com/raghavs1729/weighted-prime-races}.
\bibitem{RS94} M. Rubinstein, P. Sarnak, \emph{Chebyshev's bias},
Experiment.\ Math.\ \textbf{3} (1994), 173--197.
\todocite{Theorem~1.2 is quoted through \cite{Dev20} and
\cite{Fio14}; the original was not consulted for this article}
\bibitem{ANS14} A. Akbary, N. Ng, M. Shahabi, \emph{Limiting
distributions of the classical error terms of prime number theory},
Q. J. Math.\ \textbf{65} (2014), 743--780. arXiv:1306.1657.
\todocite{statements checked against the arXiv version; journal pages
not checked}
\bibitem{Dev20} L. Devin, \emph{Chebyshev's bias for analytic
$L$-functions}, Math.\ Proc.\ Cambridge Philos.\ Soc.\ (2020).
arXiv:1706.06394. \todocite{Theorems~2.1, 2.2, 2.3, 2.5,
Corollary~5.8 and Definition~5.3 read in the arXiv version; volume
and pages to be added; numbering in the published version not
checked}
\bibitem{Fio14} D. Fiorilli, \emph{Highly biased prime number races},
Algebra Number Theory \textbf{8} (2014). arXiv:1210.6946.
\todocite{Remark~1.3 and Lemmas~2.5, 2.7 read in the arXiv version;
pages to be added; numbering in the published version not checked}
\bibitem{Bai21} A. Bailleul, \emph{Explicit Kronecker--Weyl theorems
and applications to prime number races}, arXiv:2007.05763.
\todocite{Theorem~1.7 read in the arXiv version; publication data not
checked}
\bibitem{MN20} G. Martin, N. Ng, \emph{Inclusive prime number races},
Trans.\ Amer.\ Math.\ Soc.\ (2020). arXiv:1710.00088.
\todocite{Definition~1.3 and Theorem~1.5 read in the arXiv version;
volume and pages to be added}
\bibitem{Kac93} J. Kaczorowski, \emph{A contribution to the
Shanks--R\'enyi race problem}, Quart.\ J. Math.\ Oxford
\textbf{44} (1993). \todocite{not read; Theorem~1 is quoted from the
surveys arXiv:1202.3408 (\S7) and arXiv:1910.08983 (\S4.1); pages to
be added}
\bibitem{Hay25} M. Hayani, \emph{Chebyshev's bias without linear
independence}, arXiv:2512.23302 (2025).
\bibitem{Sheth24} A. Sheth, \emph{Euler products at the centre and
applications to Chebyshev's bias}, Math.\ Proc.\ Cambridge Philos.\
Soc.\ \textbf{179} (2025), 331--349. arXiv:2405.01512.
\todocite{entry reused from \cite{WR}; not re-checked}
\bibitem{GW18} J. Gerard, L.~C. Washington, \emph{Sums of powers of
primes}, Ramanujan J. \textbf{45} (2018), 171--180.
\todocite{not verified; cited as described in \cite{Was24}}
\bibitem{Was24} L.~C. Washington, \emph{Sums of powers of primes II},
Ramanujan J. \textbf{65} (2024), 783--795. arXiv:2209.12845.
\todocite{read in the arXiv text; journal reference not verified}
\bibitem{BBLMR24} M. Boran, J. Byun, Z. Li, S.~J. Miller, S. Reyes,
\emph{Sums of powers of primes in arithmetic progression}, PUMP J.
Undergrad.\ Res.\ \textbf{7} (2024), 29--50. arXiv:2309.16007.
\bibitem{HL18} G.~H. Hardy, J.~E. Littlewood, \emph{Contributions to
the theory of the Riemann zeta-function and the theory of the
distribution of primes}, Acta Math.\ \textbf{41} (1917), 119--196.
\todocite{details as in the reference list of \cite{PT19}, which
gives the year 1917 (the volume is also cited as 1916 or 1918);
original not consulted}
\bibitem{Lan18} E. Landau, \emph{\"Uber einige \"altere Vermutungen
und Behauptungen in der Primzahltheorie}, I and II, Math.\ Z.
\textbf{1} (1918), 1--24 and 213--219. \todocite{details as in the
reference list of \cite{PT19}; originals not consulted}
\bibitem{Fuj88} A. Fujii, \emph{Some generalizations of Chebyshev's
conjecture}, Proc.\ Japan Acad.\ Ser.~A Math.\ Sci.\ \textbf{64}
(1988), no.~7, 260--263. \todocite{details as in the reference list
of \cite{PT19}; original not consulted}
\bibitem{PT19} D.~J. Platt, T.~S. Trudgian, \emph{Fujii's development
on Chebyshev's conjecture}, Int.\ J. Number Theory \textbf{15}
(2019), 639--644. arXiv:1802.00089. \todocite{arXiv text read; journal
pages not checked}
\bibitem{Joh23} D.~R. Johnston, \emph{On the average value of
$\pi(t)-\mathrm{li}(t)$}, Canad.\ Math.\ Bull.\ \textbf{66} (2023),
185--195. arXiv:2201.06184. \todocite{arXiv version~2 read in full;
journal reference and numbering in the published version not
verified}
\bibitem{Pin91} J. Pintz, \emph{On an assertion of Riemann concerning
the distribution of prime numbers}, Acta Math.\ Hungar.\ \textbf{58}
(1991), 383--387. \todocite{not read; cited as described in
\cite{Joh23}}
\bibitem{Ing32} A.~E. Ingham, \emph{The Distribution of Prime
Numbers}, Cambridge Tracts in Mathematics and Mathematical Physics
30, Cambridge University Press, Cambridge, 1932. \todocite{not
consulted; pp.~103--104 and Theorem~27 as cited in \cite{Joh23}}
\bibitem{Zhao25} T. Zhao, \emph{On the mean values of the error terms
in Mertens' theorems}, Res.\ Number Theory \textbf{11} (2025),
paper~62. arXiv:2411.18903.
\bibitem{Suz25} M. Suzuki, \emph{On variants of Chebyshev's
conjecture}, Ramanujan J. \textbf{68} (2025), article~95.
arXiv:2411.07436.
\bibitem{LWS26} Liu, Wang, Shi, \emph{Two problems in comparative
prime number theory: the five-way race modulo $11$, and the sign of
the integrated prime counting error}, Zenodo,
doi:10.5281/zenodo.22208026 (2026). \todocite{authors' initials;
seen at abstract level only}
\bibitem{HKMN24} A. Hamieh, H. Kadiri, G. Martin, N. Ng,
\emph{Comparative prime number theory problem list},
arXiv:2407.03530. \todocite{the problem cited was read in the arXiv
text; authors' initials and publication data not verified}
\bibitem{BKLNW21} S. Broadbent, H. Kadiri, A. Lumley, N. Ng, K. Wilk,
\emph{Sharper bounds for the Chebyshev function $\theta(x)$}, Math.\
Comp.\ \textbf{90} (2021). arXiv:2002.11068.
\todocite{Corollary~11.2 and Appendix~A.1.1 read in arXiv version~2;
the journal reference is from memory and was not verified; numbering
in the published version not checked; pages to be added}
\bibitem{But18} J. B\"uthe, \emph{An analytic method for bounding
$\psi(x)$}, Math.\ Comp.\ \textbf{87} (2018). arXiv:1511.02032.
\todocite{Theorem~2 read in the arXiv version; numbering in the
published version not checked; pages to be added}
\bibitem{PT21} D.~J. Platt, T.~S. Trudgian, \emph{The Riemann
hypothesis is true up to $3\cdot10^{12}$}, Bull.\ London Math.\ Soc.
\todocite{not read; cited as in \cite{BKLNW21}; volume, year and
pages to be added}
\bibitem{RS62} J.~B. Rosser, L. Schoenfeld, \emph{Approximate
formulas for some functions of prime numbers}, Illinois J. Math.\
\textbf{6} (1962). \todocite{original not consulted; (3.15)--(3.16)
as quoted in \cite[Lemma~2.2]{Joh23}; pages to be added}
\bibitem{Pla16} D.~J. Platt, \emph{Numerical computations concerning
the GRH}, Math.\ Comp.\ \textbf{85} (2016). arXiv:1305.3087.
\todocite{\S3 and Theorem~7.1 read in the arXiv version; pages to be
added}
\bibitem{PT26} D.~J. Platt, T.~S. Trudgian, \emph{Zeros of quadratic
Dirichlet $L$-functions}, arXiv:2610.29018 (2026).
\todocite{Theorem~1 and \S2.1 read in the arXiv version; authors'
initials not verified}
\bibitem{BMOR21} M.~A. Bennett, G. Martin, K. O'Bryant,
A. Rechnitzer, \emph{Counting zeros of Dirichlet $L$-functions},
Math.\ Comp.\ \textbf{90} (2021), 1455--1482.
\todocite{entry reused from \cite{WR}, not re-checked; the inequality
is quoted from the abstract of arXiv:2005.02989}
\bibitem{BT15} D.~G. Best, T.~S. Trudgian, \emph{Linear relations of
zeroes of the zeta-function}, arXiv:1209.3843. \todocite{not read;
authors' initials and publication data not verified}
\bibitem{Tru14} T.~S. Trudgian, \emph{An improved upper bound for the
argument of the Riemann zeta-function on the critical line II},
J. Number Theory (2014). \todocite{title recalled from memory; volume
and pages to be added; the bound
$|S(T)|\le0.112\log T+0.278\log\log T+2.510$ is quoted from a working
note and was not checked against the paper}
\bibitem{Vil09} C. Villani, \emph{Optimal Transport: Old and New},
Grundlehren der mathematischen Wissenschaften 338, Springer, Berlin,
2009. \todocite{theorem number not checked against the book}
\bibitem{Dav00} H. Davenport, \emph{Multiplicative Number Theory},
3rd ed., revised by H.~L. Montgomery, Graduate Texts in Mathematics
\textbf{74}, Springer, New York, 2000.
\bibitem{MV07} H.~L. Montgomery, R.~C. Vaughan, \emph{Multiplicative
Number Theory I: Classical Theory}, Cambridge Studies in Advanced
Mathematics 97, Cambridge University Press, Cambridge, 2007.
\bibitem{Bil99} P. Billingsley, \emph{Convergence of Probability
Measures}, 2nd ed., Wiley, New York, 1999.
\bibitem{Alp26} L. Alp\"oge, note dated 20 November 2026,
\url{https://alpo.ge/remark.pdf}. \todocite{not examined by the
author of this article; title to be added; re-check before
submission}
\bibitem{Lean26} Lean FRO and Mathlib Initiative, joint statement of
21 November 2026,
\url{https://leanprover-community.github.io/blog/posts/riemann-hypothesis/}.
\todocite{not examined by the author of this article; re-check
before submission}
\end{thebibliography}

\end{document}
